\documentclass[11pt]{amsart}
\usepackage[T1]{fontenc}
\usepackage{lmodern}
\usepackage[margin=1.05in]{geometry}
\usepackage{mathtools,amssymb,amsthm,mathrsfs}
\usepackage{microtype,enumitem,booktabs}
\usepackage{xcolor}
\definecolor{linkblue}{RGB}{25,60,110}
\usepackage[colorlinks=true,linkcolor=linkblue,citecolor=linkblue,urlcolor=linkblue]{hyperref}
\usepackage[capitalise,nameinlink,noabbrev]{cleveref}
\hypersetup{pdftitle={Algebraic Kuga--Satake correspondences and Hodge conjectures on a K3 quadratic locus},pdfauthor={OpenAI},bookmarksdepth=2}
\setlength{\emergencystretch}{1.5em}
\newcommand{\Q}{\mathbb Q}
\newcommand{\C}{\mathbb C}
\newcommand{\Z}{\mathbb Z}
\newcommand{\R}{\mathbb R}
\newcommand{\PP}{\mathbf P}
\renewcommand{\AA}{\mathbf A}
\DeclareMathOperator{\NS}{NS}
\DeclareMathOperator{\CH}{CH}
\DeclareMathOperator{\KS}{KS}
\DeclareMathOperator{\Hg}{Hg}
\DeclareMathOperator{\End}{End}
\DeclareMathOperator{\Hom}{Hom}
\DeclareMathOperator{\SO}{SO}
\DeclareMathOperator{\GL}{GL}
\DeclareMathOperator{\Spin}{Spin}
\DeclareMathOperator{\cl}{cl}
\DeclareMathOperator{\id}{id}
\DeclareMathOperator{\Pic}{Pic}
\DeclareMathOperator{\Km}{Km}
\DeclareMathOperator{\rank}{rank}
\DeclareMathOperator{\sig}{sig}
\DeclareMathOperator{\Tr}{Tr}
\DeclareMathOperator{\im}{im}
\DeclareMathOperator{\coker}{coker}
\DeclareMathOperator{\Aut}{Aut}
\DeclareMathOperator{\Gr}{Gr}
\DeclareMathOperator{\Sym}{Sym}
\DeclareMathOperator{\Alb}{Alb}
\DeclareMathOperator{\Jac}{Jac}
\DeclareMathOperator{\Hdg}{Hdg}
\newtheorem{theorem}{Theorem}[section]
\newtheorem{proposition}[theorem]{Proposition}
\newtheorem{lemma}[theorem]{Lemma}
\newtheorem{corollary}[theorem]{Corollary}
\theoremstyle{definition}
\newtheorem{definition}[theorem]{Definition}
\theoremstyle{remark}
\newtheorem{remark}[theorem]{Remark}
\numberwithin{equation}{section}
\setlist[enumerate]{label=\textup{(\roman*)},leftmargin=*,itemsep=3pt}
\setlist[itemize]{leftmargin=*,itemsep=3pt}
\setcounter{tocdepth}{1}
\allowdisplaybreaks[1]
\title[Hodge conjectures on a K3 quadratic locus]{Algebraic Kuga--Satake correspondences and Hodge conjectures on a K3 quadratic locus}
\author{OpenAI}
\date{September 30, 2026}
\begin{document}
\begin{abstract}
We prove the rational Hodge and generalized Hodge conjectures in every
cohomological degree of every self-power of a projective complex K3
surface whose transcendental quadratic space, with its cup-product form,
admits a rational isometric embedding in
$\mathbb U_{\Q}^{\oplus2}\perp\langle-1\rangle^4$.
This includes every ample $P$-polarized K3 surface, including all Picard jumps,
for $P=\mathbb U\oplus D_8(-1)\oplus D_4(-1)$.
On this locus we construct algebraic correspondences inducing every
prescribed standard even-Clifford Kuga--Satake tensor.  More generally,
for any projective complex K3 surface, algebraicity of one exact standard
even-Clifford Kuga--Satake tensor implies both conjectures for every self-power.
\end{abstract}
\maketitle
\tableofcontents

\section{Introduction}
\label{intro:section}

The Hodge conjecture asks which cohomology classes on a smooth
projective complex variety come from algebraic subvarieties.  In its
rational form it asserts surjectivity of
\[
 \cl:\CH^r(X)_{\Q}\longrightarrow
 \Hdg^r(X):=H^{2r}(X,\Q)\cap H^{r,r}(X)
\]
for every $r$.  The generalized Hodge conjecture asks the corresponding
support question for whole Hodge substructures.  Write
\[
 N^pH^j(X,\Q)=
 \sum_{\substack{B\subset X\text{ closed}\\\operatorname{codim}_X B\ge p}}
 \ker\bigl(H^j(X,\Q)\longrightarrow H^j(X\setminus B,\Q)\bigr).
\]
The generalized conjecture says that this equals the largest rational
sub-Hodge structure whose types $(a,b)$ satisfy $a,b\ge p$.
We use rational coefficients for both conjectures throughout, and
abbreviate them as HC and GHC.

For a projective K3 surface $S$, let
\[
 T_S=\NS(S)_{\Q}^{\perp}\subset H^2(S,\Q)
\]
with its cup-product form $q$.  Its extreme Hodge spaces are
$T_S^{2,0}=H^{2,0}(S)$ and its conjugate, both one-dimensional.
The Hodge index theorem gives signature $(2,20-\rho(S))$.
This small Hodge level makes K3 surfaces a natural setting for the
generalized conjecture, but their powers contain tensor Hodge classes
and substructures that need not be generated by divisor classes alone.

Kuga and Satake attached abelian varieties to polarized K3 Hodge
structures through Clifford algebras \cite{KS}, and Deligne developed
the construction in families \cite{Deligne}.  The distinguished
correspondence is absolutely Hodge and motivated
\cite[Proposition~6.26\textup{(d)} and Corollary~6.27]{DeligneMilne},
\cite[Section~7.1]{Andre}.  Algebraicity is the additional assertion
needed here: we require a cycle inducing the prescribed Clifford map,
not merely some nonzero Hodge embedding into abelian cohomology.

\subsection{One exact Kuga--Satake tensor}

We specify the Kuga--Satake tensor used throughout the paper.  For a
quadratic space $(T,q)$ of K3 type we use the Clifford relation
$v^2=q(v,v)$.  Its full even algebra $W=C^+(T,q)$ carries the standard
left Kuga--Satake Hodge structure.  If $e_1,e_2$ are an oriented
orthonormal basis of the positive period plane with
$T^{2,0}=\C(e_1+ie_2)$, its complex structure is left multiplication
by $J=e_1e_2$.  Let $A$ be an abelian variety with a specified rational
Hodge identification $H^1(A,\Q)=W$.  Such a realization is unique up
to the usual isogeny choices.  For a nonisotropic $w\in T$ the
standard map is
\begin{equation}\label{main:kappa}
 j_w:T(1)\longrightarrow\End(W),\qquad j_w(v)(a)=vaw.
\end{equation}
Here $\Q(1)$ has weight $-2$.  The right factor need not be a Hodge
vector: conjugation by left spin multiplication makes
\eqref{main:kappa} a Hodge map at every period.

Let $E$ be a rational polarization of $W$ and set
$E^\flat:W\to W^\vee(-1)$, $y\mapsto E(y,-)$.
In the ordering $\End(W)=W\otimes W^\vee$, apply $(E^\flat)^{-1}$
to the second factor after twisting by $\Q(-1)$.  This gives the
specified tensor form
\begin{equation}\label{main:iota}
 \iota_{w,E}:T\longrightarrow W\otimes W
      \subset H^2(A\times A,\Q).
\end{equation}
The adjective ``standard'' below refers to this actual map, not to
an unspecified nonzero embedding into the cohomology of an abelian
variety; see \cite[Proposition~6.3]{VanGeemenKS} for the full
even-Clifford construction.

\begin{theorem}\label{ext:criterion}
Let $S$ be a projective K3 surface.  Suppose that, for one rational
nonisotropic $w\in T_S$, one rational polarization, and one full
even-Clifford realization, the tensor of $j_w(v)(a)=vaw$ is induced
by a cycle in $\CH^2(S\times A^2)_{\Q}$.
Then rational HC and rational GHC hold in every degree of every
self-power of $S$.
\end{theorem}

Section~\ref{ext:section} proves this criterion independently of the
concrete K3 family constructed later.

\subsection{The quadratic locus}

Let $\mathbb U$ be the hyperbolic plane with Gram matrix
$\left(\begin{smallmatrix}0&1\\1&0\end{smallmatrix}\right)$, and put
\[
 P=\mathbb U\oplus D_8(-1)\oplus D_4(-1),\qquad
 V_P=(\mathbb U\oplus\mathbb U(2)\oplus D_4(-1))_{\Q}
     \simeq\mathbb U_{\Q}^{\oplus2}\perp\langle-1\rangle^4.
\]
The form $V_P$ has dimension eight and signature $(2,6)$.
Lemma~\ref{main:lattice} relates it to primitive copies of $P$ in
the K3 lattice.  An \emph{ample $P$-polarization} on $S$ means a
primitive embedding $P\hookrightarrow\NS(S)$ whose image contains
an ample class.  At a Picard jump its fixed perpendicular still has
dimension eight, although it is then larger than $T_S$.

\begin{theorem}\label{main:theorem}
Let $S$ be a smooth projective complex K3 surface.  Suppose that
there is a rational isometric embedding
\[
 (T_S,q)\hookrightarrow V_P.
\]
For any full even-Clifford realization $H^1(A,\Q)=C^+(T_S,q)$, any
rational nonisotropic $w\in T_S$, and any rational polarization $E$
of this weight-one Hodge structure, there is a cycle
\[
 \Gamma_S\in\CH^2(S\times A\times A)_{\Q}
\]
whose action on every $x\in T_S$ is exactly $\iota_{w,E}(x)$ in
$H^2(A\times A,\Q)$.  The same assertion holds on a transported
isogenous model with its specified tensor.

Moreover, the rational Hodge conjecture and the rational generalized
Hodge conjecture hold in every cohomological degree of every self-power
$S^m$, for every integer $m\ge1$.
\end{theorem}

In particular the theorem applies to every ample $P$-polarized K3
surface, including all Picard jumps.  Its hypothesis concerns the
actual transcendental quadratic space with its cup-product form.
Standard motivic decompositions also extend these HC and GHC conclusions
to all self-powers of Hilbert schemes of points and to the smooth
projective moduli spaces specified in Corollary~\ref{transfer:moduli}.

\subsection{Context and proof strategy}

Geometric constructions for special K3 families include Paranjape's
curve-product and Prym construction for six-line double planes
\cite[Sections~2--4]{Paranjape}, and Bolognesi--Laterveer's treatment
of order-three non-symplectic automorphisms whose fixed locus contains
a curve \cite[Theorem~4.4 and Section~6]{BolognesiLaterveer}.
Floccari and Fu explicitly distinguish an algebraic embedding from
the prescribed Kuga--Satake map
\cite[Conjectures~2.4 and~2.6]{FloccariFuOG6}; they prove algebraicity
of the latter for the K3 surfaces associated with projective singular
OG6-varieties \cite[Corollary~4.7]{FloccariFuOG6}.

For the lattice $P$ used here, Poon computes the Clifford algebra and
constructs a modular map to quaternionic abelian eightfolds
\cite[Theorem~3.1.4 and Section~4]{Poon}.
Van Geemen and Verra study related quaternionic Prym eightfolds and
algebraic quaternionic classes \cite[Sections~3--4]{VanGeemenVerra}.
These modular and quaternionic-class constructions do not by themselves
give the prescribed K3-to-abelian correspondence.
The Clifford and half-spin descriptions are classical
\cite[Proposition~6.3 and Example~6.6]{VanGeemenKS}; our recovery
argument keeps the full even module and both tensor orientations.

For powers of a K3 surface, Zarhin's description of its Hodge group
reduces tensor Hodge classes to classical invariant theory
\cite{Zarhin1983}.  When the endomorphism field of the transcendental
Hodge structure is CM, ordinary HC is already known for every
self-power \cite[Corollary~2.3]{VarescoPowers}.
The totally real case also requires relative determinant tensors;
Varesco identifies these generators and explains why classes from
the square alone need not suffice \cite[Theorem~2.16]{VarescoPowers}.
For specified four-dimensional families with generic Picard number
$16$, he also derives ordinary HC for all powers from algebraicity
of the Kuga--Satake correspondences \cite[Theorem~0.2]{VarescoPowers}.

\emph{From the exact tensor to HC and GHC.}
Section~\ref{ext:section} proves the general criterion from its one
exact standard tensor.  Alternated compositions of $j_w$ inject
$(\bigwedge^kT_S)(k)$ into the same endomorphism algebra: Clifford
PBW and invertibility of the right factor prove injectivity.
Algebraic adjoints and a polynomial inverse give return
correspondences from $H^2(A^2)$ to $\bigwedge^kT_S$ with target
support of codimension $k-1$.
The returns first algebraize relative volume tensors in the totally
real case; together with
the endomorphism-field actions and classical invariant theory, these
give ordinary HC.  A separate highest-weight argument bounds the
number of exterior factors needed to carry a rational Hodge
substructure in terms of its Hodge coniveau.  Their returns, followed
by an HC-realized map, then supply the required geometric support.
This proves GHC for whole substructures, not just for Hodge classes.

\emph{A curve source and its K3 comparison.}
Sections~\ref{geom:section} and \ref{domination:section} construct
double covers $X$ of abelian surfaces.  Their resolved inversion
quotients $Y$ carry a rank-eight variation $U$ with quadratic form
$V_P$ and an open period image.  Resolved pullback identifies it
with $U_X\subset H^2(X)$, multiplying the cup form by two.
On a nonempty geometric open, a Gauss-pencil construction gives a
product of curves dominating $X$;
when $U_X$ has no nonzero rational $(1,1)$-class, it supplies an
algebraic injection
\[
 U_X\hookrightarrow H^1(C_1,\Q)\otimes H^1(C_2,\Q).
\]
Section~\ref{pencil:section} realizes $U$, with its $Y$-form, in an
auxiliary K3 family $S^0$ over the original pencil line.  The larger
integral anti-invariant lattice is retained to control vanishing roots
and descent.  The remaining task is to make the specified comparison
$U\to U_X$ algebraic; the curve domination alone does not do so.

Section~\ref{comp:section} accomplishes this by two quadratic
substitutions.  The first, $t=s+u^2$, gives uniruled threefold models
on both sides, whose third cohomology has algebraic splittings through
curve cohomology with a Tate twist.  The required degree-three Hodge comparison is
therefore algebraic.  The second, $s=b+v^2$, produces a conic with
two components over $t=b$.  Their difference defines a cylinder that
returns this comparison to the surface fibers.  Its class lies in
the Leray term controlled by degree-three cohomology.  Fixing the first
models and correspondence before $b$ varies makes the discarded
boundary's Hodge structures fixed, so nonconstant simple periods
eliminate the boundary error.  Composition with the curve injection
now gives an algebraic curve source for $T_{S_b^0}=U_b$.

\emph{The prescribed tensor and all periods.}
Section~\ref{final:section} upgrades a curve injection at a full
special-orthogonal period to every rational Hodge map into the full
even-Clifford tensor square.  Among the nontrivial constituents,
weight-one effectivity leaves the two half-spin representations;
a multiplicity-space calculation and
rational Hodge-Hom descent recover the prescribed tensor.
The auxiliary K3 need not be $P$-polarized.
In Section~\ref{specialization:section}, a rational change of period
coordinates puts the construction into any local ample-$P$ target
ball without altering its algebraic pencils.  Successive choices
give a target seed with $\NS=P$; Buskin's theorem transports the
auxiliary curve source there, and the half-spin argument realizes the
fixed target tensor.  Proper cycle parameters then extend that same
action on the whole rank-eight perpendicular of $P$, in the full
degree-two target, throughout the ball.  Only afterward does
Section~\ref{transfer:section} compress to the actual $T_S$ at a
Picard jump.  A separate rational period and Buskin transfer reaches
the theorem's full rational quadratic-embedding locus.

\subsection{Conventions and the lattice}

Cohomology means rational Betti cohomology, except when an integral
lattice is specified; integral cohomology is taken modulo torsion.
For smooth projective varieties $M,N$, a cycle
$Z\in\CH^r(M\times N)_{\Q}$ acts by
\[
 Z_*x=(p_N)_*(\cl(Z)\smile p_M^*x),\qquad
 Z_*:H^k(M)\longrightarrow H^{k+2r-2\dim M}(N).
\]
Degree-changing maps carry the corresponding Tate twist as Hodge maps.
We use only cohomological projectors represented by cycles, not
idempotence modulo rational equivalence.  The geometric $U$ below is
distinct from the hyperbolic plane $\mathbb U$; the auxiliary K3 is
$S^0$, the double-cover surface is $X$, and the final K3 is $S$.

\begin{lemma}\label{main:lattice}
There is a primitive embedding of $P$ in
$\Lambda_{\mathrm{K3}}=\mathbb U^{\oplus3}\oplus E_8(-1)^{\oplus2}$
with orthogonal complement
$T_P=\mathbb U\oplus\mathbb U(2)\oplus D_4(-1)$.
For every primitive embedding of $P$ in the K3 lattice, the rational
orthogonal complement is isometric to $V_P$.
\end{lemma}
\begin{proof}
For an even lattice $L$, write $A_L=L^*/L$ and
$q_L:A_L\to\Q/2\Z$ for its discriminant form.  The forms of
$D_8(-1)$ and $D_4(-1)$ are respectively the split and anisotropic
forms $u,v$ on $(\Z/2)^2$.  Their values are integral modulo $2\Z$,
so each is isomorphic to its negative.  Also
$q_{\mathbb U(2)}=u$.  Hence
\[
 (A_P,q_P)\simeq u\oplus v\simeq(A_{T_P},-q_{T_P}).
\]
An anti-isometry glues $P\oplus T_P$ to an even unimodular lattice
of signature $(3,19)$, hence to $\Lambda_{\mathrm{K3}}$, with both
summands primitive.  This gives the asserted embedding.
For another embedding, Witt's extension theorem
\cite[Corollary~7.2.8]{PetersSterk} applied to the two
copies of the nondegenerate rational space $P_{\Q}$ gives an isometry
of $\Lambda_{\mathrm{K3},\Q}$ carrying one copy to the other.  It carries
their rational perpendiculars to each other.  Finally rescaling one
isotropic basis vector identifies $\mathbb U(2)_{\Q}$ with
$\mathbb U_{\Q}$, and the coordinate description of $D_4$ gives
$D_4(-1)_{\Q}\simeq\langle-1\rangle^4$.
\end{proof}

The lattice $P$ occurs as an exact N\'eron--Severi lattice.  Choose a
period in the domain of $T_P$ outside the countably many hyperplanes
defined by integral vectors not in $P$.  Period surjectivity gives a
K3 surface with N\'eron--Severi group exactly $P$; it is projective
because $P$ contains a positive class.  Its ample cone contains a
primitive integral class of $P$.  These are standard consequences of
the K3 period and chamber theorems \cite{Huybrechts}.
\section{Hodge classes and geometric supports from one exact tensor}
\label{ext:section}

We prove Theorem~\ref{ext:criterion}.  The argument is independent of
the concrete K3 family constructed later; its only Clifford premise is
one algebraic realization of the exact standard tensor in that theorem.

\subsection{Algebraic summands and coefficient extensions}

Let $S$ be a projective K3 surface, and write
\[
 N=\NS(S)_{\Q},\qquad T=N^\perp\subset H^2(S,\Q),
 \qquad V=T(1).
\]
We denote the intersection form on $T$, and its weight-zero version
on $V$, by $q$.  Choose a point $o\in S$, a divisor basis
$D_1,\ldots,D_\rho$ of $N$, and let $(b^{ij})$ be the inverse of
$(D_i\cdot D_j)$.  The following cycles in $\CH^2(S\times S)_{\Q}$
act as the indicated cohomological projectors:
\begin{equation}\label{ext:k3-projectors}
 \pi_0=[o]\times S,\quad \pi_4=S\times[o],\quad
 \pi_N=\sum_{i,j}b^{ij}D_i\times D_j,\quad
 \pi_T=\Delta_S-\pi_0-\pi_4-\pi_N.
\end{equation}
The form on $N$ is nondegenerate by the Hodge index theorem.  In
particular, these projectors and the pairing and inverse pairing on
$T$ are algebraic.  The latter are given by intersection and the
transcendental component of the diagonal.  Throughout this section,
``projector'' means a projector on cohomology represented by a cycle;
no assertion of idempotence modulo rational equivalence is needed.

The product of $k$ copies of $\pi_T$, followed by the alternating
projector
\begin{equation}\label{ext:alternator}
 \operatorname{Alt}_k=\frac{1}{k!}
       \sum_{\sigma\in\mathfrak S_k}\operatorname{sgn}(\sigma)\sigma,
\end{equation}
is an algebraic projector $\pi_{\wedge,k}$ on $S^k$ with image
$\bigwedge^k T\subset H^{2k}(S^k,\Q)$.  Here permutations act by
their graphs, with the convention giving the ordinary permutation
of the tensor factors.  All these factors have even degree.
The projector is self-adjoint for the middle intersection pairing:
$\pi_T$ is self-transpose, and transposition replaces a permutation
by its inverse.  The induced pairing on $\bigwedge^kT$ is a nonzero
rational multiple of the exterior pairing associated with $q$,
hence is nondegenerate.

We will also use algebraic degree-one projectors on abelian
varieties.  For example, if $A$ has dimension $g$ and $a\geq2$ is an
integer, the polynomial
\[
 \pi_1^A=\prod_{\substack{0\leq j\leq2g\\j\ne1}}
       \frac{[a]^*-a^j\id}{a-a^j}
\]
in the transpose graph of multiplication by $a$ acts as the
projector onto $H^1(A,\Q)$.

An algebraic class with coefficients in a field $K\supset\Q$ will
mean an element of the scalar extension of the rational image of
the cycle-class map.  This convention makes rational descent
transparent.  If $A_{\Q}\subset H_{\Q}$ is that rational image in a
fixed cohomology group, then
\begin{equation}\label{ext:descent}
 H_{\Q}\cap(A_{\Q}\otimes_{\Q}K)=A_{\Q}
 \quad\text{inside }H_{\Q}\otimes_{\Q}K.
\end{equation}
Indeed tensoring $H_{\Q}/A_{\Q}$ with a field is faithful.  Products,
permutations and contractions by algebraic pairings preserve this
notion of algebraicity.  All varieties remain over $\C$; only the
coefficients are extended.

\subsection{Algebraic returns from exterior powers}

\begin{proposition}\label{ext:conditional}
Let $S$ be a projective K3 surface with transcendental space $(T,q)$
of rational dimension $t$, put $V=T(1)$, and let
$W=C^+(T,q)=H^1(A,\Q)$ be an even-Clifford Kuga--Satake realization.
Suppose that for one rational nonisotropic $w\in T$ the tensor form
of
\begin{equation}\label{ext:exact-kappa}
 j_w:V\longrightarrow\End(W),\qquad
 j_w(v)(a)=vaw,
\end{equation}
is induced by a cycle in $\CH^2(S\times A^2)_{\Q}$, using a
rational polarization of $W$.  Then, for every $1\leq k\leq t$,
there is an injective algebraic Hodge map
\[
 F_k:\bigwedge^kV\longrightarrow\End(W)
\]
with an algebraic left inverse on cohomology.  More precisely, put
$Z=A^2$ and $n=\dim Z$.  There are correspondences
\[
 \Gamma_k\in\CH^{k+1}(S^k\times Z)_\Q,\qquad
 R_k\in\CH^{n+k-1}(Z\times S^k)_\Q.
\]
Viewing $\bigwedge^kV$ in $H^{2k}(S^k,\Q)(k)$ and identifying
$\End(W)\simeq(H^1(A,\Q)\otimes H^1(A,\Q))(1)\subset H^2(Z,\Q)(1)$
by a rational polarization, their weight-zero actions satisfy
$(\Gamma_k)_*|_{\wedge^kV}=F_k$ and $(R_k)_*F_k=\id$.
Every rational Hodge class in $\bigwedge^kV$ is the class
of a rational algebraic cycle of codimension $k$ on $S^k$.
\end{proposition}

\begin{proof}
Choose an ample divisor class $\ell$ on $A$, and use the
polarization
\[
 e(x,y)=\int_Ax\cup y\cup\ell^{g-1},\qquad g=\dim A.
\]
Changing the rational polarization in the hypothesis preserves
algebraicity: the change on the identified dual factor is a rational
Hodge automorphism of $H^1(A,\Q)$, and hence comes from an abelian
homomorphism with rational coefficients.  The elementary argument for
this fact is given in the proof of Lemma~\ref{transfer:compression}.
Here $n=2g$; put
\[
 B_W=H^1(A,\Q)\otimes H^1(A,\Q)\subset H^2(Z,\Q).
\]
The polarization identifies $\End(W)$ with $B_W(1)$.

\emph{The Clifford calculation.}
Define $F_k$ by alternating the composition of $k$ copies of
$j_w$.  Associativity gives
\[
 j_w(v_1)\cdots j_w(v_k)(a)
       =v_1\cdots v_k\,a\,w^k.
\]
Thus, if $c_k:\bigwedge^kT\to C(T,q)$ is the antisymmetrized
Clifford multiplication, then
\begin{equation}\label{ext:clifford-evaluation}
 F_k(\alpha)(1)=c_k(\alpha)w^k.
\end{equation}
The Clifford PBW decomposition identifies $c_k$ with the degree-$k$
exterior summand, so it is injective.  Moreover
$w^{-1}=w/q(w,w)$, and right multiplication by $w^k$ is invertible.
Equation~\eqref{ext:clifford-evaluation} proves injectivity in every
degree under consideration.  For odd $k$, the left and right
Clifford multiplications are both odd and their composite still
preserves $C^+(T,q)$.

The map $j_w$ is a Hodge morphism for the left spin action;
tensor products, composition and alternation preserve this
property.  Consequently
\begin{equation}\label{ext:forward-twist}
 F_k:\bigwedge^kT\longrightarrow B_W(1-k)
\end{equation}
is the corresponding untwisted Hodge morphism.  In particular, the
construction uses the exact multiplication in
\eqref{ext:exact-kappa}; an unspecified algebraic copy of $T$ in
$H^2(A^2,\Q)$ would not imply
\eqref{ext:clifford-evaluation}.

\emph{Algebraic realization and degree.}
Represent the tensor form of $j_w$ by the cycle in the
hypothesis, projected by $\pi_T$ and $\pi_1^A\boxtimes\pi_1^A$.
Take its $k$-fold exterior product and project the input by
$\pi_{\wedge,k}$.  Composition in $\End(W)$ contracts $k-1$ pairs
of $W$-slots by $e$, in the appropriate tensor order.  Each such
contraction is algebraic: the cycle
\[
 (\Delta_A)_*(\ell^{g-1})\in\CH^{2g-1}(A^2)_{\Q},
\]
regarded as a correspondence to a point, realizes
$W\otimes W\to\Q(-1)$.  It lowers cohomological degree by two.
The two uncontracted slots give $Z$; any signs arising from
permuting degree-one cohomology factors are included in the cycle
coefficients.  Hence the resulting correspondence, denoted by
$\Gamma_k$, has
\begin{equation}\label{ext:forward-cycle}
 \Gamma_k\in\CH^{k+1}(S^k\times Z)_{\Q},\qquad
 (\Gamma_k)_*|_{\wedge^kT}=F_k.
\end{equation}
Indeed its degree shift is $2-2k$, which agrees with
$2((k+1)-\dim S^k)$.  Its output twist is $1-k$, as in
\eqref{ext:forward-twist}.

\emph{An algebraic left inverse.}
Let $L=\operatorname{pr}_1^*\ell+\operatorname{pr}_2^*\ell$ on $Z$.
The summand $B_W$ is primitive in degree two for $L$.  To see this,
expand $L^{n-1}$: every term acting on $B_W$ has odd degree in each
abelian factor, whereas the only possible top bidegree on $A^2$ is
$(2g,2g)$.  Thus $L^{n-1}B_W=0$.  The Hodge--Riemann bilinear
relations imply that
\[
 Q_Z(x,y)=\int_Zx\cup y\cup L^{n-2}
\]
is a polarization on $B_W$, up to a fixed sign.  Retwisting
\eqref{ext:forward-twist} gives a morphism of weight-two Hodge structures
$(\bigwedge^kT)(k-1)\to B_W$.  Thus its image is a Hodge substructure
of $B_W$.  The restriction of $Q_Z$ to this image is nondegenerate: the
associated definite Hermitian form restricts to a definite form
there.

Multiplication by $L^{n-2}$, followed by the transpose of
$\Gamma_k$ and by $\pi_{\wedge,k}$, defines an algebraic map
\begin{equation}\label{ext:adjoint}
 G_k:H^2(Z,\Q)\longrightarrow\bigwedge^kT,\qquad
 G_k\in\CH^{n+k-1}(Z\times S^k)_{\Q}.
\end{equation}
Here the cohomological degrees are
\[
 H^2(Z)\xrightarrow{\,L^{n-2}\,}H^{2n-2}(Z)
       \xrightarrow{\,{}^t\Gamma_k\,}H^{2k}(S^k).
\]
If $Q_k$ denotes middle intersection on $\bigwedge^kT$,
self-adjointness of $\pi_{\wedge,k}$ gives
\begin{equation}\label{ext:adjoint-identity}
 Q_k(G_kF_kx,y)=Q_Z(F_kx,F_ky).
\end{equation}
The pairing on the right is nondegenerate, and $Q_k$ is
nondegenerate by the discussion of \eqref{ext:alternator}.
Therefore $D_k=G_kF_k$ is invertible on $\bigwedge^kT$.

Its characteristic polynomial has nonzero constant term, so
Cayley--Hamilton supplies a polynomial $P_k\in\Q[X]$ such that
$P_k(D_k)D_k=\id$ on that summand.  Composition of cycles realizes
$P_k(D_k)$ in codimension $2k$ on $S^k\times S^k$, using
$\pi_{\wedge,k}$ as the identity of the summand.  The correspondence
\begin{equation}\label{ext:left-inverse}
 R_k=P_k(D_k)G_k\in\CH^{n+k-1}(Z\times S^k)_{\Q}
\end{equation}
satisfies $R_kF_k=\id$; in weight-zero notation it maps
$H^2(Z,\Q)(1)$ to $(\bigwedge^kT)(k)$.  This construction uses
multiplication by an ample class, transposition and a polynomial
inverse on a finite-dimensional summand.  It requires no algebraic
inverse to a Lefschetz operator.
The composite $\Gamma_k\circ R_k$, of codimension $n$ on
$Z\times Z$, is the corresponding cohomological projector onto
$F_k(\bigwedge^kT)\subset H^2(Z,\Q)$.

Finally, let $\alpha\in(\bigwedge^kV)^{0,0}_{\Q}$.  Its image is a
rational $(1,1)$ class in $B_W$, so Lefschetz's $(1,1)$ theorem gives
a divisor class on $Z$ representing it.  Applying $R_k$ yields
$\alpha$.  The resulting cycle has codimension
$1+(n+k-1)-n=k$ on $S^k$, as asserted.
\end{proof}

\subsection{Geometric supports for exterior powers}
\label{ghc:subsec:exterior-coniveau}

The same return that algebraizes exterior Hodge classes also places
each whole exterior power on a fixed geometric support.  We use
geometric coniveau $N^p$ as defined in the introduction and set
$N^pH^j(X,\Q)=H^j(X,\Q)$ for $p\le0$.  Recall that a sub-Hodge
structure has Hodge coniveau at least $p$ when all its types $(a,b)$
satisfy $a,b\ge p$.

\begin{lemma}[Support of a correspondence]\label{ghc:lem:cycle-support}
Let $X,Y$ be smooth projective varieties, with $\dim X=N$ and
$\dim Y=D$, and let $C\in\CH^{N+p}(X\times Y)_\Q$, where $p\ge0$.
There is a fixed closed subset $B\subset Y$ of codimension at least
$p$ supporting the image of $C_*$ on cohomology.  In particular,
\[
 C_*H^j(X,\Q)\subset N^pH^{j+2p}(Y,\Q).
\]
\end{lemma}
\begin{proof}
Choose a finite cycle representative for $C$.  Each component has
dimension $D-p$, so the union $B$ of its proper images in $Y$ has
dimension at most $D-p$.  The restriction of the correspondence to
$X\times(Y\setminus B)$ is zero.  Compatibility of its action with
restriction proves the assertion, including for rational linear
combinations of components.
\end{proof}

\begin{proposition}[Exterior-power coniveau]
\label{ghc:prop:exterior-coniveau}
Let $S$ satisfy the exact even-Clifford algebraicity hypothesis of
Proposition~\ref{ext:conditional}, and put $t=\dim_\Q T$.
For every $1\le k\le t$,
\[
 \bigwedge^kT\subset N^{k-1}H^{2k}(S^k,\Q).
\]
Moreover every rational sub-Hodge structure of this exterior power
satisfies the generalized Hodge conjecture in its ambient
$H^{2k}(S^k,\Q)$.  In particular, the top exterior line is supported
in codimension~$t$.
\end{proposition}
\begin{proof}
Retain $Z=A^2$ and $n=\dim Z$ from
Proposition~\ref{ext:conditional}.  Equation~\eqref{ext:left-inverse}
gives
\[
 R_k\in\CH^{n+k-1}(Z\times S^k)_\Q,
 \qquad R_kF_k=\id\quad\text{on }\bigwedge^kT.
\]
The support of this cycle has dimension
$n+2k-(n+k-1)=k+1$.  Lemma~\ref{ghc:lem:cycle-support} therefore
puts its entire image from $H^2(Z,\Q)$ in coniveau $k-1$,
and that image contains $\bigwedge^kT$.

The only possible types of $\bigwedge^kT$ are
$(k+1,k-1)$, $(k,k)$ and $(k-1,k+1)$, since the extreme Hodge
spaces of $T$ are one-dimensional.  Thus a sub-Hodge structure
requiring more than coniveau $k-1$ is necessarily of type $(k,k)$.
Every rational vector in such a substructure is algebraic by
Proposition~\ref{ext:conditional}.  Choosing a finite rational basis
and taking the union of the supports of its codimension-$k$ cycles
proves coniveau $k$.  This also applies to $\bigwedge^tT$, whose
only type is $(t,t)$.  These are all possible coniveau requirements.
\end{proof}

The exterior returns now have both properties needed for the
criterion.  We first use their algebraicity on Hodge classes to
prove ordinary HC for arbitrary powers of $S$.  Then
Section~\ref{ghc:subsec:k3-ghc} combines that result with the support
property to treat arbitrary rational Hodge substructures.

\subsection{The endomorphism field at an individual surface}

The Hodge structure $T$ is simple.  Indeed, semisimplicity of
polarizable Hodge structures splits any sub-Hodge structure, and
at most one summand can contain its one-dimensional $(2,0)$ part.
Any other summand would be of type $(1,1)$, contrary to the
definition of $T$ and Lefschetz's $(1,1)$ theorem.  Thus
$E=\End_{\mathrm{Hdg}}(T)$ is a division algebra; its faithful action
on $H^{2,0}(S)$ embeds it into $\C$, so it is a number field.

We use Zarhin's Hodge-group theorem \cite{Zarhin1983} in the precise
forms restated in \cite[Section~2.1 and Theorem~2.8]{VanGeemenReal}
for the field, adjoint and totally real group, and in
\cite[Section~2.1]{VarescoPowers} for the CM group.  The field
$E$ is either totally real or CM.  The involution induced by the
$q$-adjoint is respectively the identity or CM conjugation.  If
$E$ is totally real, there is a nondegenerate symmetric
$E$-bilinear form $\phi$ on $T$ whose trace is $q$, and
\begin{equation}\label{ext:real-hodge-group}
 \Hg(V)=\operatorname{Res}_{E/\Q}\SO(T,\phi).
\end{equation}
If $E$ is CM with real subfield $E^+$, the analogous form is
Hermitian and
\begin{equation}\label{ext:cm-hodge-group}
 \Hg(V)=\operatorname{Res}_{E^+/\Q}\operatorname{U}(T,\phi).
\end{equation}
Here $\Hg(V)$ is the connected rational group generated by the
Hodge circle.  Its invariant tensors are precisely the rational
Hodge tensors of the weight-zero structure $V$.

Put $d=\dim_E T$, and choose a finite Galois field $K/\Q$ containing
the images of all embeddings of $E$.  Write
\begin{equation}\label{ext:eigenspaces}
 V_K=\bigoplus_{\eta:E\hookrightarrow K}V_\eta.
\end{equation}
In the real case, the summands are mutually orthogonal and
nondegenerate for $q$, and \eqref{ext:real-hodge-group} becomes the
product of their full special orthogonal groups after extending
scalars to an algebraic closure.  In the CM case, $q$ pairs
$V_\eta$ perfectly with $V_{\bar\eta}$ and pairs it trivially with
all other summands.  Equation~\eqref{ext:cm-hodge-group} becomes a
product of full general linear groups, one per conjugate pair,
acting on $V_\eta$ and $V_{\bar\eta}$ as standard and dual
representations.


\subsection{The ordinary Hodge assertion}

We first obtain the invariant tensors needed for ordinary HC.  The
following constructions are used only inside this proof of
Theorem~\ref{ext:criterion}.

\emph{Volume tensors and real multiplication.}
Suppose first that $E$ is totally real.  Each line
\[
 \bigwedge^dV_\eta\subset(\bigwedge^dV)_K
\]
is fixed by the group in \eqref{ext:real-hodge-group}.  Formation
of invariants of a rational algebraic group commutes with
extension of scalars.  Hence these lines belong to the scalar
extension of the rational Hodge-class space of $\bigwedge^dV$.
Proposition~\ref{ext:conditional} proves their algebraicity.
These relative volume tensors are the additional special-orthogonal
generators identified in \cite[Theorem~2.16]{VarescoPowers}.

Choose a nonzero volume tensor $\alpha_\eta$ on $V_\eta$.
Contracting $\alpha_\eta\otimes\alpha_\eta$ by $q$ in $d-1$ pairs
of slots gives
\begin{equation}\label{ext:volume-contraction}
 c_\eta q_\eta^{-1}\in V_\eta\otimes V_\eta,
       \qquad c_\eta\in K^*,
\end{equation}
where $q_\eta$ is the restricted form.  For example, in an
orthogonal basis $u_1,\ldots,u_d$ with $q(u_i,u_i)=a_i\ne0$, the
remaining tensor is a nonzero scalar multiple of
$\sum_i a_i^{-1}u_i\otimes u_i$; its scalar contains the nonzero
product $\prod_i a_i$.  This verifies in particular that the
contraction cannot vanish.

These contractions are algebraic.  They are obtained by
intersecting along the relevant diagonals and pushing forward on
powers of $S$.  Two codimension-$d$ tensors initially give
codimension $2d$ on $S^{2d}$; contracting $d-1$ pairs of K3 slots
leaves codimension $2d-2(d-1)=2$ on $S^2$.  Pairing
$q_\eta^{-1}$ with $q$ in one slot gives the orthogonal projector
$p_\eta$ onto $V_\eta$ in $V_K$, so that projector is algebraic as
well.  For $a\in E$, its action is
\begin{equation}\label{ext:e-action}
 a=\sum_\eta\eta(a)p_\eta.
\end{equation}
This action is rational; \eqref{ext:descent}, applied to its
cohomological correspondence class on $S^2$, proves rational
algebraicity.

\emph{CM endomorphisms and pairings.}
Now suppose $E$ is CM.  Its norm-one elements span $E$ over $\Q$.
Here is an elementary verification.  Their rational span $F$ is a
unital subalgebra of the finite field extension $E$, closed under
multiplication, and is therefore a field.  If $\bar s=-s$, then
\[
 u=\frac{1+s}{1-s},\qquad u\bar u=1,\qquad
 s=\frac{u-1}{u+1}\in F.
\]
The denominators are nonzero, and $u+1=2/(1-s)$.  Fixing a nonzero
such $s$, every $a\in E^+$ satisfies $as\in F$, hence
$a=(as)/s\in F$.  Thus $F=E$.

Each norm-one element is a Hodge isometry of $T$.  Extending it by
the identity on $N$ gives a rational Hodge isometry of
$H^2(S,\Q)$, which is algebraic by
\cite[Theorem~1.1]{Buskin}.  Rational linear combinations and
projection by $\pi_T$ give every endomorphism in $E$.  Choose a
primitive element $b$ of $E/\Q$.  The formulas
\[
 p_\eta=\prod_{\tau\ne\eta}
       \frac{b-\tau(b)\id}{\eta(b)-\tau(b)}
\]
then realize the individual eigenprojectors with coefficients in
$K$.  Applying $p_\eta\otimes p_{\bar\eta}$ to the inverse metric
gives the separate standard--dual pairing tensors.  This argument
requires that the \emph{endomorphism field} be CM; it does not
require that $T$ be one-dimensional over that field.


\emph{Invariant generation and K\"unneth assembly.}
First work over an algebraic closure of $K$.  The first
fundamental theorem for a special orthogonal group on a standard
space with nondegenerate symmetric form says that its invariant
tensors are generated by the metric and a volume tensor.  The
first fundamental theorem for a general linear group on copies of
its standard space and dual says that invariant tensors are
generated by the standard--dual pairings.  In the multilinear
formulation the number of standard and dual slots must agree,
and complete pair contractions span the invariant space.  These
statements apply to arbitrary numbers of slots; there is no stable
range restriction on the spanning assertion
\cite[Sections~4.2 and 10.2]{KraftProcesi}.  A tensor formulation of
the orthogonal statement is also given in
\cite[Theorem~2.1, Remark~2.2 and Corollary~4.2]{LehrerZhangSO}:
the top exterior tensor is harmonic and has nonzero norm, and
adjoining it to the ordinary orthogonal contractions gives all
special orthogonal invariants.

For each tensor power of $V$, use the algebraic projectors of
the preceding construction to decompose every slot
according to \eqref{ext:eigenspaces}.  Invariants of the product
of the independent groups in
\eqref{ext:real-hodge-group} or \eqref{ext:cm-hodge-group} are then
obtained factor by factor.  In the real case all metric tensors
and each relative volume tensor are algebraic by
the preceding construction.  In the CM case the
full groups are general linear, so the algebraic pairings already
suffice.  Products, permutations and the algebraic duality on
$T$ therefore realize every invariant tensor.  Scalar extension
of invariant spaces and \eqref{ext:descent} show that every
rational Hodge tensor in $V^{\otimes r}$ is algebraic over $\Q$,
as a codimension-$r$ class on $S^r$.

Finally apply the algebraic projectors
\eqref{ext:k3-projectors} to every factor of $S^m$.  The only
non-Tate summand of $H^*(S,\Q)$ is $T$ in degree two.  Every Hodge
class in a K\"unneth component is therefore a sum of Hodge tensors
in $T^{\otimes r}$ tensored with point, divisor and unit classes.
If a component has $r$ transcendental slots, $a$ divisor slots
and $b$ top-degree slots, then its degree is
$2r+2a+4b$ and the cycles constructed above have codimension
$r+a+2b$.  They thus give all classes in the required codimension
on $S^m$.


This proves the ordinary Hodge assertion of Theorem~\ref{ext:criterion}.
The only Clifford input in this deduction was the single map in its
hypothesis, through Proposition~\ref{ext:conditional}.

\subsection{Generalized Hodge conjecture for K3 powers}
\label{ghc:subsec:k3-ghc}
\label{ghc:subsec:orthogonal-ghc}

To pass from exterior powers to arbitrary Hodge substructures, we
combine their return maps with ordinary HC on products of powers.
The latter algebraizes maps between actual rational summands, so
no rational splitting of a complex middle-exterior representation
is required.  The Hodge level of a nonzero pure Hodge structure is
$\max\{|a-b|:H^{a,b}\ne0\}$.

Let $H\subset H^{2d}(S^m,\Q)$ be an actual simple sub-Hodge
structure of Hodge coniveau at least $p$, with $0\le p\le d$.
Then $H(d)$ has Hodge level at most $2(d-p)$.  We will realize
each nontrivial $H(d)$ as a quotient of a tensor product of
$c\le d-p$ exterior powers of $V=T(1)$.  The number $c$, rather
than the sum of their exterior degrees, is what controls support.
We will show that the product of the $c$ return correspondences,
followed by an ordinary-HC realization of the quotient map, has
target support of codimension at least $d-c$.  The next two lemmas
construct the required quotient with this bound on $c$; ordinary
HC handles trivial $H(d)$.  The assembled proof below gives the
correspondence calculation.

\begin{lemma}[Orthogonal exterior carriers]
\label{ghc:lem:orthogonal-carriers}
Let $V=T(1)$ for a projective K3 surface with
$\operatorname{Hg}(V)=\SO(V,q)$ and $t=\dim V\ge3$.
Let $H$ be a nontrivial rational simple representation occurring
in a direct sum of tensor powers of $V$, endowed with its induced
weight-zero Hodge structure.  If its Hodge level is at most $2b$,
there are integers $c,s\ge1$ and $k_1,\ldots,k_c\ge1$ such that
\[
 c\le b,\qquad s=\sum_i k_i,\qquad k_i\le t,
 \qquad \bigotimes_{i=1}^c\bigwedge^{k_i}V\longrightarrow H
\]
is a surjective rational Hodge map.
\end{lemma}
\begin{proof}
Over $\C$, choose orthogonal torus coordinates so that the Hodge
cocharacter has vector weights $1,0,\ldots,0,-1$.  The integral
highest weights of tensor representations of $\SO_t$ have the form
\[
 \begin{cases}
 \lambda_1\ge\cdots\ge\lambda_l\ge0,&t=2l+1,\\
 \lambda_1\ge\cdots\ge\lambda_{l-1}\ge|\lambda_l|,&t=2l.
 \end{cases}
\]
Only integral tensor weights occur here; half-integral spin weights
are not representations in tensor powers of the standard module.
For an irreducible constituent of highest weight $\lambda$, the
largest value of the Hodge cocharacter on its weights is
$\lambda_1$: the cocharacter is dominant, and all weights lie in
the convex hull of the Weyl orbit of $\lambda$.  The largest value
is attained at $\lambda$.  The orthogonal Weyl group also gives
the value $-\lambda_1$.  Hence this constituent has Hodge level
exactly $2\lambda_1$.

Put $c=\lambda_1$.  In odd dimension take the columns of the
partition $(\lambda_1,\ldots,\lambda_l)$.  Their number is $c$.
The tensor product of their exterior highest-weight vectors is a
highest-weight vector of weight $\lambda$.  Complete reducibility
therefore places the desired irreducible representation in the
tensor product of these $c$ exterior powers.  In even dimension use
the columns of $(\lambda_1,\ldots,\lambda_{l-1},|\lambda_l|)$.
For columns of height $l$, use the highest-weight vector of the
appropriate sign in $\bigwedge^lV_\C$; the two available weights
are $(1,\ldots,1,1)$ and $(1,\ldots,1,-1)$.
Their sign is chosen to agree with $\lambda_l$.  Shorter columns
are unchanged.  The same highest-weight argument applies.  We keep
the full exterior powers, including both middle-exterior
constituents, as the carriers.

Choose a complex irreducible constituent $H_\lambda$ of $H_\C$.
It has $c=\lambda_1\le b$.  The preceding construction defines a
rational representation, for $\mathbf k=(k_1,\ldots,k_c)$,
$C_{\mathbf k}=\bigotimes_i\bigwedge^{k_i}V$ whose complexification shares
this constituent with $H_\C$.  Consequently
\[
 \Hom_{\SO(V,q)}(C_{\mathbf k},H)\otimes_\Q\C
 =\Hom_{\SO(V_\C,q)}((C_{\mathbf k})_\C,H_\C)\ne0.
\]
This scalar-extension identity follows by taking invariants in
$(C_{\mathbf k})^*\otimes H$.  Any nonzero rational map in this space is
surjective by rational simplicity of $H$, and it is a Hodge map
because the Hodge cocharacter lies in $\SO(V,q)$.
This proves the required rational assertion without descending
an individually chosen complex projector.

For $t=3$ the argument is the single-coordinate case, with columns
of height one.  For $t=4$ the coordinates are
$\lambda_1\ge|\lambda_2|$, and columns have heights one and two;
the two signs in the latter give the same construction.  Thus the
small orthogonal ranks cause no exception.
\end{proof}

\begin{lemma}[Exterior carriers with endomorphism fields]
\label{ghc:lem:k3-carriers}
Let $V=T(1)$ for a projective K3 surface and let
$G=\operatorname{Hg}(V)$.  If a nontrivial rational simple
representation $H$ occurring in a direct sum of tensor powers of
$V$ has induced Hodge level at most $2b$, then there are
$1\le c\le b$ and $1\le k_1,\ldots,k_c\le\dim_\Q V$ with a
surjective rational Hodge map
\[
 \bigotimes_{j=1}^c\bigwedge^{k_j}V\longrightarrow H.
\]
\end{lemma}
\begin{proof}
Put $E=\End_{\mathrm{Hdg}}(T)$ and $a=\dim_E T$.
We use the full groups in \eqref{ext:real-hodge-group} and
\eqref{ext:cm-hodge-group}.  Choose an algebraic closure of $\Q$
inside $\C$ over which these groups split.  Highest-weight
constituents can be described over that closure and then extended
to $\C$.  The multiset of constituents of $H$ is invariant under
the absolute Galois group because $H$ is rational.  We shall use
this rationality before applying any Hodge-level bound.

\emph{The possible orthogonal ranks.}
In the totally real case, $a\ge3$.  Indeed
$\End_G(V)=\End_{\mathrm{Hdg}}(V)=E$, because commuting with the
Hodge circle is equivalent to commuting with its rational Zariski
closure.  This equality remains true after scalar extension, by
scalar extension of invariants.  Put $e=[E:\Q]$.  If $a=1$, the
split orthogonal group is trivial, so its commutant on $V$ has
dimension $e^2>e$; here $e=\dim_\Q V\ge2$ because the two extreme
K3 Hodge lines are nonzero.  If $a=2$, each split standard orthogonal
module is the sum of two distinct characters of $\SO_2$, so the
commutant has dimension at least $2e>e$.  Both contradict the
dimension $e$ of $E$.  Thus the orthogonal factors below have the
ranks covered by Lemma~\ref{ghc:lem:orthogonal-carriers}.

\emph{Totally real widths and columns.}
Suppose first that $E$ is totally real.  Over $\C$ the group is
$\prod_{\sigma:E\hookrightarrow\C}\SO(V_\sigma,q_\sigma)$.
There is one distinguished embedding $\sigma_0$ for which
$V_{\sigma_0}$ has the types $(1,-1)$ and $(-1,1)$; the other
summands have type $(0,0)$.  Accordingly the Hodge cocharacter is
trivial on the other factors and has vector weights
$1,0,\ldots,0,-1$ on the distinguished factor.

Fix one complex irreducible constituent of $H_\C$, with tuple of
orthogonal highest weights $\lambda_\sigma$, and write
$w_\sigma=\lambda_{\sigma,1}$ in the coordinates of
Lemma~\ref{ghc:lem:orthogonal-carriers}.  Galois acts transitively
on the embeddings of $E$ and transports the actual standard
eigenspace of each orthogonal factor to the actual standard
eigenspace at the conjugate embedding.  The conjugacy class of
cocharacters with weights $1,-1,0,\ldots,0$ on that standard module
is therefore preserved.  The maximum of this class on a constituent
is its first highest-weight coordinate, so Galois transports
$w_\sigma$ to the corresponding width at $\sigma_0$.  This argument
also covers type $D_4$: a triality automorphism exchanging the vector
and a half-spin module does not preserve the actual standard module
and cannot arise from this rational eigenspace transport.  That
conjugate constituent still occurs in $H_\C$.
The Hodge-level calculation in
Lemma~\ref{ghc:lem:orthogonal-carriers} consequently gives
$w_\sigma\le b$ for every $\sigma$.

Set $c=\max_\sigma w_\sigma$.  For each factor form its signed
orthogonal columns as in that lemma, with heights
$h_{\sigma j}$ for $1\le j\le c$, padding with zeros.  In the
even case the sign of a middle column is retained.  Set
$k_j=\sum_\sigma h_{\sigma j}$.  The exterior decomposition
of a direct sum includes
\[
 \bigotimes_\sigma\bigwedge^{h_{\sigma j}}V_\sigma
       \ \subset\ \bigwedge^{k_j}V_\C.
\]
Taking the Cartan components over all columns shows that
$\bigotimes_j\bigwedge^{k_j}V_\C$ contains the chosen tuple
of irreducibles.  Columns from different embeddings thus occupy
the same exterior carrier; their widths are maximized, not added.
The ranks three and four are included by
Lemma~\ref{ghc:lem:orthogonal-carriers}.

\emph{CM widths and signed columns.}
Now suppose that $E$ is CM.  Choose one embedding $\eta$ from
each conjugate pair and write
\[
 G_\C=\prod_\eta\GL(W_\eta),\qquad
 V_\C=\bigoplus_\eta(W_\eta\oplus W_\eta^*),
 \qquad \dim W_\eta=a.
\]
Choose the representatives so that the distinguished embedding
containing $V^{1,-1}$ is represented by a standard module.  On
this factor the Hodge cocharacter is conjugate to
$\operatorname{diag}(z,1,\ldots,1)$; on the other factors it is
trivial.  Its action on the distinguished dual is the inverse.
A complex irreducible constituent of $H_\C$ has, on each factor,
an integral highest weight
\[
 \lambda_{\eta,1}\ge\cdots\ge\lambda_{\eta,a}.
\]
These are highest weights for the full general linear group:
negative entries and central determinant characters are retained.
The extrema of the distinguished cocharacter on a factor of
highest weight $\lambda_\eta$ are
$\lambda_{\eta,1}$ and $\lambda_{\eta,a}$, by the weight-polytope
description and the permutation Weyl group.

Galois acts transitively on the conjugate embedding pairs.  Since
the decomposition of $V$ comes from its rational $E$-action, this
action transports standard modules to standard modules or their
duals.  On highest weights the latter possibility replaces
$(\lambda_1,\ldots,\lambda_a)$ by
$(-\lambda_a,\ldots,-\lambda_1)$; it introduces no additional
central twist.  Transporting each factor to the distinguished
one, rationality of $H$ and its level bound therefore imply
\[
 w_\eta:=\max\{\lambda_{\eta,1},
                  -\lambda_{\eta,a},0\}\le b.
\]
In particular this bound controls signed central weights as well
as differences of adjacent weights.

Set $c=\max_\eta w_\eta$.  For $1\le j\le c$ put
\[
 h^+_{\eta j}=\#\{i:\lambda_{\eta,i}\ge j\},\qquad
 h^-_{\eta j}=\#\{i:\lambda_{\eta,i}\le-j\},\qquad
 k_j=\sum_\eta(h^+_{\eta j}+h^-_{\eta j}).
\]
The positive and negative columns can be placed in the same
exterior power, because its direct-sum decomposition contains
\[
 \bigotimes_\eta
 \left(\bigwedge^{h^+_{\eta j}}W_\eta
       \otimes\bigwedge^{h^-_{\eta j}}W_\eta^*\right)
       \ \subset\ \bigwedge^{k_j}V_\C.
\]
The highest weight of $\bigwedge^h W$ has $h$ initial entries
equal to $1$, and that of $\bigwedge^h W^*$ has $h$ final entries
equal to $-1$.  Summing these weights over $j$ gives exactly
$\lambda_\eta$.  Their tensor highest-weight vector therefore
produces the desired Cartan component, including its central
character.  Thus the product of these $c$ full exterior carriers
contains the chosen constituent.  The cost is the maximum of
the positive and negative widths, not their sum.

\emph{Rational descent of the carrier.}
In both cases $c=0$ would make the chosen constituent trivial.
A rational simple representation with a trivial complex
constituent is trivial, by scalar extension of invariants, so
our nontriviality assumption gives $c\ge1$.  A factor attaining
the maximum width has a nonempty column for each $1\le j\le c$;
therefore $k_j\ge1$.  The displayed direct-sum exterior
decompositions also give $k_j\le\dim_\Q V$.

Finally, for $\mathbf k=(k_1,\ldots,k_c)$, keep the rational carrier
$C_{\mathbf k}=\bigotimes_j\bigwedge^{k_j}V$, rather than selecting any
of its complex summands.  Sharing a complex constituent yields
\[
 \Hom_G(C_{\mathbf k},H)\otimes_\Q\C
       =\Hom_{G_\C}((C_{\mathbf k})_\C,H_\C)\ne0.
\]
The equality follows by scalar extension of invariants in
$(C_{\mathbf k})^*\otimes H$.  A nonzero rational equivariant map is
surjective because $H$ is rationally simple, and it is a Hodge
map because $G$ contains the Hodge cocharacter.  This handles
nonsplit forms, Galois-conjugate constituents and their
multiplicities without assuming any complex projector rational.
\end{proof}

\begin{proof}[Proof of the generalized Hodge assertion in Theorem~\ref{ext:criterion}]
Put $V=T(1)$ and $G=\operatorname{Hg}(V)$, with the group
descriptions~\eqref{ext:real-hodge-group} and
\eqref{ext:cm-hodge-group}.  The ordinary Hodge assertion just proved gives rational HC for every
power of $S$ under the same exact-map hypothesis.  We use that conditional conclusion in
both endomorphism-field cases.

Fix a rational sub-Hodge structure
$H\subset H^{2d}(S^m,\Q)$ of Hodge coniveau at least $p$, where
$0\le p\le d$.  By semisimplicity it suffices to treat $H$ simple.
The K3 K\"unneth decomposition shows that
$H^{2d}(S^m,\Q)(d)$ is a direct sum of tensor powers of $V$:
its other factors, coming from $H^0$, the N\'eron--Severi part
of $H^2$, and $H^4$, are Tate and become trivial after the total
normalization.  Thus $H(d)$ is a rational simple representation
of $G$, of Hodge level at most $2(d-p)$.

If $H(d)$ is trivial, $H$ has type $(d,d)$.  Ordinary HC and a
finite basis of cycle classes put $H$ in
$N^dH^{2d}(S^m,\Q)\subset N^pH^{2d}(S^m,\Q)$.
Otherwise Lemma~\ref{ghc:lem:k3-carriers} supplies a tuple
$\mathbf k=(k_1,\ldots,k_c)$ and a surjective rational Hodge map
\[
 f:C_{\mathbf k}(s)\longrightarrow H(d),\qquad
 C_{\mathbf k}=\bigotimes_{i=1}^c\bigwedge^{k_i}T
       \subset H^{2s}(S^s,\Q),\quad
 s=\sum_i k_i,\quad c\le d-p.
\]
Extend $f$ by zero on a rational Hodge complement of
$C_{\mathbf k}(s)$ in $H^{2s}(S^s,\Q)(s)$, which exists by
polarizability.  Poincar\'e duality identifies this extended map to
$H(d)$ with a rational Hodge class of codimension $s+d$
on $S^s\times S^m$.  Ordinary HC for $S^{s+m}$ therefore gives
\[
 C_f\in\CH^{s+d}(S^s\times S^m)_\Q
\]
realizing it.

Retain $Z=A^2$ and $n=\dim Z$ from
Proposition~\ref{ext:conditional}.  For each column take the return $R_{k_i}$ of
\eqref{ext:left-inverse}.  Their exterior product, with factors
reordered, gives
\[
 R_{\mathbf k}\in\CH^{nc+s-c}(Z^c\times S^s)_\Q,
 \qquad (R_{\mathbf k})_*:H^{2c}(Z^c,\Q)\longrightarrow H^{2s}(S^s,\Q),
\]
whose image contains $C_{\mathbf k}$.  Indeed the tensor product of the
$F_{k_i}$ is a right inverse on $C_{\mathbf k}$.  Composition gives
\[
 C_f\circ R_{\mathbf k}\in\CH^{nc+d-c}(Z^c\times S^m)_\Q,
\]
since $(nc+s-c)+(s+d)-2s=nc+d-c$.
Its image on $H^{2c}(Z^c,\Q)$ contains $H$, and its support has
dimension
\[
 nc+2m-(nc+d-c)=2m-d+c.
\]
Lemma~\ref{ghc:lem:cycle-support} yields
\[
 H\subset N^{d-c}H^{2d}(S^m,\Q)
       \subset N^pH^{2d}(S^m,\Q).
\]
For a nonsimple substructure, take the finite union of the supports
obtained for its simple summands.  This treats every actual
occurrence, including multiplicities and substructures embedded
diagonally across isomorphic K\"unneth summands.

For the converse, put $X=S^m$ and let $B\subset X$ be any closed
support of codimension at least $p$, including a singular or
non-equidimensional support.  For each irreducible component $B_i$
choose a smooth projective resolution and let
$q_i:\widetilde B_i\to X$ be its proper map; put
$c_i=\operatorname{codim}_X B_i\ge p$.  Deligne's exactness theorem
for a proper smooth variety and a closed support, using the disjoint
union of these resolutions as a proper smooth surjective source
\cite[Proposition~8.2.7 and Corollary~8.2.8]{DeligneHodgeIII}, gives
\[
 \ker\bigl(H^j(X,\Q)\longrightarrow H^j(X\setminus B,\Q)\bigr)
   =\sum_i(q_i)_*H^{j-2c_i}(\widetilde B_i,\Q).
\]
Each proper Gysin map is a Hodge morphism
$H^{j-2c_i}(\widetilde B_i,\Q)(-c_i)\to H^j(X,\Q)$.
The source before twisting has nonnegative Hodge indices, since
$\widetilde B_i$ is smooth projective; after twisting both indices
are at least $c_i\ge p$.  Taking sums over $B$ therefore proves that
geometric coniveau $p$ has Hodge coniveau at least $p$.

Thus the two filtrations agree.  If $p>d$, the same reverse inclusion
forces $N^pH^{2d}(S^m,\Q)=0$, since a pure structure of weight $2d$
cannot have both indices at least $p$.  Cases $p\le0$ are immediate.
Odd cohomology of $S^m$ vanishes.
The result therefore holds in every degree, not only in middle
cohomology or the pure transcendental tensor summand.
\end{proof}


\section{Abelian double covers and their periods}
\label{geom:section}

We construct a rank-eight variation of K3 type with quadratic space
$V_P$ and show that its period map has rank six.  The variation occurs
in the cohomology of a double cover of an abelian surface.  A resolved
inversion quotient supplies the quadratic form that will be used in the
K3 realization.

\subsection{The double cover}

Work on a smooth finite-level chart in the moduli space of principally
polarized abelian surfaces.  Let $B$ be an indecomposable principally
polarized abelian surface, let $M=\mathcal O_B(H)$ be a symmetric line
bundle defining the principal polarization, and choose
$\chi\in\Pic^0(B)[4]$ of exact order four.  The symmetric representative
$M$ and the torsion point $\chi$ add only finite choices.  Products of
line bundles will be written without tensor symbols when convenient.

We recall the theta facts needed below.  Ample vanishing and
Riemann--Roch on $B$ give
\[
 h^0(B,M\chi)=h^0(B,M\chi^{-1})=1,\qquad h^0(B,M^2)=4,
\]
with no higher cohomology for these bundles.  The linear system $|M^2|$
is basepoint-free: the sums of opposite translates of $H$ belong to it,
and for any point of $B$ one can choose both translates to avoid that
point.  The normalized symmetric linearization of $M^2$ fixes every
section.  For example, in period coordinates its four second-order
theta sections are
$\theta[\epsilon;0](2z,2\tau)$, with
$\epsilon\in(\tfrac12\Z/\Z)^2$; reindexing the defining series by
$n\mapsto-n-2\epsilon$ proves that they are even.  Their distinct Fourier
congruence classes give a basis.  Changing the symmetric theta
representative tensors $M$ by two-torsion and leaves its square, with
this normalized linearization, unchanged.  These descriptions agree
with the standard theta constructions in
\cite{MumfordTheta,BirkenhakeLange}.

Choose $[d]\in\PP H^0(B,M^2)$, and write $D=(d=0)$.  We restrict to the
nonempty open where $D$ is smooth and avoids $B[2]$; nonemptiness follows
from Bertini and basepoint-freeness, by avoiding the finitely many
evaluation hyperplanes at $B[2]$.  Every such $D$ is invariant under
inversion.  The parameter space is thus an open in a projective
$3$-space bundle over the three-dimensional abelian-data chart, and has
dimension six.

All families below are formed on algebraic affine charts of this
projective bundle.\label{geom:section-charts}
After trivializing its rank-four section bundle, choose a nonzero
relative functional $\ell$.  On the chart $\ell(d)\ne0$ use the
actual normalized universal section $\ell(d)=1$.  On a simply
connected product polydisc, two such normalizations differ by a
nonvanishing scalar with a holomorphic square root.  Rescaling the
cover coordinate by that square root identifies the resulting covers,
their involutions, and their rational eigenspaces.

Let
\[
 \pi:B'\longrightarrow B,\qquad
 \pi_*\mathcal O_{B'}=\mathcal O_B\oplus\chi^2
\]
be the degree-two isogeny attached to $\chi^2$.  Choose a trivialization
of $\pi^*\chi^2$ and form
\begin{equation}\label{geom:double-cover}
 L=\pi^*(M\chi),\qquad
 f:X=\operatorname{Spec}_{B'}(\mathcal O_{B'}\oplus L^{-1})
       \longrightarrow B',\qquad z^2=\pi^*d.
\end{equation}
The chosen trivialization identifies $L^2$ with $\pi^*M^2$, so
$\pi^*d$ defines the multiplication in this algebra.  The surface $X$
is smooth.  Write $\delta$ for its sheet involution.

The line bundle $\pi^*\chi$ has order two, so $L$ is symmetric.
Normalize its symmetric linearization at the origin.  Its square and
the pullback of the normalized linearization of $M^2$ agree at the
origin and hence everywhere, since an automorphism of a line bundle
on an abelian variety is scalar.  Inversion therefore lifts to an
involution $j$ of $X$ commuting with $\delta$.  Put
\[
 U_X=H^2(X,\Q)^{j=1,\,\delta=-1}.
\]

\subsection{The eigenspace and its quadratic form}

\begin{proposition}\label{geom:structure}
There is a smooth double cover $Y\to K=\Km(B')$, with $K$ the resolved
Kummer surface, satisfying the following assertions.
\begin{enumerate}[label=\textup{(\roman*)}]
\item The anti-invariant lattice
\[
 V=H^2(Y,\Z)^{\delta=-1}
\]
(with torsion removed) is even, has signature $(2,14)$, and has
$2$-elementary discriminant group.  For fixed $(B,M,\chi)$, monodromy
acts trivially on this discriminant group.
\item There are eight mutually orthogonal algebraic anti-invariant
classes $a_1,\ldots,a_8$, each of square $-4$.  Their perpendicular
\[
 U=\langle a_1,\ldots,a_8\rangle^\perp\subset V_\Q
\]
is a Hodge structure of signature $(2,6)$ and, with the intersection
form inherited from $Y$, is rationally isometric to
\[
 V_P=(\mathbb U\oplus\mathbb U(2)\oplus D_4(-1))_\Q.
\]
\item Resolved pullback gives a flat Hodge isomorphism
$\beta:U\xrightarrow{\sim}U_X$ with
\[
 (\beta x,\beta y)_X=2(x,y)_Y.
\]
The algebraic projector
\begin{equation}\label{geom:projector}
 e=\frac14(\Delta_X+\Gamma_j-\Gamma_\delta-\Gamma_{j\delta})
       \in\CH^2(X\times X)_\Q
\end{equation}
has image $U_X$ in degree two and is zero in every other degree.
\end{enumerate}
The constructions are compatible with variation on the indicated
smooth parameter opens.
\end{proposition}

\begin{proof}
We first determine the fibre signs of the normalized involution on $L$.
Write $B=\C^2/\Lambda$ and represent $\chi$ by a character
$c:\Lambda\to\mu_4$ of exact order four.  Then
$B'=\C^2/\Lambda'$, where $\Lambda'=\ker(c^2)$.  Choose
$\lambda\in\Lambda\setminus\Lambda'$ and put
$k=\lambda\bmod\Lambda'\in\ker\pi$.  Translation by $k$ pairs the
sixteen points of $B'[2]$.  The fibre sign of $\pi^*M$ is the same at
$x$ and $x+k$, since these points have the same image in $B$.  For the
flat factor, representing $x$ by $v$ with $2v\in\Lambda'$ shows that
its sign changes by
\[
 \frac{c(2v+2\lambda)}{c(2v)}=c(2\lambda)=-1.
\]
Thus the signs of $L$ in each pair are opposite: eight are positive
and eight negative.

The branch curve is absent near these sixteen points.  At a
positive-sign point, $j$ preserves each local sheet of $X$ and acts
by $(u,v)\mapsto(-u,-v)$.  The quotient has two ordinary double
points, whose resolutions give two disjoint $(-2)$-curves over the
corresponding exceptional curve of $K$.  At a negative-sign point,
$j$ exchanges the local sheets.  The quotient is smooth, and its map
to $B'/\{\pm1\}$ is locally $\C^2\to\C^2/\{\pm1\}$.
Blowing up this smooth source point resolves the map to $K$ and gives
a double cover branched along the exceptional curve of $K$.  These
local models construct the smooth double cover $Y\to K$.

Set $D'=\pi^{-1}D$.  It is smooth and connected, since it is ample,
and $(D')^2=16$.  Adjunction gives $g(D')=9$.  Inversion acts freely
on $D'$, so its quotient has genus five.  The branch divisor of
$Y\to K$ is the disjoint union of that quotient and the eight
exceptional curves belonging to the negative signs.  Therefore
\[
 e(Y)=2e(K)-\bigl((2-2\cdot5)+8\cdot2\bigr)=48-8=40.
\]
The cover algebra, Serre duality, and ample vanishing give
$H^1(B',L^{-1})=0$ and hence $H^1(X,\Q)=f^*H^1(B',\Q)$.
Inversion is $-1$ on this space.  Taking quotient invariants, and
using invariance of $H^1$ under the intervening blowups, gives $q(Y)=0$.
For holomorphic two-forms,
\[
 H^{2,0}(X)=H^{2,0}(B')\oplus H^0(B',L),\qquad
 H^0(B',L)=H^0(B,M\chi)\oplus H^0(B,M\chi^{-1}).
\]
The first summand is invariant under both involutions. Inversion
replaces $\chi$ by $\chi^{-1}$; the trivialization of $\pi^*\chi^2$
used to identify the pullbacks has the nontrivial deck character, so it
exchanges the two one-dimensional summands in $H^0(B',L)$. They are
sheet-anti-invariant.  Thus $p_g(Y)=2$.  Consequently
$b_2(Y)=38$, and Hodge index gives signature $(5,33)$ on $H^2(Y)$.

The sheet-invariant rational cohomology is the pullback of
$H^2(K,\Q)$, with pairing multiplied by two, and has signature
$(3,19)$.  Its perpendicular $V_\Q$ has signature $(2,14)$.
The canonical class $K_Y$ is invariant.  For $v\in V$ we have
$v\cdot K_Y=0$, and Wu's characteristic parity formula
\cite[Theorem~11.14 and Problem~14-B]{MilnorStasheff}
\[
 v^2\equiv v\cdot K_Y\pmod2
\]
shows that $V$ is even.

For the integral assertions, let
$\mathcal H=H^2(Y,\Z)/\mathrm{torsion}$ and let $\mathcal H^+$ be
its invariant lattice.  Integral Poincaré duality makes $\mathcal H$
unimodular.  Both eigensublattices are primitive, and
\[
 2x=(x+\delta x)+(x-\delta x)
 \quad\text{for every }x\in\mathcal H.
\]
The quotient $\mathcal H/(\mathcal H^+\oplus V)$ is therefore killed
by two.  Unimodular gluing identifies this quotient with each of
$A_{\mathcal H^+}$ and $A_V$, proving that their discriminant groups
are $2$-elementary.  For fixed $B'$, pullback from the fixed Kummer
surface makes the invariant rational cohomology pointwise constant.
Its integral sublattice and its discriminant group are consequently
fixed as well.  The gluing identification is equivariant, so
monodromy acts trivially on $A_V$.

At a positive-sign point, denote the two exceptional curves on $Y$
by $E_i^+$ and $E_i^-$.  They are exchanged by $\delta$, and
\[
 a_i=[E_i^+]-[E_i^-],\qquad a_i^2=-4.
\]
The eight differences are mutually orthogonal.  Parallel transport
can permute their indices or exchange the two curves, but preserves
their span, so its perpendicular $U$ is flat.  Near a degeneration
at which the branch passes through a positive-sign point, the two
labels are governed by a scalar square root; the local action on
$a_i$ changes its sign.

The same local models identify $U$ with $U_X$.  Resolving the two
quotient nodes at a positive-sign point adds precisely one
sheet-anti-invariant exceptional direction, namely $a_i$; the blowup
at a negative-sign point adds only a sheet-invariant direction.
On a common resolution of $X\dashrightarrow Y$, pullback of a class
in $U$ is orthogonal to all exceptional classes over $X$, and hence
comes from $H^2(X,\Q)$.  Its image is exactly the
$(j,\delta)=(1,-1)$ eigenspace.  The degree-two projection formula
gives
\[
 (\beta x,\beta y)_X=2(x,y)_Y.
\]
In particular, $U$ has signature $(2,6)$.  We retain its $Y$-pairing
throughout the rational quadratic-form and K3 arguments.

The correspondence \eqref{geom:projector} is the product of the
commuting projectors $(1+j)/2$ and $(1-\delta)/2$.  The sheet
involution fixes $H^0$ and $H^4$, and the calculation of
$H^1(X,\Q)$ shows that it fixes $H^1$; by duality it fixes $H^3$
as well.  Thus $e$ vanishes outside degree two and has image exactly
$U_X$ there.

We now determine the rational form on $U$ with its $Y$-pairing.
The invariant lattice
$\mathcal H^+$ contains the pullback lattice $H^2(K,\Z)(2)$ with
finite index, necessarily a power of two.  The latter has absolute
determinant $2^{22}$, so
$|\det\mathcal H^+|$ is a square in $\Q^*/\Q^{*2}$.
Unimodular gluing gives the same absolute determinant for $V$.
Since the negative index of $V$ is fourteen, its determinant is
positive.  Splitting off the eight orthogonal lines of square $-4$
therefore gives
\[
 \det(U)=1\quad\text{in }\Q^*/\Q^{*2}.
\]
For an odd prime $p$, the projectors $(1\pm\delta)/2$ split
$\mathcal H\otimes\Z_p$ orthogonally, so $V\otimes\Z_p$ is
unimodular.  The $a_i$ span a unimodular $\Z_p$-lattice, since their
Gram determinant is $(-4)^8$.  Orthogonal projection using its
inverse Gram matrix is integral over $\Z_p$.  Its perpendicular
thus gives a unimodular $\Z_p$-lattice in $U\otimes\Q_p$.

The lattice $T_P$ has rank eight, signature $(2,6)$, and determinant
$(-1)(-4)(4)=16$; it too is unimodular at every odd prime.
A unimodular quadratic form over $\Z_p$, for odd $p$, diagonalizes
with unit entries, and the Hilbert symbol of two units is $1$.
Hence both rational forms have Hasse invariant $1$ at every odd
prime.  Their real Hasse invariants are
$(-1)^{\binom62}=-1$.  Hilbert reciprocity then gives the same
invariant, namely $-1$, at $2$.  The local classification and the
rational local-to-global isometry criterion now give
$U\simeq V_P$.  For a symmetric pairing $q$ we use the quadratic
polynomial $Q(x)=q(x,x)$, whose diagonal coefficients are those of
$q$ itself in the conventions of
\cite[Equations~(3.1)--(3.2)]{PetersSterk}.  The precise Hilbert-symbol,
reciprocity, and classification results are
\cite[Theorems~A.4.4, A.4.6, 3.3.1, and~3.3.5]{PetersSterk}.

The covers, blowups, quotients, and graph correspondences are all
relative after the finite choices and restriction to the stated
smooth opens.
\end{proof}

\subsection{The character eigenspace and the period differential}

\begin{lemma}\label{geom:twisted-hodge}
Put $E=T_B(-\log D)$.  There is a flat complex identification of
$U_X$ with a character eigenspace having Hodge components
\begin{equation}\label{geom:hodge-eigensheaves}
 U_X^{p,q}\simeq
 H^q\bigl(B,\Omega_B^p(\log D)\otimes M^{-1}\chi\bigr),
 \qquad p+q=2.
\end{equation}
In particular, $U_X^{2,0}=H^0(B,M\chi)$ is one-dimensional.
Fix a nonzero translation-invariant volume form $\Omega$ on $B$, and let
$\theta\in H^0(B,M\chi)$ represent this line.  Under
\eqref{geom:hodge-eigensheaves}, the period differential is the map
on $H^1$ induced by
\begin{equation}\label{geom:derivative}
 E\xrightarrow{\ \theta\ }E\otimes M\chi
       \simeq\Omega_B^1(\log D)\otimes M^{-1}\chi.
\end{equation}
The Kodaira--Spencer map identifies the tangent space of the full
six-dimensional family $(B,M,\chi,[d])$ with $H^1(B,E)$.
\end{lemma}

\begin{proof}
Let $\alpha:\widetilde B\to B$ be the cyclic degree-four isogeny
trivializing $\chi$, with generator $g$, and let $Z\to B$ be the
untwisted double cover $w^2=d$ defined using $M$.  On
$\widetilde X=Z\times_B\widetilde B$, let $\delta$ change the sign of
$w$, and let $g$ act through the isogeny.  A trivialization of
$\alpha^*\chi$ identifies
\[
 X=\widetilde X/\langle\delta g^2\rangle.
\]
On sheet-anti-invariant cohomology, invariance under $\delta g^2$
selects exactly the $g$-eigenvalues $i$ and $-i$.  Inversion
conjugates $g$ to $g^{-1}$ and exchanges these eigenspaces.
Projection to either one therefore identifies their $j$-invariant
part with that eigenspace over $\C$.  Choose the generator so that
the selected character sheaf is $\chi$; the other convention
replaces $\chi$ by $\chi^{-1}$.  These finite-group projectors
commute with parallel transport.

The eigensheaves can be read directly from the local double-cover
equation.  If $D=(x=0)$ and $w^2=x$, the sheet-odd holomorphic forms
are generated over the base by
\[
 w;\qquad w\,dy,\quad w\,\frac{dx}{x};
 \qquad w\,\frac{dx}{x}\wedge dy
\]
in degrees zero, one, and two.  The odd generator has the transition
rule of $M^{-1}$.  Thus the sheet-odd direct image of holomorphic
$p$-forms is $\Omega_B^p(\log D)M^{-1}$; selecting the character
tensors this by $\chi$.  Its logarithmic flat differential has
local term $\tfrac12d\log x$, together with the flat connection on
$\chi$.  The Hodge decomposition of the smooth
projective finite cover, with the character projector applied, gives
\eqref{geom:hodge-eigensheaves}.

Infinitesimal changes of local coordinates for the pair $(B,D)$
give a \v{C}ech cocycle of logarithmic vector fields in $E$.
Locally, a logarithmic field $v$ satisfying $v(f)=af$ lifts on
$w^2=f$ to $v+(a/2)w\partial_w$.  This lift is unique: an
infinitesimal vertical change preserving $w^2=f$ has coefficient
$b$ with $2b=0$.  Contraction with the sheet-odd
two-form computes the Kodaira--Spencer action on the Hodge line.
The local generators just displayed identify this contraction with
multiplication by $\theta$ in \eqref{geom:derivative}.  The bundle
isomorphism there is the rank-two determinant identity
\[
 E\otimes\mathcal O_B(D)\simeq\Omega_B^1(\log D),
 \qquad \mathcal O_B(D)=M^2.
\]
The chosen finite-order line bundles have unique infinitesimal
extensions with their torsion data, so they add no tangent directions.

It remains to identify all tangent directions of our family.
Adjunction and the logarithmic normal sequence give
\[
 0\longrightarrow E\longrightarrow T_B
       \longrightarrow\mathcal O_D(D)=K_D\longrightarrow0.
\]
An infinitesimal translation preserving $D$ is zero: its flow would
lie in the stabilizer of $\mathcal O_B(D)$, the finite kernel of
the polarization isogeny.  Thus $H^0(E)=0$, and the cohomology
sequence yields
\[
 0\longrightarrow
 \frac{H^0(K_D)}{H^0(T_B)}
 \longrightarrow H^1(E)\longrightarrow
 \ker\bigl(H^1(T_B)\longrightarrow H^1(K_D)\bigr)
 \longrightarrow0.
\]
Here $h^0(K_D)=5$, $h^0(T_B)=2$, and $h^1(T_B)=4$.
The sequence
\[
 0\longrightarrow\mathcal O_B\xrightarrow{\ d\ }M^2
       \longrightarrow K_D\longrightarrow0
\]
identifies $H^1(K_D)$ with $H^2(\mathcal O_B)$.  Under this
identification both normal-sequence boundary maps are contraction
with $c_1(D)$.  Indeed, choose local frames $e_j=g_{ij}e_i$ of
$M^2$ and write $d=f_i e_i$, so $f_i=g_{ij}f_j$.
Pulling the displayed extension back along $T_B\to K_D$, the
local lift of a vector field $v$ is $v(f_i)e_i$.
Differences of these lifts, divided by $d$, are
$v(\log g_{ij})$.  The pulled-back extension therefore has
class $c_1(D)\in H^1(\Omega_B^1)$, and its connecting maps in
degrees zero and one are the asserted contractions.
In degree one the target is the one-dimensional $H^2(\mathcal O_B)$.
This contraction is nonzero and its three-dimensional kernel is exactly the tangent
space to polarized deformations of $B$.  Hence $h^1(E)=3+3=6$.

Extending $M$ imposes the same condition, since
$c_1(M)=\tfrac12c_1(M^2)$.  Among its infinitesimal extensions
the symmetric representative is unique: differences form an
$H^1(\mathcal O_B)$-torsor on which inversion acts by $-1$.
The order-four point $\chi$ and the associated isogenies extend
uniquely because torsion of the relative Picard variety is finite
étale.  Thus the root and torsion data create no additional
infinitesimal condition.

For fixed $B$, the section directions are
$H^0(M^2)/\langle d\rangle$.  Their images in $H^0(K_D)$ have zero
boundary in $H^1(\mathcal O_B)$.  On a translation vector $v$, that
same boundary is $v\mathbin{\lrcorner}c_1(D)$, and the map
\[
 H^0(T_B)\longrightarrow H^1(\mathcal O_B),\qquad
 v\longmapsto v\mathbin{\lrcorner}c_1(D)
\]
is an isomorphism because $D$ is ample.  Section directions
therefore identify with the entire quotient
$H^0(K_D)/H^0(T_B)$.  Every polarized first-order deformation lifts
to our projective bundle of sections, by ample vanishing and base
change.  The induced map on the right-hand term is the polarized
Kodaira--Spencer isomorphism.  The two identifications in the exact
sequence prove that the full Kodaira--Spencer map is an isomorphism.
\end{proof}

\subsection{Full rank of the period map}

\begin{proposition}\label{geom:period}
The period map of $U$, equivalently of $U_X$, has rank six on a
nonempty Zariski-open subset of the six-dimensional parameter space.
For a very general parameter, its rational Hodge structure is simple,
has no nonzero rational $(1,1)$-class, and has Hodge group $\SO(U)$.
\end{proposition}

\begin{proof}
Let $F=(\theta=0)$.  It is a translate of the theta curve of $B$,
hence is a smooth genus-two curve with $\Jac(F)\simeq B$.
For clarity, these facts use indecomposability as follows.
An integral singular theta curve would have normalization of
genus at most one and could not generate the abelian surface.
If the theta divisor is not integral, ampleness and square two
force two reduced elliptic components meeting once, and hence the
product polarization: each component has positive integral
intersection with $H$ and nonnegative even self-intersection.
These possibilities are excluded.  Adjunction gives genus two,
and the restriction
pairing on $H^1(B,\Z)$, obtained by cup product with the principal
theta class, is unimodular.  The induced map on Jacobians, of equal
dimension, is therefore an isomorphism.

Set $K=K_F$ and $\eta=\chi^{-2}|_F$.  Restriction
$\Pic^0(B)\to\Pic^0(F)$ is an isomorphism, so $\eta$ is nontrivial
two-torsion.  Adjunction gives
\[
 (M\chi)|_F=K,\qquad M^2|_F=K^{\otimes2}\otimes\eta.
\]
For a transverse $D$, multiplication by $\theta$ yields
\begin{equation}\label{geom:quotient-sequence}
 0\longrightarrow E\longrightarrow E\otimes M\chi
       \longrightarrow(E\otimes M\chi)|_F\longrightarrow0.
\end{equation}
We shall prove that the last sheaf has no global sections for
general $d$.  By the cohomology sequence and
Lemma~\ref{geom:twisted-hodge}, this implies injectivity of the
period differential on all six tangent directions.

Put $Z=D\cap F$, a reduced divisor of degree four.  Restricting the
logarithmic normal sequence to $F$ is exact because the intersection
is transverse.  After tensoring by $K$, its map on global sections is
\begin{equation}\label{geom:tangency-matrix}
 A_d:H^0(F,K)^{\oplus2}\longrightarrow
 H^0\bigl(Z,(K^{\otimes3}\otimes\eta)|_Z\bigr),\qquad
 v\longmapsto (v\mathbin{\lrcorner}dd)|_Z.
\end{equation}
Its kernel is $H^0(F,(E\otimes M\chi)|_F)$, and source and target
both have dimension four.  Differentiating a section at its zero
scheme is intrinsic, so $A_d$ is independent of local bundle frames.

The restriction sequence for $F$, together with ample vanishing, gives
a surjection
\[
 H^0(B,M^2)\longrightarrow H^0(F,K^{\otimes2}\otimes\eta)
\]
whose kernel is generated by $\theta\theta'$, where
$0\ne\theta'\in H^0(B,M\chi^{-1})$.  Its restriction
$t=\theta'|_F$ is a nonzero section of $K\otimes\eta$:
a section vanishing along $F$ would give a section of the nontrivial
flat bundle $\chi^{-2}$.  The bundle $K^{\otimes2}\otimes\eta$
is basepoint-free by Riemann--Roch on the genus-two curve.  We may
therefore fix a general restriction $s=d|_F$ whose four zeros are
simple and avoid the zeros of $t$.

Replacing $d$ by $d+\lambda\theta\theta'$ fixes $s$ and $Z$ and
replaces $A_d$ by $A_d+\lambda A'$.  Choose a basis $v_1,v_2$ of
translation-invariant vector fields on $B$.  The normal derivatives
\[
 \theta_i=(v_i\mathbin{\lrcorner}d\theta)|_F,\qquad i=1,2,
\]
form a basis of $H^0(F,K)$: the map to this space is injective
because the theta divisor has finite translation stabilizer, and
both spaces have dimension two.  In these coordinates,
\[
 A'(a,b)=t(a\theta_1+b\theta_2)|_Z.
\]
The canonical pencil of a genus-two curve gives an isomorphism
$\Sym^2H^0(K)\simeq H^0(K^{\otimes2})$.  Moreover,
$K^{\otimes2}(-Z)\simeq\eta$ has no section, so restriction to $Z$
is injective.  Since $t$ does not vanish on $Z$, we obtain
\[
 \im A'=t\,H^0(K^{\otimes2})|_Z,\qquad \rank A'=3.
\]
The kernel is spanned by the skew gradient
$k=(\theta_2,-\theta_1)$.  The canonical pencil has no base point,
so this is a nowhere-zero tangent section of
$T_F\otimes K\simeq\mathcal O_F$.

Under $T_F\otimes K\simeq\mathcal O_F$, the vector $k$ is a
constant $c\ne0$.  With the adjunction convention
$\operatorname{Res}(\Omega/h)$, where $\theta=h e$, this constant is
$c=-\Omega(v_1,v_2)$: writing $\Omega=J\,dh\wedge du$ identifies
$k$ with $-\Omega(v_1,v_2)\partial_u\otimes du$.
Consequently $A_d(k)=c\,d_Fs|_Z$.  Here $d_Fs|_Z$ is an intrinsic section of
$(K^{\otimes3}\otimes\eta)|_Z$.  To show that this vector is outside
$\im A'$, consider the sequence
\[
 0\longrightarrow\eta\xrightarrow{\ \cdot s\ }K^{\otimes2}
       \longrightarrow K^{\otimes2}|_Z\longrightarrow0,
\]
where $\eta^{\otimes2}\simeq\mathcal O_F$ is understood.  The
connecting class of $(d_Fs)/t|_Z$ pairs by Serre duality with
$t\in H^0(K\otimes\eta)$.  This pairing is computed by local
principal parts.  Near $z\in Z$, choose a coordinate $u$ and
an $\eta$-frame $\epsilon$ with $\epsilon^2=1$, and write
$s=f(u)(du)^2\epsilon$ and $t=g(u)du\,\epsilon$.
Here $f(z)=0$, $f'(z)\ne0$, and $g(z)\ne0$.  A local lift of
$(d_Fs)/t|_Z$ is $(f'/g)(du)^2$.  Dividing by $s$ gives the
$\eta$-valued principal part $[f'/(gf)]\epsilon$; pairing it with
$t$ gives $(f'/f)du$.  Thus
\[
 \left\langle\partial\bigl((d_Fs)/t|_Z\bigr),t\right\rangle
   =\sum_{z\in Z}\operatorname{Res}_z\frac{df}{f}=4,
\]
up to the fixed nonzero normalization of Serre duality.  Each zero
is simple.  Changing a frame adds a regular differential and leaves
the residues unchanged.  The formula concerns the four local
principal parts; it does not require a global meromorphic
differential with those four poles.  The connecting class is
nonzero, and hence $A_d(\ker A')$ is not contained in $\im A'$.

Choose source and target bases in which $A'$ has an invertible
$3\times3$ block and is zero in the remaining row and column.
The coefficient of $\lambda^3$ in $\det(A_d+\lambda A')$ is the
determinant of that block multiplied by the induced map
$\ker A'\to\coker A'$.  It is nonzero by the preceding pairing.
Thus $A_{d+\lambda\theta\theta'}$ is invertible for general
$\lambda$.  This construction shows that invertibility is a
nonempty open condition in the space of sections.  Intersecting
with the nonempty open of smooth divisors avoiding $B[2]$ gives
the required general $d$.

It follows from \eqref{geom:quotient-sequence} that
$H^1(E)\to H^1(E\otimes M\chi)$ is injective.  By
Lemma~\ref{geom:twisted-hodge}, its source is the entire
six-dimensional tangent space.  The type-IV period domain of a
quadratic space of signature $(2,6)$ has complex dimension six;
the period map is therefore a local isomorphism at these points.
The rank condition is Zariski open, since it is the rank condition
on the algebraic Kodaira--Spencer map between Hodge bundles.

Finally, on a type-IV period domain the Hodge circles rotate its
oriented positive two-planes.  These rotations generate the
identity component of $\SO(U_\R)$, and rotations from any nonempty
open have the same Zariski closure.  Hence the Hodge circles in our
open period chart are not all contained in any proper rational
algebraic subgroup of $\SO(U)$.

We recall why the exceptional conditions may be taken complex
analytic.  A proper rational algebraic subgroup stabilizes a
rational line in a tensor representation that is not stabilized by
$\SO(U)$.  Indeed, take a finite-degree representation of regular
functions containing generators of the subgroup's defining ideal,
intersect it with that ideal, and take the top exterior power.
The resulting line has precisely the required stabilizer.
Polynomial representations and semisimplicity place this
representation in a sum of tensor powers and duals of $U(1)$.
A Hodge circle stabilizing a rational line fixes
it pointwise, because it is connected and compact.  Thus it forces
a specified rational tensor to have type $(0,0)$.  Each such
condition is a closed complex-analytic Hodge locus, and it is
proper unless that tensor is fixed by all the Hodge circles in
the domain.  There are only countably many rational tensors.
Outside their proper loci, $\Hg(U)=\SO(U)$.  The standard
representation of $\SO(U)$ is absolutely irreducible and has no
invariants.  Simplicity of $U$ and the absence of nonzero rational
$(1,1)$-classes follow.
\end{proof}
\section{Coherent half sums and domination by a product of curves}
\label{domination:section}

We retain the data $(B,M,\chi,D)$ of Section~\ref{geom:section}, with
$M=\mathcal O_B(H)$ and $D\in|2H|$ smooth and disjoint from $B[2]$.
The regular Gauss fibers of $D$ consist of four opposite pairs.  Coherently
taking half the sum of one point from each pair produces two curves
mapping to $B$.  We prove that their addition map has total degree
eight and pulls $D$ back to an even divisor.  Separate covers of the
two curves then give a morphism to the surface $X$ of
Equation~\eqref{geom:double-cover}.

\subsection{The generic Gauss pencil}

\begin{lemma}\label{geom:gauss-generic}
For $(B,D)$ in a nonempty Zariski-open subset of the parameter space,
$D$ has no nontrivial translation stabilizer.  Its quotient
$C=D/\{\pm1\}$ is a smooth genus-three curve, and the connected
\'{e}tale cover $D\to C$ is defined by a nontrivial line bundle
$\epsilon\in\Pic^0(C)[2]$.  The Gauss map is
the composite
\[
 D\longrightarrow C\xrightarrow{|K_C\otimes\epsilon|}\PP^1,
\]
of degrees two and four.  No fiber of the degree-four map has all its
multiplicities even.  The good open remains nonempty on every
connected component obtained by adjoining the finite theta,
torsion, and level choices used to define $X$.
\end{lemma}

\begin{proof}
The smooth ample divisor $D$ is connected: the divisor sequence and
$H^1(B,\mathcal O_B(-D))=0$ give $H^0(D,\mathcal O_D)=\C$.
Thus it is irreducible.  Adjunction gives $2g(D)-2=D^2=8$.
Inversion acts freely because $D$ avoids $B[2]$, so Riemann--Hurwitz
gives $g(C)=3$.  Connectedness of the double cover makes $\epsilon$
nontrivial.

The curve $D$ generates $B$: containment in a translate of a proper
abelian subvariety would make $D$ an elliptic curve with square zero.
Hence $J(D)\to B$ is surjective and
invariant one-forms on $B$ restrict injectively to $D$.  They are
anti-invariant under inversion and span the two-dimensional space
$H^0(C,K_C\otimes\epsilon)$, whose dimension follows from
Riemann--Roch.  They have no common zero on $D$, since its tangent
line is a nonzero subspace of $T_0B$.  They therefore define the
basepoint-free pencil $|K_C\otimes\epsilon|$, of degree four.
Identifying its target with the projective line of tangent directions
identifies the composite with the Gauss map.

We justify that these pairs $(C,\epsilon)$ are general.  The open
projective bundle of smooth curves in $|2H|$ over the three-dimensional
level moduli of principally polarized abelian surfaces has dimension
six and maps algebraically to the moduli space $\mathcal R_3$ of
genus-three curves with nonzero two-torsion.  Its fibers are countable.
In fact, fixing $(C,\epsilon)$ fixes the double cover $D$; any map of
that curve to an abelian surface is a homomorphism from $J(D)$ followed
by a translation.  Surjective homomorphisms, up to isomorphism of
their targets, form a countable set: the connected kernel is
determined by a rational subspace of $H_1(J(D),\Q)$, and the further
finite isogeny choices are countable.  For a fixed homomorphism,
requiring inversion on the image to be the given
involution on $D$ imposes an equation $2b=\text{constant}$ on the
translation $b$, giving finitely many choices.  Polarizations are
discrete, and $[D]=2[H]$ fixes their numerical class; the remaining
symmetric theta and level choices are finite.  A countable fiber of a
finite-type morphism over $\C$ has dimension zero.  Since
$\mathcal R_3$ is irreducible of dimension six, the map is dominant
on each component of its six-dimensional source.  Irreducibility of
$\mathcal R_3$ follows from the transitive symplectic monodromy on
nonzero two-torsion.  The Torelli theorem for curves, with their
integral principal polarization
\cite[Chapter~3, Theorem~2.2]{Huybrechts}, likewise gives
zero-dimensional fibers for the Jacobian map from $\mathcal R_3$
to the moduli space $\mathcal A_3[\epsilon\ne0]$ of principally
polarized abelian threefolds with nonzero two-torsion.  This target
is irreducible of dimension six: it is a quotient of connected
Siegel space, and the symplectic group is transitive on nonzero
vectors modulo two.  Hence the Jacobian map is dominant.

Choose a theta characteristic $\kappa$ on a genus-three curve and set
\[
 \Theta_\kappa=
 \{\alpha\in J(C):h^0(C,\kappa\otimes\alpha)>0\}.
\]
The Riemann theta theorem identifies $\Theta_\kappa$ with a symmetric
theta divisor; see \cite[Chapter~II, Section~3, Corollary~3.6 and
Theorem~3.10]{MumfordTheta}.  An all-even member of
$|K_C\otimes\epsilon|$ exists precisely when
some $\alpha$ with $2\alpha=\epsilon$ lies on $\Theta_\kappa$:
an effective divisor $E$ with $2E\in|K_C\otimes\epsilon|$ corresponds
to $\alpha=\mathcal O_C(E)\kappa^{-1}$.  This condition is independent
of $\kappa$: translating a symmetric theta divisor by two-torsion
permutes the set of all sixty-four halves.  The condition is closed
on $\mathcal A_3[\epsilon\ne0]$: after a finite theta-level cover,
the theta divisor and the finitely many half sections are defined,
and incidence with each section is closed.  Their finite union
descends to a closed subset.

It is a proper condition.  On a product of three elliptic curves,
choose $\epsilon$ nonzero on every factor.  Every half then has
coordinates of order four, whereas a symmetric product theta divisor
is a union of coordinate divisors at two-torsion points.  No half lies
on that divisor.  The resulting nonempty open in the moduli of
principally polarized threefolds meets the dominant Jacobian locus,
and hence pulls back to a nonempty open in our parameter space.  This
proves the assertion about fibers.  The order-four choice $\chi$ on
$B$ is independent of $D$, $C$, and $\epsilon$.  Every connected
component of its finite étale torsion/level cover surjects onto the
connected base, so the same open remains nonempty in each component.

Finally, a translation preserving $D$ belongs to the finite kernel
$K(M^2)=B[2]$.  Each nonzero element of that kernel acts
nontrivially on $\PP H^0(B,M^2)$: otherwise its commutator with all
theta-group elements would be trivial, contrary to the nondegenerate
theta-group pairing \cite{BirkenhakeLange}.  Its fixed locus is proper.
Avoiding the union of these finitely many loci gives a nonempty
open on each section fiber.  Its intersection with the preceding
open proves the lemma.
\end{proof}

\subsection{Coherence, parity, and signed relations}

Fix data in the open of Lemma~\ref{geom:gauss-generic}.
Let $R\subset\PP^1$ be the regular-value open of the Gauss pencil.
Locally on $R$ its fibers on $D$ can be written
$\{\pm x_1,\ldots,\pm x_4\}$.  A \emph{selection} means an unordered
subset containing exactly one point of each opposite pair.  The
selections form a finite étale cover $\mathcal S\to R$ of degree
sixteen.  If $s$ is a selection containing $p$, write $s[p]$ for the
selection obtained by replacing $p$ by $-p$.

This selection cover and its splitting by the parity of sign changes
form the combinatorial part of Donagi's tetragonal construction
\cite[Sections~2.1--2.2 and~2.5]{DonagiPrym}.
The argument below constructs compatible half sums in the given
abelian surface $B$ and lifts their addition map through separate
curve covers to $X$.

\begin{lemma}\label{geom:halves}
There is an algebraic map $h:\mathcal S\to B$ such that
\begin{equation}\label{geom:coherent}
 2h(s)=\sum_{p\in s}p,\qquad h(s[p])=h(s)-p.
\end{equation}
The selections split globally into two degree-eight covers
$\mathcal S_+\sqcup\mathcal S_-$ according to the parity of the number
of sign changes.  Their smooth projective completions $C_\pm$, which
may be disconnected, carry extensions $a_\pm:C_\pm\to B$ of $h$.
Wherever its differential is nonzero, either map is tangent to the
direction recorded by its parameter in $\PP^1$.
\end{lemma}

\begin{proof}
For a fixed regular fiber, a coherent assignment is a function $h$ on
its sixteen selections satisfying Equation~\eqref{geom:coherent}.
Choosing one selection $s_0$ and one half $b$ of its sum determines
it uniquely by
\[
 h(s)=b-\sum_{p\in s_0\setminus s}p.
\]
Thus coherent assignments form a torsor under the constant group
$B[2]$.  This torsor is algebraic and finite étale over $R$: start
with the cover of pairs $(s,b)$ satisfying $2b=\sum_{p\in s}p$, and
identify $(s,b)$ with $(s[p],b-p)$.  After \'{e}tale-locally labeling
the four pairs, the flips commute, are involutions, and their
equivalence classes have sixteen elements;
the resulting finite étale quotient is precisely the assignment
torsor.  Parallel transport commutes with the constant $B[2]$-action,
so its monodromy acts by translations, despite possible permutations
of the selections.

Consider a small disc about any omitted value, including infinity
if it is omitted.
At the central fiber choose one point from each pair of opposite
points of $D$, and near each chosen point take the entire cluster
specializing to it.  Their union is a monodromy-invariant selection
on the punctured disc.  Its sum extends holomorphically to the disc.
Because multiplication by two on $B$ is étale, that sum has a
holomorphic half on the simply connected disc.  The corresponding
coherent assignment is fixed by local monodromy.  A translation of a
torsor with a fixed point is the identity.  Hence the assignment
torsor has trivial monodromy about every puncture, extends étale over
$\PP^1$, and is trivial.  Any of its global sections gives the
algebraic map $h$.

Two selections have the same parity if they differ by an even number
of flips.  The two parity classes define a degree-two \'{e}tale cover
of $R$.  The invariant selection over each punctured disc fixes its
parity class, hence both classes.  This cover also extends \'{e}tale
over $\PP^1$ and is trivial.  Labeling its two sections gives the
global covers $\mathcal S_\pm$.  Properness of $B$ extends their maps
uniquely to the smooth projective completions.  Finally,
differentiation of Equation~\eqref{geom:coherent}, using the invariant
tangent space of $B$, gives locally
\begin{equation}\label{geom:half-derivative}
 da_\pm=\tfrac12\sum_{i=1}^4\sigma_i\,dx_i,
 \qquad \sigma_i\in\{1,-1\}.
\end{equation}
All the $dx_i$ lie in the same Gauss direction.  Their sum, whenever
nonzero, lies in that direction as well.  Continuity gives the same
conclusion at the omitted values wherever the extended map has
nonzero differential.
\end{proof}

The common assignment on both parities gives an exact identity for
any two selections over the same regular value:
\begin{equation}\label{geom:agreed-sum}
 h(s)+h(t)=2h(s)-\sum_{p\in s\setminus t}p
          =\sum_{p\in s\cap t}p.
\end{equation}
Here we used Equation~\eqref{geom:coherent}.  In particular,
$h(s^c)=-h(s)$ for the opposite selection $s^c$.  Changing the
common assignment by $b_0\in B[2]$ changes the left side of
Equation~\eqref{geom:agreed-sum} by $2b_0=0$, so the identity is
independent of that choice.

\begin{lemma}\label{geom:relations}
On a regular-value disc, no signed sum of a nonempty subset of the
four local branches $x_i$ is constant.  Every connected component of
$C_+$ and $C_-$ maps nonconstantly to $B$.  No image component from
$C_+$ is an opposite translate of an image component from $C_-$;
that is, their image curves cannot satisfy $Z_+=b-Z_-$ for any
$b\in B$.
\end{lemma}

\begin{proof}
A one-term relation is impossible because a branch of $D$ is
nonconstant.  A two-term relation identifies open subsets of $D$
and a translate of $D$ or $-D$.  Since $D=-D$, triviality of its
translation stabilizer forces the translate to be zero; this would
identify two of the distinct opposite pairs, also impossible.

Suppose there is a three-term relation.  The monodromy on the four
unsigned pairs is transitive, since $C$ is connected.  Analytic
continuation therefore supplies four constant relations, one omitting
each index.  Their coefficient matrix has zero diagonal and entries
$\pm1$ off the diagonal.  Let $J$ denote the all-ones matrix.  Modulo two
the coefficient matrix is $J-I$ in dimension four,
whose square is the identity over $\mathbb F_2$.  Its determinant is
therefore odd, in particular nonzero.  The corresponding homomorphism
$B^4\to B^4$ is an isogeny.  The four moving branches lie in a fixed
finite fiber of this isogeny and must all be constant, a contradiction.

Two four-term relations whose sign vectors are neither equal nor
opposite differ in between one and three positions.  Subtracting
them makes twice a nonempty proper signed partial sum constant.
That partial sum is itself locally constant, since $[2]$ has finite
fibers, contrary to the cases just proved.  Thus a four-term relation,
if it exists, determines its sign selection uniquely up to complement.
Monodromy must preserve that selection or exchange it with its
complement.
Every local monodromy has a fixed signed point: a partition of four
which is not all-even is one of $1+1+1+1$, $2+1+1$, or $3+1$, and
there is a singleton in each case.  The étale cover $D\to C$ gives
two fixed signed points above that singleton.  The selection therefore
cannot be exchanged with its opposite.  The local monodromies
generate the full monodromy, so the four chosen signed points form a
global invariant subset of the eight sheets of $D\to\PP^1$.
This contradicts connectedness of $D$.

A constant component of $a_\pm$ would give a constant four-term
relation by Equation~\eqref{geom:coherent}.  Suppose image components
from opposite parities satisfy $Z_+=b-Z_-$.  Choose a general smooth
point $z_+\in Z_+$ and its partner $z_-=b-z_+\in Z_-$, and lift
them along local inverse branches of $a_+$ and $a_-$ away from
ramification and the exceptional Gauss parameters.  The tangent
directions at these paired image points agree.
Lemma~\ref{geom:halves} identifies both directions with their Gauss
parameters, so the corresponding local selections $s_+,s_-$ lie
over the same parameter and satisfy
$h(s_+)+h(s_-)=b$.  Opposite parities have either one or three agreed
signed points.  Equation~\eqref{geom:agreed-sum} then gives one of
the excluded partial relations.
\end{proof}

\subsection{The degree-eight addition map}

\begin{proposition}\label{geom:domination}
For every parameter in a nonempty Zariski-open subset of the
parameter space, there exist smooth projective connected curves
$C_1,C_2$ and a dominant generically finite morphism
\[
 g:C_1\times C_2\longrightarrow X.
\]
Whenever additionally $U_X$ has no nonzero rational $(1,1)$-class,
an algebraic correspondence induces an injection
\[
 U_X\hookrightarrow H^1(C_1,\Q)\otimes H^1(C_2,\Q).
\]
In particular, this injection exists at a very general parameter.
\end{proposition}

\begin{proof}
Take the open of Lemma~\ref{geom:gauss-generic} and the common
coherent assignment of Lemma~\ref{geom:halves}.

\smallskip\noindent
\emph{Image classes and finiteness.}
For either parity, its eight sign vectors satisfy
\[
 \sum_\sigma\sigma_i\sigma_j=8\delta_{ij}.
\]
If $\omega$ is an invariant real two-form on $B$,
Equation~\eqref{geom:half-derivative} therefore gives, over a
regular-value disc,
\[
 \sum_{s\in\mathcal S_\pm}h(s)^*\omega
       =2\sum_{i=1}^4x_i^*\omega.
\]
The right side is exactly the sum over the eight branches $\pm x_i$
of $D$, since inversion fixes invariant two-forms.  Integrating over
the regular-value open, whose omitted points have measure zero,
proves
\begin{equation}\label{geom:image-cycle}
 (a_+)_*[C_+]=(a_-)_*[C_-]=[D]
 \quad\text{in }H_2(B,\R).
\end{equation}
Invariant forms represent all of $H^2(B,\R)$, so this equality of
cycle classes includes the degrees of all image maps and every
connected component of both sources.

Consider addition
\[
 F:C_+\times C_-\longrightarrow B,\qquad
 (p,q)\longmapsto a_+(p)+a_-(q).
\]
It is finite on every product of connected components.  Otherwise
a positive-dimensional fiber would contain a curve.  Neither
projection of that curve could be constant, since each factor map is
nonconstant and hence finite onto its image.  Both projections would
therefore be surjective, forcing the two image curves to be opposite
translates, contrary to Lemma~\ref{geom:relations}.  Properness now
gives finiteness.  Each component product is two-dimensional and its
finite image is all of $B$.  Its total degree, summed over the
component products, is, by Equation~\eqref{geom:image-cycle},
\begin{equation}\label{geom:addition-degree}
 \deg F=[D]^2=(2H)^2=8.
\end{equation}
One can obtain this formula directly by intersecting one image cycle
with a general translate of the negative of the other; the
intersection multiplicities count the fibers of addition with the
image-map degrees.

\smallskip\noindent
\emph{The pullback divisor.}
Fix a signed point $x$ of a regular Gauss fiber.  For each selection
$s\in\mathcal S_+$ containing $x$, let $\tau=s^c[-x]$.
It differs from $s$ in three signs, so $\tau\in\mathcal S_-$,
and Equation~\eqref{geom:agreed-sum} gives
\begin{equation}\label{geom:four-pairs}
 h(s)+h(\tau)=x.
\end{equation}
There are exactly four such selections $s$.  The resulting four
points of $C_+\times C_-$ are distinct, since their first selections
differ.

As $x$ varies locally on $D$, each construction defines a smooth
local curve in the source, parametrized by the common regular Gauss
parameter and mapped isomorphically to an open subset of $D$.
At its points both factor differentials have image in the tangent
line of $D$ at $x$, so $dF$ has rank at most one.  Its restriction
to the constructed local curve is nonzero.  Choose target coordinates
$(r,s)$ with $D=(s=0)$ and source coordinates $u=r\circ F,v$, with
the local curve given by $v=0$.  Then
\[
 F(u,v)=(u,G(u,v)),\qquad G(u,0)=0,\qquad G_v(u,0)=0.
\]
Thus the multiplicity of $F^*D$ along each constructed local curve
is at least two.

Write $F^*D=\sum_\alpha e_\alpha E_\alpha$, with $E_\alpha$
running through the prime divisors on all component products.
Every $E_\alpha$ dominates $D$, since $F$ is finite.  Put
$f_\alpha=[\C(E_\alpha):\C(D)]$.  The valuation degree formula
and Equation~\eqref{geom:addition-degree} give
\[
 \sum_\alpha e_\alpha f_\alpha=8.
\]
Above a general complex point of $D$, the normalization of
$E_\alpha$ has $f_\alpha$ distinct local branches.  The four
branches constructed above already contribute at least $4\cdot2$
to this sum.  They exhaust it, and each has multiplicity exactly
two.  Several of these local branches may belong to the same prime
divisor.  This proves, on every component product, the Cartier
divisor identity
\begin{equation}\label{geom:even-divisor}
 F^*D=2E_0
\end{equation}
for an effective integral divisor $E_0$; the smooth source makes
$E_0$ Cartier.

\smallskip\noindent
\emph{Lifting to the double cover.}
Choose connected components $C_1^0\subset C_+$ and
$C_2^0\subset C_-$.  Pull back $\pi:B'\to B$ separately along
their maps and choose a connected component of each pullback.  This
gives finite \'{e}tale surjective maps $C_i'\to C_i^0$ and lifts
$a_i':C_i'\to B'$.  Their sum
\[
 F'=a_1'+a_2':C_1'\times C_2'\longrightarrow B'
\]
is finite and surjective: its composite with $\pi$ is the finite
surjective map obtained from $F$ by the factor covers.  Pullback of
Equation~\eqref{geom:even-divisor} gives $F'^*(\pi^*D)=2E$ on the
smooth product $S'=C_1'\times C_2'$.

Let $s_E$ be the canonical section of $\mathcal O_{S'}(E)$ and put
\[
 N=F'^*L\otimes\mathcal O_{S'}(-E),\qquad
 u=F'^*(\pi^*d)/s_E^2\in H^0(S',N^2).
\]
The section $u$ is everywhere nonzero.  It defines an \'{e}tale
double cover
\[
 \widetilde S'=\operatorname{Spec}_{S'}
        (\mathcal O_{S'}\oplus N^{-1})\longrightarrow S'.
\]
This is the normalization of $S'\times_{B'}X$ and has an
everywhere-defined morphism to that fiber product.  Explicitly,
where $E$ has local equation $f=0$, the pulled-back equation of $X$
is $z^2=uf^2$.  Its normalization is $y^2=u$, and the morphism back
is $z=fy$.  This formula also applies at singular points and
intersections of $E$ and glues by the definition of $N$.

The \'{e}tale double cover $\widetilde S'\to S'$ corresponds to a
character
\[
 \rho:\pi_1(C_1')\times\pi_1(C_2')\longrightarrow\Z/2.
\]
It is the sum of its two factor restrictions $\rho_1,\rho_2$.
Let $C_i\to C_i'$ be the connected finite \'{e}tale cover
corresponding to $\ker\rho_i$, using the identity cover if
$\rho_i=0$.  These covers are algebraic and projective.  The
pullback of $\widetilde S'$ to $C_1\times C_2$ is trivial, so
choosing either section and composing the preceding morphisms gives
$g:C_1\times C_2\to X$.  Its composite to $B'$ is finite
surjective, so its image has dimension two.  The surface $X$ is
irreducible: its nonempty reduced branch divisor makes its
function-field square class nontrivial.  Hence $g$ is dominant and
generically finite.

\smallskip\noindent
\emph{The cohomological correspondence.}
The projection formula gives $g_*g^*=(\deg g)\id$ on rational
cohomology, so $g^*$ is injective.  Suppose $U_X$ has no nonzero
rational $(1,1)$-class.  Then it has no nonzero Hodge map to
$\Q(-1)$: dualizing such a map with a polarization would give a
rational $(1,1)$-class in $U_X$.  Thus the two outer projections of
$g^*U_X$ in the decomposition
\[
 \begin{split}
 H^2(C_1\times C_2,\Q)={}&
 \bigl(H^2(C_1)\otimes H^0(C_2)\bigr)
 \oplus\bigl(H^1(C_1)\otimes H^1(C_2)\bigr)\\
 &\oplus\bigl(H^0(C_1)\otimes H^2(C_2)\bigr)
 \end{split}
\]
vanish.  Choose $p_i\in C_i$ and the algebraic projectors
\[
 \pi_1^{C_i}=\Delta_{C_i}-[p_i\times C_i]-[C_i\times p_i].
\]
With $e$ from Equation~\eqref{geom:projector}, the correspondence
\[
 \bigl(\pi_1^{C_1}\boxtimes\pi_1^{C_2}\bigr)
       \circ{}^t\Gamma_g\circ e
 \in\CH^2(X\times C_1\times C_2)_\Q
\]
induces the asserted injection on the whole rational space $U_X$.
Proposition~\ref{geom:period} supplies the additional Hodge
hypothesis at a very general parameter.
\end{proof}

\subsection{Successive generic choices}

We call a set residual if its complement is a countable union of
closed nowhere-dense subsets.  In the connected coordinate domains
below, every residual set is dense by Baire's theorem.

\begin{lemma}[Ordered construction choices]\label{geom:genericity}
Every nonempty regular period chart admits a smaller normalized product
chart $A\times Q$ of abelian data and affine section coordinates as in
Section~\ref{geom:section-charts},
mapping into it and lying in the simultaneous geometric open of
Propositions~\ref{geom:period} and \ref{geom:domination}.
For any prescribed countable collection of proper closed analytic
period loci in the original chart, the following choices are available
in this order:
\begin{enumerate}
\item abelian data $a$ in a residual subset of $A$;
\item for each such $a$, an actual section pencil in a residual subset
of each prescribed nonempty admissible pencil open meeting $Q$;
\item after that pencil is fixed, all but countably many points of its
intersection with $Q$ and with a nonempty algebraic good open.
\end{enumerate}
The selected points avoid the prescribed period loci.  The geometric
requirements hold on the algebraic good open, and the exceptional points
for proper marked rational Hodge-tensor conditions on its whole connected
smooth part form a countable set.  The pencils in
\textup{(ii)} have the form $d_t=d_0+td_1$ with $\ell(d_0)=1$ and
$\ell(d_1)=0$, and their leading section $d_1$ is smooth and avoids
two-torsion.  The final countable subset may include any further
countable exceptions arising after the pencil is fixed; no uniform
exceptional list across pencils is required.
\end{lemma}

\begin{proof}

Start with any nonempty regular period chart.  The simultaneous
geometric open of Propositions~\ref{geom:period} and
\ref{geom:domination} is dense, so choose a datum $(a_*,[d_*])$
in its intersection with this chart.  In the same section fiber,
choose an independent section $d_\infty$ whose divisor is smooth
and avoids two-torsion.  Trivialize the section bundle, choose
$\ell(d_*)=1$ and $\ell(d_\infty)=0$, extend these choices in a
local frame, and shrink the abelian base.  Smoothness and avoidance
of two-torsion are open conditions, so $\PP(\ker\ell)$ contains
a good section over every point of this base.

Use the normalization $\ell(d)=1$ of
Section~\ref{geom:section-charts} and choose a product polydisc
$A\times Q$ within the original chart and the simultaneous
geometric open.  Its period map is a local isomorphism.  These
choices precede any actual pencil.  Affine lines in each section
chart are the pencils $d_t=d_0+t d_1$ on the original affine
$t$-line, with $\ell(d_0)=1$ and $\ell(d_1)=0$.
The good directions form a nonempty dense Zariski-open subset of
$\PP(\ker\ell)$.  Since the direction map on lines meeting $Q$
is a surjective submersion, pencils with a smooth leading section
avoiding two-torsion form a dense open in every nonempty admissible
pencil open.  Every pencil meeting $Q$ lies in the geometric open
at all but finitely many parameters.  Its period map is nonconstant,
since the period differential is injective on section directions.

Pull back the prescribed analytic conditions to
$H_m\subset A\times Q$.  For each $m$, the set
\[
 A_m=\{a\in A:\{a\}\times Q\subset H_m\}
\]
is closed: it is the intersection, over $q\in Q$, of the closed
sets $\{a:(a,q)\in H_m\}$.  It has empty interior, since otherwise
$H_m$ would contain an open product and hence the connected chart
by the identity theorem.  The complement of $\bigcup_m A_m$ is
therefore residual in $A$.  For every $a$ in this complement, each
$H_{m,a}$ is a proper closed analytic subset of $Q$.

Fix such $a$.  In the open space of affine lines meeting $Q$, put
\[
 \mathcal L_{m,a}=\{\Lambda:\Lambda\cap Q\subset H_{m,a}\},
\]
where $\Lambda$ denotes an affine line.
This set is closed: every point of a limiting line inside $Q$ is a
limit of points on the approximating lines inside $Q$.  It has empty
interior, since an open set of lines meeting $Q$ sweeps out an open
subset of $Q$; equivalently, evaluation from the incidence space of
lines and points is a submersion.  Excluding these countably many
closed nowhere-dense sets leaves a residual set in every nonempty
open of admissible pencils.  For any remaining line, $\Lambda\cap Q$
is connected, and each condition restricts to a proper analytic
subset of this one-dimensional domain, hence to a discrete and
countable set.  Their union is countable.

For the exceptional Hodge-tensor conditions on the connected smooth
open of a whole pencil, enlarge the initial list to include every
proper Hodge locus of a rational tensor in the marked local system.
Continue markings on the universal cover.  A tensor condition
holding identically on any open would, by analytic continuation,
hold on the selected pencil disc for one of the marked rational
tensors already excluded there.  Thus every such condition has
discrete zeros on marked pencil charts.  There are only countably
many rational tensors, continuation choices, and charts, so only
countably many exceptional pencil parameters.

Intersect the residual pencil choices with the dense open of good
directions.  The geometric bad parameters are finite, while the local
analytic exceptions and the continued Hodge-tensor exceptions are
countable.  Any further pencil-specific countable set can be added
after fixing that pencil.  This proves the stated order and conclusions.
\end{proof}
\section{Algebraic K3 pencils with controlled compactifications}
\label{pencil:section}

We keep the notation of Proposition~\ref{geom:structure}.  The integral lattice
\(V\) has signature \((2,14)\), whereas its rank-eight rational summand \(U\)
has signature \((2,6)\).  It is useful to embed the larger lattice: the local
vanishing cycles below are roots in \(V\), although their projections to
\(U\) need not be roots in an integral K3 lattice.  We first realize the
periods by polarized K3 surfaces with rational double points.  Their
canonical contractions retain enough descent data to work over the
original parameter curve.  A smooth K3 family is obtained on a smaller
open of that curve, and its quadratic form on \(U\) is the form inherited
from \(Y\).

\begin{proposition}[K3 realization of a section pencil]
\label{pencil:realization}
Choose the abelian data and the pencil \(d_t=d_0+t d_1\) in the generic
sets of Lemma~\ref{geom:genericity}.  There is a nonempty algebraic open
\(O\subset\AA^1_t\) and a smooth projective
family of K3 surfaces \(S^0\to O\) with a flat isometric embedding of
rational Hodge variations
\[
 U\ \hookrightarrow\ R^2(S^0/O)_*\Q.
\]
The form on \(U\) here is its cup-product form inherited from \(Y\).
For very general \(t\), this summand is the transcendental part of
\(H^2(S^0_t,\Q)\).  For every \(s\) in a nonempty Zariski-open subset
of \(\AA^1_s\), every smooth projective model of the total family
obtained by \(t=s+u^2\) is uniruled.  The K3
family over \(O\) is defined over the original \(t\)-line; an additional
unspecified finite base change is not part of the conclusion.
\end{proposition}

\subsection{The lattice extension}

In this section \(M\) denotes the lattice enlargement, and
\(\ell(A)\) denotes the minimum number of generators of a finite abelian
group \(A\).

\begin{lemma}[Enlargement of the lattice]
\label{pencil:lattice}
Let \(V\) be an even lattice of signature \((2,14)\) with elementary
\(2\)-primary discriminant group.  Suppose a group \(G\subset O(V)\) acts
trivially on that group.  There are an even finite-index overlattice \(M\)
of \(V\), a primitive embedding \(M\hookrightarrow\Lambda_{\mathrm{K3}}\), and a
primitive vector \(h\in M^\perp\) with \((h,h)>0\), such that the action of
\(G\) extends to \(\Lambda_{\mathrm{K3}}\) and is the identity on \(M^\perp\).
\end{lemma}

\begin{proof}
Put \(A_V=V^*/V\), and denote its discriminant quadratic and bilinear forms by
\(q:A_V\to\Q/2\Z\) and \(b:A_V\times A_V\to\Q/\Z\).
On \(A_V\simeq\mathbf F_2^n\), the form \(\beta=2b\), valued in
\(\Z/2\Z\), is nondegenerate and symmetric.  Let
\(E=\ker(x\mapsto\beta(x,x))\), so \(\dim E\geq n-1\).
The restriction of \(\beta\) to \(E\) is alternating.  The inverse
image in \(E\) of a Lagrangian in its symplectic quotient by the radical
is a totally isotropic subspace of dimension
\[
 \frac{\dim E+\dim\operatorname{rad}(\beta|_E)}2
 \ \geq\ \left\lfloor\frac{n-1}{2}\right\rfloor.
\]
On this subspace, \(b(x,x)=0\) implies that \(q(x)\) is integral
modulo \(2\Z\); moreover
\(q(x+y)-q(x)-q(y)=2b(x,y)=0\) modulo \(2\Z\).
Thus \(q\) restricts to a linear functional with values in
\(\mathbf F_2\).  If \(n\geq5\), the subspace has dimension at least
two, and this functional has a nonzero isotropic vector.

Adjoin the corresponding isotropic element to the lattice.  The new
lattice is even and has discriminant form \(H^\perp/H\), where \(H\)
is the order-two subgroup generated by that vector.  This reduces the
discriminant length by two and leaves the rational quadratic space
unchanged.  Repetition gives an even overlattice
\(M\) with \(\ell(A_M)\leq4\).  Each enlargement is \(G\)-stable, since
\(G\) acts trivially on the preceding discriminant group; its action on
the new discriminant group is again trivial.

The sufficient primitive-embedding criterion of Nikulin
\cite[Chapter~14, Theorem~1.12]{Huybrechts} applies, since
\[
  3>2,\qquad 19>14,\qquad
  22-16=6\geq\ell(A_M)+2.
\]
Choose such an embedding.  Its complement has signature \((1,5)\), and
therefore contains a primitive positive integral vector \(h\).  The
complement is primitive in the unimodular K3 lattice, so \(h\) is also
primitive in \(\Lambda_{\mathrm{K3}}\).
Unimodular gluing identifies the discriminant forms of \(M\) and
\(M^\perp\) with opposite signs.  Consequently an isometry of \(M\)
trivial on \(A_M\), together with the identity on \(M^\perp\), extends
uniquely to \(\Lambda_{\mathrm{K3}}\).  These extensions respect composition.

In particular, a root \(r\in V\), with \(r^2=-2\), remains an integral
root throughout the construction.  Its reflection is
\[
 s_r(x)=x+(x,r)r.
\]
It is integral on every even overlattice under consideration and on
\(\Lambda_{\mathrm{K3}}\), and is the identity on \(M^\perp\).
Whenever a monodromy operator on \(V\) is a product of these
reflections, its extension is the same product of K3 root reflections.
\end{proof}

\subsection{Local monodromy of the section pencil}

We analyze the normalized original-line pencil \(d_t=d_0+t d_1\)
already chosen in Proposition~\ref{pencil:realization} through
Lemma~\ref{geom:genericity}.  Its abelian data \((B,H,\chi)\) are fixed.
The leading member \(d_1\) is
smooth and avoids two-torsion.  The corresponding variations \(V\)
and \(U\) are those of Proposition~\ref{geom:structure}.
The next lemma distinguishes the finite root monodromies from the
sign at infinity; the substitution \(t=s+u^2\) will kill the latter.

\begin{lemma}[Local monodromy of the section pencil]
\label{pencil:monodromy}
At every finite singular parameter of this general pencil, the
monodromy on \(V\) is a product of reflections in mutually orthogonal
integral \((-2)\)-roots in \(V\).  These roots are perpendicular to the
limiting Hodge line.  The periods extend to interior periods after
killing these finite monodromies.  At infinity the monodromy on \(V\)
is \(-\id\).
\end{lemma}

\begin{proof}
\emph{The possible singular members.}
The morphism defined by \(|2H|\) factors through inversion and realizes
the quotient as the classical quartic Kummer surface.  A general line in
the dual projective space meets its discriminant either at a simple
Lefschetz contact in the smooth locus or at a hyperplane through one
ordinary Kummer node.  Here is the transversality assertion being used.
The hyperplanes containing each Kummer node form a plane in the dual
space.  A general point of such a plane avoids every other node and is
transverse to the quadratic tangent cone at that node.  For tangencies
in the smooth locus, the conormal incidence has generically separable
projection to the dual variety.  If its image has dimension less than
two, a general line avoids it.  If it has dimension two, its general
smooth point represents a unique contact, because that contact point
is recovered as the tangent hyperplane to the dual variety.
Separability makes the quadratic contact nondegenerate.  The smooth-contact dual component cannot be a node plane, by biduality.
The higher
contacts, intersections of these discriminant components, and their
singular loci have codimension at least two in the dual projective
space.  A general line avoids them and meets each remaining component
transversely.  Upstairs on \(B\), the
first event gives two ordinary nodes exchanged by inversion.  At a
two-torsion point in the second event, the invariant local equation has
a nondegenerate quadratic term and its constant term vanishes to first
order in \(t\).  Thus the branch has an ordinary node with transverse
smoothing.  The base points of a general pencil are smooth on every
member, since the differentials of \(d_0\) and \(d_1\) there are
independent.  Pulling back along the isogeny to \(B'\) gives disjoint
copies of these events.  We examine their inversion quotients.

\emph{An exchanged pair.}
At an exchanged pair away from the fixed points, the branch curve on
\(K=\Km(B')\) acquires one ordinary node.  The local double surface is
\(z^2=xy+t\).  Its vanishing sphere has square \(-2\), and the deck
involution \(\delta\) negates it.  It therefore gives an integral root
of \(V\) and its Picard--Lefschetz reflection.

\emph{A positive-sign fixed point.}
At a fixed point with nonbranch exceptional curve on \(K\), choose
coordinates in which the upstairs model is
\[
 z^2=x^2+y^2+t,\qquad
 j(x,y,z)=(-x,-y,z).
\]
These coordinates can be chosen equivariantly: the branch equation
is even in \((x,y)\), and the holomorphic Morse lemma, applied to
its nondegenerate even quadratic-coefficient matrix, preserves that
parity.  The parameter is replaced locally by its constant term,
which is a coordinate because the discriminant is met transversely.
On the chart \(a=x^2,\ v=y/x\) of the resolved base quotient, the
double surface has equation
\[
 z^2=a(1+v^2)+t.
\]
At \(t=0\), its branch is the exceptional curve \(a=0\) together with
a residual curve meeting it at \(v=i,-i\).  At each crossing the equation
is \(z^2=ab+t\).  There are consequently two disjoint ordinary node
smoothings, and two orthogonal integral anti-invariant roots in \(V\).
The local monodromy is their product.  The difference of the two lifts
of the nonbranch exceptional curve, of square \(-4\), need not be fixed
by this monodromy; it is not being used as one of the individual roots.

\emph{A negative-sign fixed point.}
At a fixed point with branch exceptional curve, the involution in the
upstairs quadratic model is instead
\[
 j(x,y,z)=(-x,-y,-z).
\]
It acts antipodally on the real vanishing sphere \(S^2\), and hence by
\(-1\) on its degree-two class.  There is no vanishing class in the
\(j\)-invariant quotient cohomology.  To check the resolution contribution,
for \(t\ne0\) the points \((0,0,\pm\sqrt t)\) are exchanged by \(j\).
Their quotient is one smooth point.  On the chart represented near one
chosen square root, the induced sheet involution is
\((x,y)\mapsto(-x,-y)\).  Blowing up this point therefore gives a
\(\PP^1\) fixed pointwise by \(\delta\); its cohomology class is
\(\delta\)-invariant.  More explicitly, on the Kummer-resolution chart
\(a=x^2,\ v=y/x\), put \(\eta=xz\).  The finite double-cover model is
\[
 \eta^2=a\bigl(a(1+v^2)+t\bigr),\qquad
 \delta(\eta)=-\eta.
\]
For \(t\ne0\), the preimage of \(a=0\) is that ramification curve;
its square is \(-1\), because the base exceptional curve has square
\(-2\).  Its class is the only added resolution class at this point.
Thus this event contributes no negative eigenspace to \(V\).

\emph{The full monodromy operator.}
All the displayed roots are integral classes on the nearby smooth
\(Y\).  They are mutually orthogonal because their Milnor
neighborhoods are disjoint, and the local monodromy negates each of
them.  To see that there is no further monodromy, use the orthogonal
decomposition of variations in Proposition~\ref{geom:structure},
\[
 V_\Q=U\oplus\langle a_1,\ldots,a_8\rangle_\Q,
 \qquad U\simeq U_X.
\]
On \(X\), the complement of the disjoint node neighborhoods is
locally topologically trivial, so Picard--Lefschetz gives the product
of the corresponding commuting reflections.  An exchanged pair has
a one-dimensional \(j\)-invariant vanishing space.  At a positive
fixed point, \(j\) has degree \(+1\) on the vanishing sphere and
contributes one dimension to \(U_X\).  At a negative fixed point it
has degree \(-1\) and contributes none.  On the difference lines
\(\Q a_i\), the two labels are exchanged by \(\sqrt t\mapsto-\sqrt t\)
exactly at the affected positive fixed points; all other difference
lines are fixed.  Consequently the full operator on \(V_\Q\) is an
involution whose negative eigenspace has one dimension per exchanged
pair and two per positive fixed point.  This is precisely the number
of mutually orthogonal roots found above.  Those roots span that
negative eigenspace, proving that the entire operator is their
product of reflections.  The same argument includes all disjoint
lifts of a single event through the isogeny.

\emph{The limiting periods.}
After a quadratic cover at each ordinary node, the upstairs surface
family has a simultaneous local analytic resolution.  In the
\(A_1\) case used here, write \(xy=(z-v)(z+v)\) and blow up the
non-Cartier ideal \((x,z-v)\).  On the chart \(z-v=xa\), the equations
are \(z=v+xa\) and \(y=a(2v+xa)\), so \((x,a,v)\) are smooth
coordinates.  On the other chart \(x=(z-v)b\), they are
\(z=by-v\) and \(x=b(by-2v)\), with smooth coordinates \((b,y,v)\).
Both projections to \(v\) are smooth and their central fibers resolve
the node.  The relative two-form \(dx\wedge dz/x\) becomes
\(dx\wedge da\) on the first chart, while the equal expression
\(-dy\wedge dz/y\) becomes \(db\wedge dy\) on the second.  Thus the
resolution is crepant.  These are the local \(A_1\) charts of the
classical simultaneous-resolution constructions
\cite{Brieskorn,ArtinResolution}.  They glue to the unchanged locus
because the exceptional neighborhoods are disjoint.

Choose a nonzero
\(\theta_+\in H^0(B',L)^{j=+1}\).  Using an invariant volume form
on \(B'\), the Hodge line of \(U_X\) has frame
\[
 \frac{\theta_+\,dx\wedge dy}{z}.
\]
Crepantness extends it regularly after the local base change and
resolution.  It remains a nonzero global form on the resolved central
surface: test at a point of \(B'\) outside the branch divisor and the
zero divisor of \(\theta_+\).  At a minus-sign node its coefficient
is odd and vanishes at that node, but this does not make the global
Hodge-line frame vanish over the parameter.

The chosen simultaneous resolution need not carry the involutions.
Their flat actions on cohomology nevertheless extend after the base
change, preserve the limiting Hodge filtration by continuity, and
retain the prescribed eigenspaces.  The nonzero limiting holomorphic
form on the projective resolved central surface has positive norm.
Thus \(U_X\), and hence \(U\), has an interior limiting period.
The section just used was single-valued before the local base change,
so inertia fixes its limiting vector.  In particular the negative
eigenspace of local monodromy is perpendicular to that vector.  The
additional exceptional-difference classes remain algebraic, giving
the same conclusion for the roots in \(V\).


\emph{Infinity.}
Near infinity write \(q_\infty=t^{-1}\).  The equation is
\(z^2=q_\infty^{-1}(d_1+q_\infty d_0)\).
The leading member \(d_1\) is smooth and avoids the fixed points.
After adjoining \(\sqrt{q_\infty}\) and rescaling \(z\), the family
extends smoothly.  The only twist around \(q_\infty=0\) is the change
of sign of this square root.
This is \(\delta\), which acts as \(-\id\) on \(V\).
\end{proof}


\subsection{The period stack and polarized contraction}

The monodromies just calculated describe integral period data; they
do not yet give an algebraic K3 family on the original curve.  To
construct that family, we pass through polarized contractions, for
which reflections in contracted roots act trivially.  A
\emph{polarized ADE K3 surface} is a normal surface with rational double
points whose minimal resolution is a K3 surface, together with a
primitive ample polarization.  Here a polarization in a family is a
relative Picard class, locally represented by a line bundle.  An ADE K3
surface is the canonical polarized contraction of its resolution with
the pulled-back big and nef class; this contraction is given by a
sufficiently high power of the class.  The relative-Picard convention
removes scalar automorphisms of a representing line bundle from the
moduli problem.

Fix a primitive positive vector \(h\in\Lambda_{\mathrm{K3}}\).  Choose a connected
component \(\mathscr D_h\) of
\[
 \{[\sigma]\in\PP(h^\perp\otimes\C):
       (\sigma,\sigma)=0,\ (\sigma,\overline\sigma)>0\},
\]
and let \(\Gamma_h\) be the isometries of \(\Lambda_{\mathrm{K3}}\) fixing \(h\)
and preserving this component.  Unimodular gluing identifies \(\Gamma_h\)
with the stable orthogonal group of \(h^\perp\) preserving
\(\mathscr D_h\): restriction acts trivially on the discriminant group,
and every such restricted isometry extends by the identity on \(\Z h\).
Thus \([\Gamma_h\backslash\mathscr D_h]\) is the integral period stack
for primitive degree \((h,h)\) polarizations.  This identification is
independent of the choice of the primitive vector of that square; the
usual transitivity follows from the Eichler criterion for the K3
lattice, which contains two hyperbolic planes; see
\cite[Chapter~14, Corollary~1.10]{Huybrechts}.

\begin{lemma}[The contraction morphism from the period stack]
\label{pencil:periodstack}
Let \(\mathscr P_h=[\Gamma_h\backslash\mathscr D_h]\), and let
\(\mathscr M_h\) denote the moduli stack of primitively polarized ADE K3
surfaces of degree \((h,h)\).  There is a morphism
\[
 q:\mathscr P_h\longrightarrow\mathscr M_h
\]
whose composition with the period map of quasi-polarized smooth K3
surfaces is their canonical polarized contraction.  The pullback of
the Hodge line of \(\mathscr M_h\) is the tautological period line,
including its \(\Gamma_h\)-linearization.

For \([\sigma]\in\mathscr D_h\), let \(W_\sigma\) be generated by the
reflections in the roots
\[
 R_\sigma=\{r\in\Lambda_{\mathrm{K3}}:r^2=-2,\ (r,h)=0,\ (r,\sigma)=0\}.
\]
The kernel of the homomorphism on stabilizers induced by \(q\) is exactly
\(W_\sigma\).
\end{lemma}

\begin{proof}
The existence and factorization of \(q\) are
\cite[Lemma~1.12]{DiLorenzoFringuelliVistoli}.  We use the analytic
form of that morphism.  Its hypotheses are primitive ample polarization
classes on ADE K3 surfaces, primitive big and nef classes on smooth K3
surfaces, and the full integral period stack over \(\C\).  The stable
group and the representable, surjective \'etale period map of the
quasi-polarized stack are described in
\cite[Sections~1.6 and~1.8; proof of Lemma~2.6 for surjectivity]{DiLorenzoFringuelliVistoli}.  Taking one
component of the period domain with its stabilizer gives the equivalent
stack used here.  We now prove the two stated compatibilities.

Let \(\mathscr K_h\) be the smooth quasi-polarized K3 stack.  Over a
local \'etale period chart its universal smooth family contracts to the
pullback of the ADE family.  The latter is flat with Gorenstein fibers,
so its relative dualizing sheaf is invertible.  Off the exceptional
curves, contraction identifies the relative two-forms.  In the total
smooth family over this full period chart the exceptional locus has
codimension at least two: the curves occur only over root walls.
The identification of the two line bundles therefore extends over
the exceptional locus.  On each fiber it is the usual crepant
identification for resolution of rational double points.  This gives
the claimed comparison of Hodge lines after pullback to
\(\mathscr K_h\).

For descent, consider
\(\mathscr K_h\times_{\mathscr P_h}\mathscr K_h\).  It is locally
smooth over the period domain, and the ample locus is dense in every
component.  There each period arrow comes from an actual polarized
isomorphism by Torelli, so the two induced Hodge-line identifications
agree.  Equality extends over this reduced relation space.  Hence the
comparison descends with its entire \(\Gamma_h\)-linearization,
including on the root mirrors.

Fix a root \(r\in h^\perp\) and consider its mirror.  Its reflection
\(s_r\) fixes the period vector, so its image under \(q\) acts
trivially on the Hodge line.  At a very general point of the mirror,
the resolution has rational Picard space \(\Q h\oplus\Q r\).
The polarized automorphism induced by \(q(s_r)\) lifts uniquely to
that resolution.  It fixes \(h\), and preservation of effectivity
forces it to fix the effective choice of \(r\).  It is symplectic.
The transcendental Hodge structure is simple, being polarizable of
K3 type with no nonzero rational \((1,1)\)-classes.  A Hodge endomorphism
fixing its \((2,0)\)-line pointwise is the identity: the kernel of the
difference from the identity is a nonzero Hodge substructure.
The lifted automorphism thus acts identically on all second
cohomology and is the identity by Torelli.

On the whole mirror, \(q(s_r)\) and the identity are sections of a
separated isomorphism space.  Their equality on the dense locus just
considered extends everywhere.  Since root reflections generate
\(W_\sigma\), this group is killed even at points with several or
persistent contracted roots.  This agrees with the generic
root-reflection kernel identified in
\cite[Lemma~3.6]{DiLorenzoFringuelliVistoli}.

Conversely, an integral Hodge isometry fixing \(h\) takes the ample
chamber of a chosen resolution to another chamber whose closure
contains \(h\).  The Weyl group generated by the algebraic roots
perpendicular to \(h\) acts transitively on these chambers: only walls
through \(h\) separate chambers incident to \(h\).  After composition
with an element of \(W_\sigma\), strong Torelli realizes the isometry
by an automorphism of the resolution.  By the factorization defining
\(q\), its image is the induced automorphism of the contraction.
If the image is the identity, its lift is the identity on the dense
complement of the exceptional curves and hence everywhere.  The
original isometry then belongs to \(W_\sigma\).
\end{proof}

\subsection{Descent and the relative canonical class}

The use of the integral period stack in Lemma~\ref{pencil:periodstack}
is essential: the local system and its monodromy supply a map to that
stack, whereas a map to the coarse period space need not supply a
family.  The following curve statement also keeps track of ramified
descent and of the relative canonical bundle.

\begin{lemma}[Realization over the original curve]
\label{pencil:descent}
Let \(C\) be a smooth projective complex curve, let \(D\subset C\) be
finite, and put \(C^\circ=C\setminus D\).  Suppose the following data are
given:
\begin{enumerate}[label=\textup{(\roman*)}]
\item a local system \(\mathcal L\) of lattices isometric to
\(\Lambda_{\mathrm{K3}}\), with a primitive flat section \(h\) of positive square;
\item a holomorphic line \(\mathcal F\subset\mathcal L\otimes\mathcal O_{C^\circ}\)
with
\[
 (\mathcal F,\mathcal F)=0,\qquad
 (\sigma,\overline\sigma)>0\quad(0\ne\sigma\in\mathcal F),\qquad
 (\mathcal F,h)=0;
\]
\item finite local monodromy at every point of \(D\), and an interior
limit of the period line after a local cover killing that monodromy.
\end{enumerate}
Then there is a normal projective model \(f:\mathcal Z\to C\) whose
restriction to \(C^\circ\) is a family of polarized ADE K3 surfaces
realizing these integral period data modulo contracted roots.

Suppose, in addition, that the local monodromy at each \(c\in D\), viewed
on the limiting marked lattice, belongs to the Weyl group generated by
\[
 R_c=\{r\in\mathcal L_c:r^2=-2,\ (r,h)=0,\ (r,\mathcal F_c)=0\}.
 \tag{\(*\)}\label{pencil:rootcondition}
\]
Then \(\mathcal Z\to C\) can be chosen to be a flat family of polarized
ADE K3 surfaces on all of \(C\).  Its extended Hodge line descends to a
line bundle \(\lambda\) on \(C\), and
\begin{equation}
 \omega_{\mathcal Z/C}\simeq f^*\lambda,
 \qquad K_{\mathcal Z}=f^*(K_C\otimes\lambda).
 \label{pencil:canonical}
\end{equation}

Under hypotheses \textup{(i)}--\textup{(iii)}, independently of
\eqref{pencil:rootcondition}, let
\(\mathcal U\subset\mathcal L_\Q\) be a nondegenerate rational
sub-local system containing \(\mathcal F\), and suppose every flat root
which is of type \((1,1)\) on an open set is perpendicular to \(\mathcal U\).
After shrinking a nonempty algebraic open in \(C^\circ\), there is a
smooth projective K3 family \(\mathcal S^0\) there and an identification
of \(\mathcal U\) with an isometric Hodge summand of
\(R^2\mathcal S^0_*\Q\).  Here a flat root means a locally flat
section; it need not be a globally invariant vector.  No change of the
base curve occurs in this conclusion.
\end{lemma}

\begin{proof}
\emph{The family and its algebraization.}
The torsor of markings of \((\mathcal L,h)\), with a component chosen by
\(\mathcal F\), and the equivariant period map define a holomorphic map
\(C^\circ\to\mathscr P_h\).  Its monodromy lies in \(\Gamma_h\): it fixes
\(h\), has trivial restricted discriminant action by unimodular gluing,
and preserves the period component by continuity.  Pullback along
\(q\) in Lemma~\ref{pencil:periodstack} gives an actual analytic family
of polarized ADE K3 surfaces on \(C^\circ\), with its compatible descent
data.

Choose a finite Galois cover \(\widetilde C\to C\) whose ramification
indices at the points above \(D\) are divisible by the local monodromy
orders.  Such a cover exists: choose a common multiple \(N\) of these
orders, use Riemann--Roch with sufficiently large poles away from \(D\)
to prescribe a simple zero at each point of \(D\), and normalize the
cyclic cover obtained by adjoining an \(N\)-th root of that function.
Additional branch points lie over points where the original family is
already defined.  If \(D\) is empty the identity cover suffices.
The pulled-back local system and period extend across the punctures:
in a flat trivialization the latter is the holomorphic interior
extension in hypothesis \textup{(iii)}.  They define a
map \(\widetilde C\to\mathscr P_h\), and hence a proper analytic family
\(\widetilde{\mathcal Z}\to\widetilde C\) of polarized ADE K3 surfaces.
The original monodromy data equip it, functorially through \(q\), with
an action of the Galois group.  All identifications and cocycles are
provided by this stack morphism; no choices of smooth simultaneous
resolutions have to glue.

This compact analytic family is algebraic and normal.  For normality,
use flatness and the normal ADE fibers over the smooth base.  To check
the projectivity
needed for GAGA, represent its relative polarizing class by line bundles
on small open sets of the base.  A common sufficiently divisible power
is relatively ample.  On an overlap, the ratio of two such representatives
is trivial on each fiber.  Its direct image is a line bundle, and
evaluation is an isomorphism, by cohomology and base change and
\(H^1(\mathcal O)=0\) on the fibers.  Thus the ratio comes from the
base.  After local trivializations, the remaining gluing obstruction
is a scalar cocycle in \(H^2(\widetilde C,\mathcal O_{\widetilde C}^*)\).
The exponential sequence makes this group zero, since
\[
 H^2(\widetilde C,\mathcal O_{\widetilde C})=0,
 \qquad H^3(\widetilde C,\Z)=0.
\]
The polarizing power consequently glues to a relatively ample line bundle.
For a sufficiently high power its direct image is a vector bundle, and
the family embeds into the associated projective bundle on
\(\widetilde C\).  That vector bundle and the closed embedded family
are algebraic by GAGA \cite[Theorems~2--3]{SerreGAGA}.  More explicitly,
first algebraize the vector bundle on the projective curve, then
algebraize the coherent ideal of the family in its projective bundle.
The finite group action algebraizes by the same fully faithful
comparison.  Twist the relative ample bundle by a sufficiently positive
base bundle to make it ample on the total space.  The tensor product
of its translates by the Galois group is ample and has its natural
linearization; invariant sections of a sufficiently high power give a
projective finite quotient.  That quotient is a normal projective model
\(\mathcal Z\to C\).  Over \(C^\circ\) its construction is descent of
the pullback of the initial family, also at any auxiliary branch points.
This proves
the first assertion, with no cover remaining in its conclusion.

\emph{Descent at a Weyl puncture.}
Assume \eqref{pencil:rootcondition}.  On a disc killing the monodromy,
write the map to the original disc as \(t=w^m\).  The cyclic inertia
group acts on the central polarized surface via the homomorphism on
stabilizers induced by \(q\).  Its image is the identity by
Lemma~\ref{pencil:periodstack}.  Thus the group acts pointwise on the
whole central surface as the identity, including its singular points.

Here is the corresponding scheme-theoretic descent calculation.  Let
\(A\) be a local completed ring of \(\widetilde{\mathcal Z}\) along that
fiber.  Flatness makes \(w\) a nonzerodivisor.  If \(g\) generates inertia,
then \(g(w)=\zeta w\), with \(\zeta\) primitive of order \(m\), and
\(g\) is the identity on \(A/(w)\).  Set \(B=A^{\langle g\rangle}\).
Reynolds averaging makes \(B\to A/(w)\) surjective.  If \(a=wb\) is
invariant, then \(g(b)=\zeta^{-1}b\).  For \(m>1\), reduction modulo
\(w\) forces \(b\) to be divisible by \(w\).  At each subsequent
division the character changes by \(\zeta^{-1}\); until the \(m\)-th
division it is nontrivial and the reduction must again vanish.  This
gives, also trivially when \(m=1\),
\[
 B\cap wA=w^mB,\qquad B/(t)\simeq A/(w).
\]
There is also an exact base-change statement.  Write
\(A_j=\{a:g(a)=\zeta^j a\}\), for \(0\leq j<m\).
Successive division by \(w\) gives \(A_j=w^jB\); character projection
and the fact that \(w\) is a nonzerodivisor then give
\[
 A=\bigoplus_{j=0}^{m-1}w^jB,
 \qquad
 B\otimes_{\C[[t]]}\C[[w]]\xrightarrow{\ \sim\ }A,
 \quad t\longmapsto w^m.
\]
These assertions pass from completed rings to the local rings by
faithful flatness of completion.  Thus the quotient central fiber is
exactly the original reduced ADE surface, and the upstairs family is
its flat pullback.  The quotient is flat over the parameter disc,
since \(t\) is a nonzerodivisor in \(B\), a torsion-free module over
the discrete valuation ring of the disc.  This proves the assertion
about all fibers and rules out a multiple central fiber.

The underlying ADE family has now descended.  To retain its primitive
polarization, choose on a smaller invariant
analytic \(w\)-disc a line bundle \(L\) representing its relative Picard
class; the preceding scalar-cocycle argument applies on the disc as well.
The same base-change argument used above shows that
\(g^*L\otimes L^{-1}\) comes from the disc and that automorphisms of
\(L\) are base units; the base line is trivial after shrinking.  For a
\(g\)-semilinear lift \(T\), write \(T^m=c\,\id\), where commutation with
\(T\) makes \(c=c(t)\) an invariant base unit.  Replacing \(T\) by
\(aT\) replaces \(c\) by \(c\prod_{i=0}^{m-1}g^i(a)\).  Every unit in
\(\C\{t\}\) has an \(m\)-th root after shrinking, whose norm as an
invariant upstairs unit is its \(m\)-th power.  Taking such a root of
\(c^{-1}\) therefore arranges \(T^m=\id\).  On the proper connected
reduced central fiber, \(T\) is one scalar \(\chi\in\mu_m\).  Replacing
\(T\) by \(\chi^{-1}T\) makes it the identity on \(L/(w)L\), without
changing the relative Picard class or the separately prescribed
Hodge-line linearization.

At any central point, including an ADE point, work with the actual
analytic local rings and let \(s\) be a local generator of \(L\).  Then
\(s_*:=m^{-1}\sum_{i=0}^{m-1}T^i(s)\) is invariant and congruent to
\(s\) modulo \(w\), hence is still a generator.  The invariant local
modules are consequently free of rank one over the invariant local
rings and pull back to \(L\).  Their invariant transition units glue
them to an actual invertible sheaf on the analytic quotient.  Under
\(B/(t)\simeq A/(w)\) its central restriction is exactly \(L_0\), so
the primitive ample class, not a power, descends.  These local classes
agree with the original one off the punctures and glue on \(C\); the
preceding scalar-cocycle argument glues their primitive representatives,
which GAGA algebraizes on the projective quotient.

The Hodge line must descend with its prescribed character, separately
from the scalar adjustment of the polarizing bundle.  Its limiting
inertia action is trivial: root reflections
fix a vector perpendicular to their roots.  Averaging a local frame
which is nonzero at the center gives an invariant frame with that
same central value, and hence a nowhere-vanishing frame on a smaller
disc.  Consequently
the extended Hodge line descends to a line bundle \(\lambda\) on \(C\).
By Lemma~\ref{pencil:periodstack}, it is the Hodge line of the descended
family, with precisely the linearization prescribed by \(\mathcal L\).
No character twist has been discarded.

Every ADE K3 fiber is Gorenstein with trivial dualizing sheaf.  A flat
family with these Gorenstein fibers over a smooth curve is a Gorenstein
morphism, so \(\omega_{\mathcal Z/C}\) is invertible.  Its direct image
is a line bundle and formation commutes with base change, since every
fiber has a one-dimensional space of canonical forms.  Evaluation is
an isomorphism on each fiber and therefore everywhere.  The identification
of this direct image with \(\lambda\) proves \eqref{pencil:canonical}.
For the subsequent intersection argument, its restriction to codimension
one would already suffice.

\emph{A smooth family on an open of the original curve.}
This part uses the normal contraction family from the first assertion,
not the additional Weyl condition.  We resolve its generic fiber over
\(k=\C(C)\), then recover \(\mathcal U\) by comparing local period
lifts through the cohomology of that same contraction.
The generic polarized ADE surface has a projective
minimal resolution over \(k\).  One way to see the field assertion is
to resolve its rational double points over a finite Galois extension.
Every Galois translate is a minimal resolution of the same surface;
the unique isomorphism over that surface supplies a descent datum and
automatically satisfies its cocycle.  A norm of a relatively ample
bundle makes this descent projective.  The descended resolution is
defined over \(k\) itself.  Spreading the morphism and its projective
embedding, and then shrinking an algebraic open of \(C^\circ\), gives
a smooth projective K3 family
\(\mathcal S^0\to C_0\) resolving the fixed contraction family.  All
fibers and exceptional configurations can be assumed to be of the
generic type on this open.

Here are the details of the cohomological comparison on \(C_0\).
For the minimal resolution \(\pi:S\to Z\) of a complex ADE K3
surface, let \(E_Z\subset H^2(S,\Q)\) be spanned by its exceptional
curves.  The rational cohomology sequence for the resolution gives
\[
 \pi^*:H^2(Z,\Q)\hookrightarrow H^2(S,\Q),
 \qquad \im(\pi^*)=E_Z^\perp.
\]
Indeed, the exceptional fibers are trees of rational curves, with
\(H^1=0\), and the restriction map in degree two records intersection
with their components.  It is onto on their span because the ADE
intersection matrix is negative definite.  This proves both the
injectivity and the stated image.  In a resolving family these
pullbacks are maps of higher direct-image sheaves.  Shrinking \(C_0\)
once more makes the sheaves local systems, by constructibility for
the algebraic proper morphisms.  They preserve Hodge structures and
the cup-product forms on the indicated orthogonal complements.

Locally lift the marked period map through the \'etale map
\(\mathscr K_h\to\mathscr P_h\).  Its K3 family has an identification
of its contraction with the fixed family obtained from \(q\).
The exceptional root lattice is locally constant on \(C_0\); all its
roots are therefore persistent in the sense of the hypothesis, and
are perpendicular to \(\mathcal U\).  The marked copy of
\(\mathcal U\) in this local lift consequently lies in the image of
\(H^2\) of the contraction.  The displayed pullback description gives
a flat, isometric map from it to \(R^2\mathcal S^0_*\Q\).
For two such local lifts, compare their minimal resolutions over the
identity of their common contraction on each fiber.  The resulting
marking discrepancy has trivial image under \(q\), so belongs to the
Weyl group in Lemma~\ref{pencil:periodstack}.  This group fixes
\(\mathcal U\) pointwise.  Thus the maps just constructed agree on
overlaps, including their flat structures, and glue over \(C_0\).
Nondegeneracy of \(\mathcal U\) makes its image an orthogonal Hodge
summand.  This proves the final assertion on the original curve.

For clarity about the shrinking in this argument, the mirrors of the
\(h\)-perpendicular roots are locally finite in \(\mathscr D_h\).
On a compact set of periods, a root orthogonal to
\(\langle h,\Re\sigma,\Im\sigma\rangle\) has bounded norm in a
fixed positive definite metric, since its square is \(-2\).
Only finitely many lattice vectors can occur.  A finite collection
of such local period charts covers the compact finite cover used
above.  In each chart a nonpersistent mirror has only isolated
intersections with the period curve.  Local finiteness and compactness
therefore leave only finitely many such intersection parameters.
Removing their images, as well as the exceptional points
from spreading, gives the open used here.  The argument uses constant
root type there; it requires no uniqueness assertion for simultaneous
resolutions across jumps in the root configuration.
\end{proof}

\begin{corollary}[A canonical-degree criterion]
\label{pencil:uniruledcriterion}
Under the Weyl condition of Lemma~\ref{pencil:descent}, suppose
\(\deg(K_C\otimes\lambda)<0\).  Every smooth projective model of the
function field of \(\mathcal Z\) is uniruled.
\end{corollary}

\begin{proof}
Let \(p:\widetilde{\mathcal Z}\to\mathcal Z\) be a smooth projective
resolution, and let \(A\) be sufficiently ample on \(\mathcal Z\).
Since \(\mathcal Z\) is normal and Gorenstein, the difference
\(K_{\widetilde{\mathcal Z}}-p^*K_{\mathcal Z}\) is supported on the
exceptional divisors.  Every such divisor has zero intersection with
\((p^*A)^2\): its image has codimension at least two, and the projection
formula applies.  Geometrically, two general members of \(|A|\) have
a complete-intersection curve avoiding all these images.
Equation~\eqref{pencil:canonical} consequently gives
\[
 K_{\widetilde{\mathcal Z}}\cdot(p^*A)^2
   =\deg(K_C\otimes\lambda)\,(A|_{\mathcal Z_c})^2<0.
\]
Here \(\mathcal Z_c\) is any general fiber, and its displayed
self-intersection is positive by ampleness.  The divisor \(p^*A\)
is nef, so every effective divisor, and by continuity every
pseudo-effective divisor, has nonnegative intersection with
\((p^*A)^2\).  Thus \(K_{\widetilde{\mathcal Z}}\) is not
pseudo-effective.  The theorem of Boucksom--Demailly--P\u{a}un--Peternell
\cite[Corollary~0.3]{BDPP} applies to this smooth projective threefold
and gives uniruledness.  This is a birational
property, proving the assertion for every smooth projective model.
\end{proof}

\subsection{Realization and the Hodge line}

\begin{proof}[Proof of Proposition~\ref{pencil:realization}]
\emph{Realization over the \(t\)-line.}
By Proposition~\ref{geom:structure}, the monodromy of \(V\) is trivial
on its discriminant group.  Apply Lemma~\ref{pencil:lattice}, extending
that monodromy to the K3 lattice and choosing its fixed positive class
\(h\).  The lattice embedding preserves the form inherited from \(Y\).
Give this lattice the period line of \(V\), which is also that of
\(U\).  Over \(\Q\) its orthogonal decomposition is
\[
 \Lambda_{\mathrm{K3},\Q}=U\oplus\langle a_1,\ldots,a_8\rangle_\Q
                       \oplus M_\Q^\perp.
\]
The last two summands, of dimensions eight and six, are of type
\((1,1)\), so the complement of \(U\) has dimension fourteen.  For the
choices specified in
Lemma~\ref{geom:genericity}, and a very general parameter, \(U\) has
no rational \((1,1)\)-classes by Proposition~\ref{geom:period}.
If a locally flat rational root remains algebraic on a nonempty
open, its projection to \(U\) in the displayed orthogonal decomposition
is a rational flat class of type \((1,1)\) there.  Analytic continuation
keeps it of that type wherever its flat continuation is defined.
At a very general point this projection is zero, so it is identically
zero.  Every persistent root is therefore perpendicular to \(U\).
Lemmas~\ref{pencil:monodromy} and
\ref{pencil:descent} now give, over a nonempty open of the original
line, the asserted smooth projective K3 family with the specified
\(U\)-variation.  Its orthogonal complement has dimension fourteen
and is entirely of type \((1,1)\).  Lefschetz's \((1,1)\)-Theorem
and the absence of rational \((1,1)\)-classes in \(U\) show that its
very general transcendental part is exactly \(U\).  This determines
the rational Hodge summand and its pairing; the integral
N\'eron--Severi lattice of the auxiliary K3 surface is not specified
by this construction.

\emph{The quadratic substitution.}
Let \(\Delta\) be the finite union of the critical values and
\(\AA^1_t\setminus O\).  For every \(s\notin\Delta\), make the
substitution \(t=s+u^2\).  Its finite branch value is \(t=s\), a
point of the smooth open \(O\).  It is unramified over every finite
critical value, so those monodromies remain the products in
Lemma~\ref{pencil:monodromy}.  At any other removed finite value the
original period data extend, and the local monodromy is the identity.
At infinity the ramification index is two and kills the square-root
twist.  Thus the full extended K3 lattice has trivial monodromy there:
the square is the identity on \(M\), and the action on \(M^\perp\)
was already the identity.  Every remaining finite monodromy is a
product of integral root reflections fixing \(h\) and the limiting
period vector, by Lemma~\ref{pencil:lattice}.  All local monodromies
therefore satisfy the Weyl condition in
Lemma~\ref{pencil:descent}.  Apply that lemma on \(\PP^1_u\) to
obtain a flat projective ADE K3 family
\(f:\mathcal Z\to\PP^1_u\) and its Hodge line \(\lambda\).

\emph{The Hodge line and uniruledness of the ADE model.}
We first prove uniruledness for this model using its Hodge line, whose
degree is exactly one.  We then identify its function field with the
square pullback of the fixed \(S^0\)-family.

Indeed, the compactified
upstairs double cover is defined over \(B'\times\PP^1_u\) using
\(L\boxtimes\mathcal O_{\PP^1}(1)\).  In homogeneous coordinates
\([u_0:u_1]\) its branch section is
\((d_0+s d_1)u_0^2+d_1u_1^2\), with pullback to \(B'\) understood.
On the affine chart \(u=u_1/u_0\) its equation is
\[
 z^2=d_0+(s+u^2)d_1.
\]
The double-cover formula for the relative dualizing sheaf, with
\(K_{B'}\simeq\mathcal O_{B'}\), gives the anti-sheet part of the
direct image of relative canonical forms as
\[
 H^0(B',L)\otimes\mathcal O_{\PP^1}(1).
\]
Taking \(j=+1\) picks out the fixed one-dimensional eigenspace of
\(H^0(B',L)\) from Proposition~\ref{geom:structure}, hence gives
\(\mathcal O_{\PP^1}(1)\).
Choose a nonzero \(\theta_+\in H^0(B',L)^{j=+1}\).
For completeness, with \(v=1/u\) and \(w=z/u\) the equation becomes
\[
 w^2=d_1+(d_0+s d_1)v^2,
\]
which is smooth at infinity.  Local relative forms transform as
\[
 \frac{\theta_+\,dx\wedge dy}{z}
     =v\,\frac{\theta_+\,dx\wedge dy}{w},
\]
the transition for \(\mathcal O(1)\).  In particular the displayed
frame on the finite chart has a simple zero, measured against the
nonzero frame on the infinity chart, at \(v=0\).
At every finite nodal fiber, Lemma~\ref{pencil:monodromy} shows that
the frame extends after the local resolving cover and is a nonzero
global form on the central resolution.  There is therefore no
additional integral zero or pole at a finite parameter.  The limiting
inertia acts trivially on that form: its vanishing roots are
perpendicular to the limiting period vector.  The invariant-frame
argument in Lemma~\ref{pencil:descent} consequently introduces no
fractional local weight either.  The Hodge-line identification in
Lemma~\ref{pencil:periodstack} includes these exact characters, so
the period realization transfers this extension to \(\lambda\).
We obtain \(\lambda\simeq\mathcal O_{\PP^1}(1)\), as line bundles
with the stated local extensions.

The canonical bundle of this normal polarized contraction model is
\[
 f^*(\mathcal O_{\PP^1}(-2)\otimes\lambda)
      =f^*\mathcal O_{\PP^1}(-1).
\]
Thus \(\deg(K_{\PP^1}\otimes\lambda)=-2+1=-1\), and
Corollary~\ref{pencil:uniruledcriterion} shows that every smooth
projective model of its function field is uniruled.

\emph{The algebraic square-pullback function field.}
We use a common finite cover for this construction and the fixed
construction over the \(t\)-line.  Put \(K=\C(t)\),
\(K_s=\C(u)\) with \(t=s+u^2\), and
\(\widetilde K=\C(\widetilde C)\), where \(\widetilde C\to\PP^1_t\)
is the finite Galois cover chosen in Lemma~\ref{pencil:descent}.
Choose the fields in a common algebraic closure of \(K\), put
\(\widetilde K_s=\widetilde K K_s\), and let \(\widetilde C_s\) be the normalization of
\(\PP^1_u\) in \(\widetilde K_s\).  This is the normalization of an irreducible
component of \(\widetilde C\times_{\PP^1_t}\PP^1_u\).
Since \(\widetilde K/K\) is Galois, so is \(\widetilde K_s/K_s\); its group \(G_s\)
is identified by restriction with a subgroup of \(G=\operatorname{Gal}(\widetilde K/K)\).

This common cover is admissible in the second application of
Lemma~\ref{pencil:descent}.  Indeed, at a point where the original
cover has ramification index \(e\), the substitution has index \(d\),
and the original local monodromy has order \(m\mid e\), normalization
of the fiber product gives index \(e/\gcd(e,d)\).  The pulled-back
monodromy has order \(m/\gcd(m,d)\), which divides that index.
This also covers auxiliary branch points, where \(m=1\).
Thus all pulled-back local monodromies are killed.  We take the
extended period map on \(\widetilde C_s\) to be the pullback of the
previous extended map, with the restricted equivariance data.

Write \(\widetilde{\mathcal Z}_t\to\widetilde C\) for the
algebraized family in the first construction.  Functoriality of the
stack morphism \(q\) identifies the analytic family on
\(\widetilde C_s\), together with its \(G_s\)-action, with the
analytification of the cartesian algebraic pullback
\[
 \widetilde{\mathcal Z}_t\times_{\widetilde C}\widetilde C_s
       \longrightarrow \widetilde C_s.
\]
This pullback is already flat and projective with ADE K3 fibers.
If the new analytic family is algebraized separately as in the lemma,
the fully faithful GAGA comparison used there, applied to the
projective total spaces over \(\widetilde C_s\), algebraizes the
displayed analytic identification and its compatible group action.
We may therefore use this literal algebraic pullback,
with the induced \(G_s\)-action, in forming the second quotient.

Let \(Z_K\) be the generic fiber of the original contraction over
\(K\).  Original descent identifies
\(\widetilde{\mathcal Z}_{t,\widetilde K}\) with \(Z_K\times_K\widetilde K\), including
its \(G\)-descent datum.  After extending scalars to \(\widetilde K_s\) and
restricting to \(G_s\), this is exactly the descent datum of
\(Z_K\times_K K_s\).  Uniqueness of effective finite Galois descent
therefore identifies the generic fiber of the second quotient with
\(Z_K\times_K K_s\).  The fixed smooth generic K3 surface \(S^0_K\)
is a resolution of \(Z_K\); its base change is still birational to
\(Z_K\times_K K_s\), since geometric integrality ensures that the
base change retains the dense open on which the resolution is an isomorphism.
Consequently the second model \(\mathcal Z\) has the function field
of the actual algebraic square pullback of the fixed \(S^0\)-family.
The uniruledness just proved therefore applies to every smooth
projective model of that fixed pulled-back family, as required.
\end{proof}
\section{An algebraic comparison by two square substitutions}
\label{comp:section}

We construct a correspondence between two surface fibers by passing
through the third cohomology of uniruled threefolds.  The first
substitution, $t=s+u^2$, makes degree-three Hodge maps algebraic.  The
second, $s=b+v^2$, produces a reducible conic whose difference of
components recovers degree two with the explicit factor $-4$.
Two points require care: the degree-three comparison must apply to the
particular Leray class furnished by the conic, and the omitted boundary
must have fixed Hodge structures before $b$ varies.  We formulate the
argument for one fixed pencil, including this order of choices.

\begin{theorem}[Two-substitution comparison]\label{comp:twosquares}
Let $O\subset\AA^1_t$ be a nonempty Zariski open set, and let
\[
 f:\mathcal S\longrightarrow O,
 \qquad g:\mathcal X\longrightarrow O
\]
be smooth projective families of connected surfaces, with the fibers of
$f$ K3 surfaces.  Suppose that the following conditions hold.
\begin{enumerate}
\item A polarized rational subvariation $\mathscr U\subset R^2f_*\Q$ of
K3 type, with nondegenerate restricted cup product, is given.
Its period map is nonconstant, and $U_t=\mathscr U_t$ is simple for all
but countably many $t\in O$.
\item There are commuting involutions $j,\delta$ of $\mathcal X$ over
$O$ such that
\[
 e=\tfrac14(1+j)(1-\delta)
\]
annihilates $R^qg_*\Q$ for $q\ne2$, together with an isomorphism of
rational variations of Hodge structure
\[
 \varphi:\mathscr U\xrightarrow{\sim}eR^2g_*\Q.
\]
\item For $s$ in a nonempty Zariski open subset of $\AA^1$, the total
spaces of the two surface families pulled back by $t=s+u^2$ have smooth projective uniruled
threefold models.
\end{enumerate}
Then there are a nonempty Zariski open set $G\subset O$ and a countable
set $E\subset G$ such that, for every $b\in G\setminus E$, a correspondence
\[
 \Theta_b\in\CH^2(S_b\times X_b)_\Q
\]
induces $\varphi_b$ on $U_b$.  The open set $G$ and all auxiliary
correspondences used before the choice of $b$ can be fixed for the given
pencil.
\end{theorem}

The isomorphism $\varphi$ in the theorem is the specified map of
variations; it need not preserve the two cup-product forms.
A cycle on a product acts from its first factor to its second, with the
convention of the Introduction.  In particular, a codimension-three
cylinder from a surface to a fourfold gives
$H^2(-)(-1)\to H^4(-)$, and its transpose gives
$H^4(-)\to H^2(-)(-1)$.  We suppress Tate twists on individual vectors.

\subsection{Third cohomology of uniruled threefolds}

\begin{lemma}\label{comp:uniruled}
If $Z$ is a smooth projective uniruled threefold over $\C$, then $H^3(Z)$
is a direct summand of a finite sum of spaces $H^1(C_i)(-1)$, with both
maps of the splitting induced by algebraic correspondences.  Consequently,
every rational Hodge map between the third cohomologies of two such
threefolds is algebraic.
\end{lemma}

\begin{proof}
Uniruledness gives a generically finite dominant rational map
$T\times\PP^1\dashrightarrow Z$ for a smooth projective surface $T$.
Indeed, choose a two-dimensional parameter space of rational curves
dominating $Z$ and, if necessary, make a generically finite change of its
parameter surface to trivialize the generic rational curve.  Principalize
the base ideal of the rational map in the smooth product, using
projective blowups with smooth centers
\cite[Theorem~1.0.1]{Wlodarczyk}.  This gives a morphism
$a:\widetilde{T\times\PP^1}\to Z$.  The identity
$a_*a^*=(\deg a)\id$ gives an algebraic splitting of $H^3(Z)$ from the
third cohomology of the source.

The product formula and the blowup formula give
\[
 H^3(T\times\PP^1)=H^3(T)\oplus H^1(T)(-1),
\]
and each blowup along a curve adds its $H^1(-)(-1)$; a point adds nothing
in degree three.  The product projections are algebraic, and the full
blowup decompositions are given by algebraic projectors
\cite[Appendix~A, equation~(9) and Proposition~A.1]{LiBlowup};
see also \cite[Section~6.7]{Fulton} for the Chow blowup formalism.
To see the sign in the only new summand,
let $\pi:\widetilde Z\to Z$ be a blowup along a smooth curve $C$,
with exceptional divisor $j:E\hookrightarrow\widetilde Z$ and
$q:E\to C$.  Its inclusion is $j_*q^*$ and its retraction is
$-q_*j^*$, since
\[
 q_*j^*j_*q^*y
 =q_*\bigl(c_1(\mathcal O_E(-1))\smile q^*y\bigr)=-y.
\]
The other summand has inclusion $\pi^*$ and projection $\pi_*$.
These maps and the product projections provide the required algebraic
splittings at every step of the resolution.

For completeness, the needed splittings for a surface use only divisor
classes.  Choose a smooth ample curve $i:C\hookrightarrow T$.
Hard Lefschetz says that multiplication by $[C]$ is an isomorphism
$H^1(T)(-1)\to H^3(T)$.  Since this is $i_*i^*$, restriction
$i^*:H^1(T)\to H^1(C)$ is injective and
$i_*:H^1(C)(-1)\to H^3(T)$ is surjective.
Polarizability gives a Hodge retraction $r$ of $i^*$ and a Hodge section
$s$ of $i_*$.  Poincar\'e duality identifies their tensors with rational
classes of type $(0,0)$ in, respectively,
\[
 H^1(C)\otimes H^1(T)(1)
 \quad\hbox{and}\quad
 H^1(T)\otimes H^1(C)(1).
\]
The underlying untwisted classes have type $(1,1)$ on $C\times T$ and
$T\times C$, so the Lefschetz $(1,1)$ theorem represents them by rational
divisors.  Their actions on the indicated degrees are $r$ and $s$.
Together with restriction and Gysin, they split both surface summands
through $H^1(C)(-1)$.  Taking products with the point and fundamental
classes of $\PP^1$ realizes these splittings in the product decomposition
above.  Composing with $a^*$ and $(\deg a)^{-1}a_*$ proves the first claim.

Finally, a Hodge map between two finite sums of $H^1(C_i)(-1)$ is given
by divisor correspondences on products of the curves.  Composing with the
splittings gives a cycle in $\CH^3(Z\times Z')_\Q$ with exactly the
prescribed action on $H^3$.
\end{proof}

\subsection{Coefficient projections on affine curves}

The degree-three comparison uses the canonical $H^1(R^2)$ term on
the affine square-pullback curves.  We first fix the coefficient
projections on the original curve $O$, so that the same projections
remain available after both substitutions.  We work in the rational
derived category of sheaves on a complex curve; Tate twists record
the Hodge structures on the corresponding cohomology groups.

\begin{lemma}\label{comp:affine}
Let $B$ be a smooth affine complex curve.
\begin{enumerate}
\item A bounded complex on $B$ with finite-rank local-system cohomology
only in even degrees has a unique splitting as the direct sum of its
cohomology sheaves, shifted to their degrees, inducing the identity on
each cohomology sheaf.
\item If $A$ is a bounded complex with local-system cohomology, then
there is a natural exact sequence
\[
 0\longrightarrow H^1(B,\mathcal H^{n-1}A)
 \longrightarrow \mathbb H^n(B,A)
 \longrightarrow H^0(B,\mathcal H^nA)\longrightarrow0.
\]
For a morphism of such complexes, its action on the left-hand term is
determined by its action on $\mathcal H^{n-1}$.
\end{enumerate}
\end{lemma}

\begin{proof}
An affine complex curve has the homotopy type of a one-dimensional CW
complex.  For local systems $L,M$,
\[
 \operatorname{Ext}^r_B(L,M)
   =H^r(B,L^\vee\otimes M)=0\qquad(r>1).
\]
For two distinct nonzero even degrees $i<j$, the obstruction to splitting
the corresponding truncation triangle lies in an $\operatorname{Ext}^{j-i+1}$
group, and the ambiguity lies in an $\operatorname{Ext}^{j-i}$ group.
Their degrees are at least three and two, respectively, so both vanish.
Induction proves (i), including uniqueness with the specified action on
cohomology.  In (ii), the hypercohomology spectral
sequence has only columns zero and one; its degeneration and functoriality
give the assertion.
\end{proof}

For the K3 family of Theorem~\ref{comp:twosquares}, this gives a canonical
splitting of rational complexes
\begin{equation}\label{comp:k3split}
 Rf_*\Q\simeq
 \Q\oplus(R^2f_*\Q)[-2]\oplus\Q(-2)[-4].
\end{equation}
We make this splitting on the original affine curve $O$ and pull it back
under subsequent base changes.  On the other side, the group-algebra
idempotent has a single nonzero cohomology sheaf and therefore gives
canonically
\begin{equation}\label{comp:xsplit}
 eRg_*\Q\simeq (eR^2g_*\Q)[-2]\simeq\mathscr U[-2].
\end{equation}
Use the restricted cup product to write
$p:R^2f_*\Q\to\mathscr U$ for the orthogonal projection.  Pullback of
\eqref{comp:k3split}, followed by $p$, and
\eqref{comp:xsplit}, followed by $\varphi^{-1}$, define coefficient
projections after any base change from $O$.
These are morphisms of rational sheaf complexes.  Their use on a surface
base requires neither a new derived splitting there nor algebraicity of
the projector $p$.

\subsection{The geometric Hodge comparison in degree three}

The next lemma gives the precise coefficient comparison that we need.
It uses a Hodge class on a relative product to construct a Hodge map;
algebraicity of that map will then follow from
Lemma~\ref{comp:uniruled}.

\begin{lemma}\label{comp:hodgelift}
Let $I$ be a smooth affine curve, let $h_S:E^S\to I$ be a smooth
projective K3 family, and let $h_X:E^X\to I$ be a smooth projective
surface family with a relative algebraic idempotent $e$ acting on the direct-image complex whose action on
$R^q h_{X*}\Q$ vanishes for $q\ne2$.  Suppose that
\[
 a:R^2h_{S*}\Q\longrightarrow eR^2h_{X*}\Q
\]
is a morphism of rational Hodge variations.  Let $Z^S,Z^X$ be smooth
projective compactifications which extend the two maps to a smooth
projective compactification of $I$ and are unchanged over $I$.
Write $\overline I$ for that compactified curve.  There is a rational Hodge map
$c:H^3(Z^S)\to H^3(Z^X)$ satisfying
\[
 e\,\operatorname{res}_X(cz)
   =H^1(I,a)\bigl(\operatorname{res}_S z\bigr),
 \qquad z\in H^3(Z^S),
\]
where $H^3(E^S)$ and $eH^3(E^X)$ are identified canonically with
$H^1(I,R^2h_{S*}\Q)$ and $H^1(I,eR^2h_{X*}\Q)$.
\end{lemma}

\begin{proof}
Put $R=E^S\times_I E^X$, and resolve its closure in
$Z^S\times_{\overline I}Z^X$, preserving $R$.
This gives a smooth projective fivefold $\overline R$ with proper maps
$\pi_S:\overline R\to Z^S$ and $\pi_X:\overline R\to Z^X$.
On each fiber, Poincar\'e duality represents $a$ by the unique tensor
in the $H^2(E^S_t)\otimes H^2(E^X_t)$ summand of $H^4(R_t)$ with that
correspondence action.  The tensors form a flat rational section
$\alpha$ of $R^4(R/I)_*\Q$ and have type $(2,2)$.

The global invariant cycle theorem
\cite[Theorem~4.1.1]{DeligneHodgeII} identifies the image of
$H^4(\overline R)\to H^4(R_t)$ with the monodromy invariants.
Restriction is a morphism of polarizable pure Hodge structures.
Its image therefore contains $\alpha_t$, and semisimplicity gives a
rational lift $\widetilde\alpha\in H^4(\overline R)$ of type $(2,2)$.
For example, split restriction onto its image as a Hodge map and apply
the splitting to $\alpha_t$.  The restriction of this lift to every
fiber is $\alpha$, because both define flat sections.
Define
\[
 c(z)=(\pi_X)_*
       \bigl(\widetilde\alpha\smile\pi_S^*z\bigr).
\]
Cup product raises degree by four and the proper pushforward lowers
degree by four, so this is a rational Hodge map from $H^3$ to $H^3$.

Its restriction to $I$ is the relative cohomological correspondence
defined by $\widetilde\alpha|_R$.  Pullback, cup product, and integration
along the smooth proper relative surface projection give a morphism
$Rh_{S*}\Q\to Rh_{X*}\Q$; on degree-two cohomology sheaves it is $a$.
For the source, $R^3h_{S*}\Q=0$, and for the projected target,
$eR^3h_{X*}\Q=0$.  Lemma~\ref{comp:affine}(ii), with $n=3$,
therefore identifies the entire source and projected target $H^3$ with
their canonical $H^1(R^2)$ terms.  Functoriality shows that the
projected open action of $c$ is exactly $H^1(I,a)$, as asserted.
The proof uses no claim that $\widetilde\alpha$ is algebraic.
\end{proof}

\subsection{Compatible projective models}

The following construction will be used for both surface families.
It fixes the first compactifications before the relative correspondence
and the point $b$ are chosen.

\begin{lemma}\label{comp:models}
For a smooth projective surface family $\mathcal Y\to O$, its first
square pullbacks admit models $Z$ projective over $\AA^1_s$ and a nonempty
open set $G\subset O$ such that $Z_G\to G$ is smooth.  These models may
be chosen with the following additional property.  For every $b\in G$
there is a smooth projective fourfold $M_b$, with morphisms to
$\PP^1_t$ and $\PP^1_v$, whose inverse images of $V_0\subset\PP^1_v$
and $O\subset\PP^1_t$ are, respectively,
\begin{enumerate}
\item $Z_G\times_GV_0$, where $V_0\to G$ is $v\mapsto b+v^2$ and
$V_0=\{v\in\AA^1:b+v^2\in G\}$;
\item $\mathcal Y\times_OQ_b$, where
\begin{equation}\label{comp:conic}
 Q_b=\{u^2+v^2=(t-b)w^2\}\subset O\times\PP^2.
\end{equation}
\end{enumerate}
The two identifications agree on their intersection.  In particular,
both specified spaces are preserved as actual open subvarieties.
Shrinking $G$ before choosing $b$ preserves these conclusions.
\end{lemma}

\begin{proof}
Choose a smooth projective compactification
$\overline{\mathcal Y}\to\PP^1_t$ whose inverse image over $O$ is
$\mathcal Y$.  It is obtained by taking a projective closure of the
family and resolving away from its smooth open.  Take the component containing the given open. Its source is smooth and
integral, and its map to the nonsingular curve is flat. The generic fiber
is geometrically integral: a smooth geometrically connected surface is
geometrically integral. Hence the base changes below have the dense
unchanged open and the same function field as the indicated pullback.  The map
\[
 C=\PP^1_u\times\AA^1_s\longrightarrow\PP^1_t,
 \qquad t=s+u^2,
\]
is a morphism: in homogeneous coordinates on $\PP^1_u$, with $u=U_0/U_1$, it is
$[U_0:U_1]\mapsto[U_0^2+sU_1^2:U_1^2]$.
Set
$W=\overline{\mathcal Y}\times_{\PP^1_t}C$.  Choose a projective
resolution $Z\to W$ as a sequence of blowups, all supported over
$t\notin O$.  The flat base change is generically smooth and therefore
reduced, with no component supported only over the boundary; its open
over $O$ is smooth.  Resolution is performed on each component if needed.
Strong projective resolution in characteristic zero, preserving the
regular locus, permits this construction
\cite[Theorems~1.0.2--1.0.3]{Wlodarczyk}.
The map $Z\to\AA^1_s$ is projective,
and is smooth after deleting finitely many values of $s$.  Fix this entire
construction, and take $G$ inside its smooth locus and inside $O$.

Let $B_b$ be the normalization of $\PP^1_t\times\PP^1_v$ in the field
\[
 \C(t,v)\bigl(\sqrt{t-b-v^2}\bigr).
\]
It is a projective surface finite over the product.  Its finite-$v$ open
is $\PP^1_u\times\AA^1_v$, with $t=b+u^2+v^2$.  Its open over $O$ is
$Q_b$: the latter is smooth and the map
\[
 Q_b\longrightarrow O\times\PP^1_v,
 \qquad (t,[u:v:w])\longmapsto(t,[v:w])
\]
has no base point: $v=w=0$ would force $u=0$.
It is proper and has finite fibers, hence is finite of degree two, with
the required function field.  The same verification on the finite-$v$
open uses $\PP^1_u\times\AA^1_v$ and the morphism
$[U_0:U_1]\mapsto[U_0^2+(b+v^2)U_1^2:U_1^2]$.
These assertions identify both opens by the universal property
of normalization.  Smoothness of $Q_b$ follows in the affine chart $w=1$
from the derivative in $t$, and at infinity from the derivatives in
$u,v$.

Start with the projective fourfold
$W_b=\overline{\mathcal Y}\times_{\PP^1_t}B_b$.  On its finite-$v$ open,
pull back the chosen blowup sequence $Z\to W$ by $s=b+v^2$.  This map of
affine lines is finite flat, so blowup commutes with the base change.
Each blowup extends to the whole current projective model without
altering its open over $O$.  Explicitly, if $j:U\hookrightarrow W_b^{(r)}$
is its finite-$v$ open and $\mathcal I\subset\mathcal O_U$ is the next
center ideal, extend it by
\[
 \mathcal J=\ker\bigl(\mathcal O_{W_b^{(r)}}
          \longrightarrow j_*(\mathcal O_U/\mathcal I)\bigr).
\]
This is a coherent ideal: the schemes are noetherian, the open immersion
is quasi-compact, and the kernel is a quasi-coherent submodule of
$\mathcal O_{W_b^{(r)}}$.  It restricts to $\mathcal I$ and is the unit
ideal off the closed inverse image of $\PP^1_t\setminus O$.  Thus its
blowup has the asserted properties.

After extending the finite sequence, both opens in the statement are
present and smooth.  A final projective resolution which is an
isomorphism on the regular locus preserves their union.  It gives $M_b$
and the two morphisms to the parameter lines.  The construction uses one
fixed set of discarded $s$-values, independent of $b$.
\end{proof}

\subsection{Proof of the comparison theorem}

\begin{proof}[Proof of Theorem~\ref{comp:twosquares}]
We retain the notation $\mathscr U$, $p$, and $\varphi$ above.

\paragraph{\emph{The fixed part.}}
We first prove the vanishing that will make each cylinder class restrict
to zero on the projective threefold fibers:
\begin{equation}\label{comp:noinvariants}
 H^0(O,\mathscr U)=0.
\end{equation}
By the fixed-part theorem, the invariant part of the geometric polarized
variation is a constant Hodge subvariation; see
\cite[Theorem~4.1.1 and Corollary~4.1.2]{DeligneHodgeII}.  A nonzero such
part would equal $U_t$ at a simple fiber, and hence everywhere as a local
system.  The fixed-part theorem would then make its period constant,
contrary to the hypothesis.

For $s\in O$, put $I_s=\{u\in\AA^1:s+u^2\in O\}$.  The pullback of
$\mathscr U$ to $I_s$ has the same monodromy image as $\mathscr U$ on
$O$.  To see this, remove the branch value $s$ temporarily.  A loop about
$s$ changes the sheet of the double cover and has trivial monodromy on
$\mathscr U$; it can therefore be used to adjust the parity of any loop
without changing its monodromy.  Consequently,
\begin{equation}\label{comp:noinvariantsI}
 H^0(I_s,\mathscr U)=0\qquad(s\in O).
\end{equation}

\paragraph{\emph{A relative correspondence in degree three.}}
Apply Lemma~\ref{comp:models} to both families, obtaining models
$Z^S\to\AA^1_s$ and $Z^X\to\AA^1_s$.  Shrink a common nonempty open
$G\subset O$ so that both families $\zeta^k:Z_G^k\to G$ are smooth
projective and their fibers are uniruled.  Write $E_s^S$ and $E_s^X$ for
their unchanged opens
over $I_s$.  The Leray spectral sequence over the affine curve $I_s$
and the degree-two projections give maps
\begin{align*}
 r_s^S &:H^3(Z_s^S)\longrightarrow
 H^3(E_s^S)=H^1(I_s,R^2f_*\Q)
 \xrightarrow{p}H^1(I_s,\mathscr U),\\
 r_s^X &:H^3(Z_s^X)\longrightarrow
 eH^3(E_s^X)=H^1(I_s,eR^2g_*\Q)
 \xrightarrow{\varphi^{-1}}H^1(I_s,\mathscr U).
\end{align*}
All coefficients in these expressions have been pulled back to $I_s$.

Here is the Hodge-theoretic content of this comparison.  Give
$H^1(I_s,\mathscr U)$ the mixed Hodge structure transported from
$eH^3(E_s^X)$ by the second identification above.  The weight-three part
\[
 T_s=W_3H^1(I_s,\mathscr U)
\]
is pure and polarizable.  Restriction from a smooth projective
compactification has image equal to the lowest weight
\cite[Corollary~3.2.17]{DeligneHodgeII}; the algebraic projector on the
open commutes with the weight filtration.  Thus $r_s^X$ is a surjective
Hodge map onto $T_s$.  This argument uses $e$ only on $E_s^X$.

Apply Lemma~\ref{comp:hodgelift} with $a=\varphi p$, after pullback to
$I_s$.  It gives a rational Hodge map
\[
 c_s:H^3(Z_s^S)\longrightarrow H^3(Z_s^X),
 \qquad r_s^Xc_s=r_s^S.
\]
In particular, $r_s^S$ is also a Hodge map into this same pure space
$T_s$.  This proves the coefficient compatibility at exactly the
strength required, using geometric invariant cycles and the canonical
Leray term over a curve.  Lemma~\ref{comp:uniruled} now realizes $c_s$
by a cycle in $\CH^3(Z_s^S\times Z_s^X)_\Q$.

We next spread the cycles just obtained.  Write
$\Delta=\AA^1\setminus O$.  The punctures of $I_s$ are the roots of
$u^2=c-s$ for $c\in\Delta$, together with infinity in its compactification.
Because $s\in G\subset O$, they remain distinct.  Locally in $s$ their
motion gives a topologically trivial punctured curve with its pulled-back
coefficient system.  Consequently $H^1(I_s,\mathscr U)$ forms a local
system over $G$, and $r^S,r^X$ are morphisms of local systems.

The projective family $Z^S_G\times_G Z^X_G$ has relative dimension six.
Use the relative Hilbert schemes for polynomials of degree three to
parameterize its codimension-three subschemes, and take all finite
products together with all rational coefficient lists
\cite[Theorem~3.2]{GrothendieckHilbert}.
This is a countable collection of finite-type parameter spaces and
includes a representative of every rational codimension-three cycle.
Stratify them into connected smooth locally closed varieties.
The degree-six fundamental cycle class of a universal flat subscheme
is locally constant in the cohomology of the ambient fibers.
Indeed, over a smooth stratum the ambient total space is smooth, so
the universal structure sheaf has a finite locally free resolution.
Flatness over the stratum makes derived restriction to any fiber equal
to ordinary restriction.  The degree-six part of its Chern character
restricts to the fundamental cycle class, with generic multiplicities.
This follows from Grothendieck--Riemann--Roch, d\'evissage, and
the compatibility of algebraic and topological Chern classes
\cite[Theorem~15.2, Example~15.2.16\textup{(a)}, and Proposition~19.1.2]{Fulton}.  Since these
characteristic classes come
from a class on the total space, their fiber restrictions are flat.
Lemma~\ref{final:hilbert} proves the corresponding assertion for the
cycle parameter spaces used in specialization.  Thus on each stratum
the equality
\begin{equation}\label{comp:rcomparison}
 r^X C_*=r^S\quad\hbox{on }R^3\zeta^S_*\Q
\end{equation}
holds everywhere or nowhere.  Successful strata cover $G$ pointwise,
since a cycle has been constructed for every $s\in G$.  A
nondominating finite-type stratum has finite image in the irreducible
curve $G$.  If all successful strata were nondominating, their images
would be countable, whereas $G(\C)$ is uncountable.

Choose a successful dominating stratum.  A closed point of its generic
fiber gives a generically finite multisection; normalize its closure
and shrink $G$ so that it is a finite \emph{\'etale} cover
$\nu:G'\to G$, of degree $d$.  The universal rational cycle restricts
to a cycle $C'$ over $G'$.  Its trace
\begin{equation}\label{comp:relativeC}
 C=\frac1d\nu_*C'
 \in\CH^3(Z^S_G\times_GZ^X_G)_\Q
\end{equation}
satisfies \eqref{comp:rcomparison}: over every $s$, its action is the
average of the $d$ successful actions over $\nu^{-1}(s)$.
Generic flatness permits a final shrinking so that the chosen cycle
representatives are flat over $G$.  From this point onward, fix $G$,
both models, and this relative cycle $C$.

\paragraph{\emph{The second substitution and its cylinders.}}
Choose $b\in G$, subject to the countable exclusions specified below,
and set $s=b+v^2$.  Lemma~\ref{comp:models} supplies smooth projective
fourfolds $M^S,M^X$ with the two required opens.  Let
\[
 V_0=\{v\in\AA^1:b+v^2\in G\},\qquad
 Y^k=\mathcal Y^k\times_OQ_b\quad(k=S,X),
\]
where $\mathcal Y^S=\mathcal S$ and $\mathcal Y^X=\mathcal X$.
In the chart $w=1$, put
\[
 Q_b^\circ=\{(u,v)\in\AA^2:v\in V_0,\ b+u^2+v^2\in O\}.
\]
The substitutions are $t=s+u^2$ and $s=b+v^2$.  The two projections
$\pi_t:Q_b\to O$ and $\pi_v:Q_b\to\PP^1_v$ are given by $t$ and
$[v:w]$.  For each $k=S,X$, the preserved opens form the square
\[
 \begin{array}{ccc}
 \mathcal Y^k\times_OQ_b^\circ &\hookrightarrow& Z_G^k\times_GV_0\\
 \big\downarrow && \big\downarrow\\
 Y^k=\mathcal Y^k\times_OQ_b &\hookrightarrow& M^k.
 \end{array}
\]
All arrows are open immersions, and the upper-left open is the
intersection of the two opens in $M^k$.  The relative product over $V_0$
of the two upper-right opens carries the pullback of the relative cycle
$C$.  The lower-left open contains the complete cylinders
$F_b^k\times L_\pm$ defined below, and the upper-left open is the
coefficient-comparison locus.  The choice order is
summarized by
\begin{center}
\small
\begin{tabular}{@{}p{.45\linewidth}p{.45\linewidth}@{}}
\toprule
Fixed before $b$ & Depending on $b$\\
\midrule
Original pencil, $G$, $Z_G^S,Z_G^X$, relative cycle $C$
& $Q_b,V_0,M^k$ and the cycles $\widehat C,D_k$ constructed below\\
\bottomrule
\end{tabular}
\end{center}
Pull \eqref{comp:relativeC} back to $V_0$ and close its cycle
representatives in $M^S\times_{\PP^1_v}M^X$.  Viewed in the absolute
product, the result is a codimension-four cycle $\widehat C$, hence acts
degree-preservingly.  Its support is over the diagonal of the $v$-line.
Restriction of its action to $V_0$ is consequently the original relative
action: restricting the target restricts the source as well, and proper
pushforward commutes with this open restriction.

The fiber of $Q_b\to O$ at $b$ is the union of the two lines
\[
 L_+=\{u=iv\},\qquad L_-=\{u=-iv\}.
\]
Write $F_b^S=S_b$ and $F_b^X=X_b$ for the original surface fibers.
The inclusions $i_{k,\pm}:F_b^k\times L_\pm\hookrightarrow M^k$,
together with the projections $p_{k,\pm}:F_b^k\times L_\pm\to F_b^k$,
define cylinder correspondences from $F_b^k$ to $M^k$.
Let $D_k$ be their difference, with the plus component first.  Thus
\[
 D_k\in\CH^3(F_b^k\times M^k)_\Q,\qquad
 d_k=(i_{k,+})_*p_{k,+}^*-(i_{k,-})_*p_{k,-}^*,
\]
and its cohomological action is
\[
 d_k:H^2(F_b^k)(-1)\longrightarrow H^4(M^k).
\]
The full lines, including their points at infinity, are contained in the
preserved open $Y^k$.  Since these cylinders are projective, they are
closed in $M^k$.

In the chart $w=1$, the two lines meet at $(u,v)=(0,0)$ with independent
tangent directions.  Their intersection multiplicity is therefore one.
Their sum is the divisor of the base parameter $t-b$, whose restriction
to either compact line has degree zero.  Thus, on the smooth surface
$Q_b$,
\[
 L_+\cdot L_-=1,\qquad L_+^2=L_-^2=-1.
\]
The normal bundles of the cylinder divisors are pulled back from these
normal bundles on $Q_b$.  The self-intersection and excess-intersection
formulas \cite[Theorem~6.3, Corollary~6.3, and Theorem~19.2]{Fulton}
therefore give
\begin{equation}\label{comp:minusfour}
 {}^td_k\,d_k=-4\id
 \quad\hbox{on }H^2(F_b^k)(-1).
\end{equation}
More explicitly, each self-term is the identity on $H^2(F_b^k)$ multiplied
by $\deg N_{L_\pm/Q_b}=-1$, and each cross-term is the identity multiplied
by the intersection multiplicity one.  The difference convention gives
$(-1)+(-1)-1-1=-4$.  Every operation on the surface factor is the
identity, so the formula holds on all of $H^2(F_b^k)$.

\paragraph{\emph{Coefficient classes of the cylinders.}}
The relative comparison acts on coefficient cohomology over the affine
conic chart.  We must therefore identify the coefficient images of the
two cylinder classes.  For either complete component $L_\pm\subset Q_b$,
its inverse image in $Y^k$ is $F_b^k\times L_\pm$.  We shall show that
the coefficient projection carries its total-space Gysin class to the
Gysin class of the same coefficient on $Q_b$.  We verify the compatibility
on a general surface base before restricting to the comparison chart.

Let $B_2$ be a smooth surface with a map
$\tau:B_2\to O$, and let $\sigma:D\hookrightarrow B_2$ be a smooth divisor
whose map to $O$ is constant, with value $b$.  Write
$\mathcal S_{B_2}=\mathcal S\times_OB_2$ and similarly for $\mathcal X_{B_2}$.
The inverse image of $D$ is $S_b\times D$, respectively $X_b\times D$.
For $x\in U_b$, its pullback to this inverse image has a Gysin class in
$H^4(\mathcal S_{B_2})$, with domain understood as $U_b(-1)$.  Under the
coefficient projection its image is
\begin{equation}\label{comp:thom}
 \sigma_*(x)\in H^2(B_2,\tau^*\mathscr U).
\end{equation}
The same formula holds for the $e$-part on $\mathcal X_{B_2}$, with coefficient
$\varphi_b(x)$.

To verify the formula, make the coefficient Thom construction explicit.
The complex orientation of the normal line of $\sigma$ gives a morphism
\[
 u_\sigma:\sigma_*\Q_D[-2]\longrightarrow\Q_{B_2}.
\]
On a normal disc, the relative cochain complex of the disc and its
punctured disc has one cohomology group, generated in degree two by
this orientation.  Tensoring this local calculation with a bounded
locally constant rational complex $K$ gives the natural purity
identification $\sigma^!K\simeq\sigma^*K[-2]$.  More precisely, the
globally defined map
\[
 \mu_K:\sigma^*K\otimes\sigma^!\Q_{B_2}\longrightarrow\sigma^!K
\]
is adjoint to the projection-formula/counit composite
\[
 \sigma_*(\sigma^*K\otimes\sigma^!\Q_{B_2})
 \simeq K\otimes\sigma_*\sigma^!\Q_{B_2}\longrightarrow K.
\]
The normal-disc calculation shows that $\mu_K$ is an isomorphism for
such $K$, and the complex orientation identifies
$\sigma^!\Q_{B_2}\simeq\Q_D[-2]$.  This gives the stated natural purity
isomorphism globally.  Equivalently, the coefficient Gysin morphism is
\[
 \sigma_*\sigma^*K[-2]
   \simeq (\sigma_*\Q_D[-2])\otimes K
   \xrightarrow{\ u_\sigma\otimes\id\ }K.
\]
The first isomorphism is the projection formula, and this construction
is natural in $K$.  In the rational sheaf category the displayed shifts
are only $[-2]$; the Tate $(-1)$ records the Hodge degree of the complex
normal orientation on cohomology.

For completeness, let $h:Y\to B_2$ be either smooth proper surface family,
let $Y_D=Y\times_{B_2}D$, and write $j:Y_D\hookrightarrow Y$ and
$h_D:Y_D\to D$.  The normal line of $j$ is $h_D^*N_{D/B_2}$ with its
pulled complex orientation.  The normal-disc construction therefore
identifies the total-space Thom morphism with the pullback of
$u_\sigma$.  Under the natural proper base-change identification
\[
 Rh_*j_*\Q_{Y_D}[-2]
   \simeq\sigma_*Rh_{D*}\Q_{Y_D}[-2]
   \simeq\sigma_*\sigma^*Rh_*\Q_Y[-2],
\]
the projection formula consequently identifies $Rh_*(u_j)$ with the
coefficient Gysin morphism for $K=Rh_*\Q_Y$.  The formal base-change and
projection-formula statements used here are natural for these bounded
rational complexes \cite[Theorems~5.10 and~6.3]{SchnuererSoergel}.

Finally $D\to O$ factors through $b$.  Proper base change identifies
the restricted splitting with the canonical splitting of
$R\Gamma(S_b,\Q)$, respectively $eR\Gamma(X_b,\Q)$, pulled back from a
point.  The class pulled back from $x$ is therefore its degree-two product
coefficient.  Naturality of the coefficient Gysin map makes it commute
with the splitting pulled from \eqref{comp:k3split}, its degree-two
projection, and the idempotent \eqref{comp:xsplit}.  This proves
\eqref{comp:thom}.  The construction applies to each divisor separately,
including when two such divisors intersect.

\paragraph{\emph{Comparison on the affine conic chart.}}
For $x\in U_b$, put $a=\widehat C_*d_S(x)$.  Our immediate aim is
to prove $e\,a=d_X(\varphi_bx)$ on $\mathcal X\times_OQ_b^\circ$,
using only the degree-three action of $C$.
On $Q_b^\circ$, write $q:Q_b^\circ\to V_0$ for projection to $v$ and
$\tau(u,v)=b+u^2+v^2$.
Put
$c=d_S(x)|_{Z^S_G\times_GV_0}$.  The restriction of $c$ to every
threefold fiber over $v$ is zero.  Each cylinder meets that fiber
transversally: the map $L_\pm\to\PP^1_v$ is an isomorphism.
Its restriction is therefore the ordinary surface Gysin map at
$u=\pm iv$, with no excess class, including at $v=0$.
Pair either surface Gysin image of $x$ with an arbitrary class in
the second cohomology of the threefold.  This testing class restricts
to a monodromy-invariant class on surface fibers over
$I_{b+v^2}$.  Its orthogonal projection to $\mathscr U$ is zero by
\eqref{comp:noinvariantsI}, so the pairing is zero.  Poincar\'e duality on
the projective threefold proves the claimed vanishing of each term.
At $v=0$ the two surface inclusions also coincide.

Write $\rho_k:Z^k_G\times_GV_0\to V_0$ and
$A_k=R\rho_{k*}\Q$.  Lemma~\ref{comp:affine}(ii) places $c$ in the
canonical Leray subspace
\begin{equation}\label{comp:lerayone}
 H^1(V_0,\mathcal H^3A_S)
 \subset\mathbb H^4(V_0,A_S).
\end{equation}
Set $K=Rq_*\tau^*\mathscr U[-2]$.  Its cohomology sheaves are local systems,
by the same noncollision of punctures used over $G$, now pulled back to
$V_0$.  Restriction to the open surface family
over $Q_b^\circ$, followed by the coefficient projections
\eqref{comp:k3split}--\eqref{comp:xsplit}, gives morphisms of rational
complexes
\[
 \beta_S:A_S\longrightarrow K,
 \qquad \beta_X:A_X\longrightarrow K.
\]
The relative correspondence gives $\gamma:A_S\to A_X$: on the
relative product, take cup product with its class of degree six and
integrate along the proper relative threefold projection.  Its action
on $R^3$ is the action prescribed by $C$.
Equation~\eqref{comp:rcomparison} says exactly that
\[
 \mathcal H^3(\beta_X\gamma)=\mathcal H^3(\beta_S).
\]
Functoriality in Lemma~\ref{comp:affine}(ii), applied to the subspace
\eqref{comp:lerayone}, now yields
\begin{equation}\label{comp:derivedcomparison}
 (\beta_X\gamma)_*c=(\beta_S)_*c
 \quad\hbox{in }\mathbb H^4(V_0,K)=H^2(Q_b^\circ,\tau^*\mathscr U).
\end{equation}
The reason degree-three agreement suffices is precisely the membership
\eqref{comp:lerayone}: the left term of the functorial exact sequence
in Lemma~\ref{comp:affine}(ii) only sees $\mathcal H^3$.
No equality on the $H^0(\mathcal H^4)$ quotient is used.
The $u$-fibers and the $v$-base have cohomological dimension one, so the
last group is also canonically
\[
 H^1\bigl(V_0,R^1q_*\tau^*\mathscr U\bigr).
\]

The restrictions of the cylinders to $Q_b^\circ$ are the two smooth section
divisors $\sigma_\pm(v)=(\pm iv,v)$.  Formula~\eqref{comp:thom} shows
that the right side of \eqref{comp:derivedcomparison} is
\[
 \sigma_{+*}(x)-\sigma_{-*}(x).
\]
The same formula on the $X$ side identifies this with the coefficient
class of $d_X(\varphi_bx)$.  Although the sections meet at $v=0$, purity
has been applied to each regular embedding separately.  By \eqref{comp:xsplit}, the
coefficient projection identifies the entire $e$-part of the fourth
cohomology over $Q_b^\circ$ with $H^2(Q_b^\circ,\tau^*\mathscr U)$.  Hence
we have proved
\begin{equation}\label{comp:smallopen}
 e\,a=d_X(\varphi_bx)
 \quad\hbox{on }\mathcal X\times_OQ_b^\circ.
\end{equation}

\paragraph{\emph{Excluding the boundary error.}}
We must prove \eqref{comp:smallopen} on all of $Y^X$, since the complete
cylinders also contain their points at infinity.  Let
$F=\AA^1\setminus G$, a fixed finite set.  The complement of
$\mathcal X\times_OQ_b^\circ$ in $Y^X$ is a disjoint union of smooth divisors:
\begin{enumerate}
\item the two infinity sections $[u:v:w]=[\pm i:1:0]$, each carrying the
fixed family $\mathcal X\to O$;
\item for each $c_0\in F$ and each of the two roots $a_0$ of
$a_0^2=c_0-b$, the divisor $v=a_0$, carrying the fixed open threefold
\[
 E_{c_0}^X
   =\{(x,u):g(x)=c_0+u^2\in O\}.
\]
\end{enumerate}
The second description follows by substituting $v=a_0$ in the affine
conic equation.  The projective equation of a finite divisor is
$v=a_0w$; a point with $w=0$ on it would also have $v=u=0$, which is
impossible.  Thus it has no point on either infinity section.
Different finite values of $v$ give disjoint divisors, and all the roots
listed above are distinct since $b\in G$.
Consequently the entire boundary is a disjoint union, with the fixed
isomorphism types just listed.
Their smoothness follows from smoothness of $g$ over $O$, including when
$c_0$ is not a good value for a compactified threefold.

Denote their union by $D$.  Topological localization and the
complex-oriented Thom isomorphism for each smooth closed divisor give
the underlying exact sequence
\begin{equation}\label{comp:localization}
 H^2(D)(-1)\longrightarrow H^4(Y^X)
 \longrightarrow H^4(\mathcal X\times_OQ_b^\circ).
\end{equation}
This is exact in mixed Hodge structures as well.  The closed immersion
$D\hookrightarrow Y^X$ is proper even though both varieties may be
nonproper.  For smooth separated complex algebraic varieties, the
adjunction and duality constructions of restriction and proper Gysin
lift to mixed Hodge modules; the smooth duality normalization gives
exactly the codimension-one shift $[-2](-1)$
\cite[Sections~1.2--1.4]{SaitoMHMIntro}.  Their cohomology mixed Hodge
structures agree with Deligne's
\cite[Theorems~0.1--0.2 and the following paragraph]{SaitoMHC}.
Thus the two arrows in \eqref{comp:localization} are morphisms for the
canonical mixed Hodge structures.  Their underlying sequence is exact
by the Thom calculation, hence it is exact in that abelian category.
The first term's weight-four part is isomorphic, independently of $b$,
to $H_F(-1)$, where the fixed polarizable pure Hodge structure $H_F$ is
\[
 H_F=
 \bigl(W_2H^2(\mathcal X)\bigr)^{\oplus2}
 \oplus
 \bigoplus_{c_0\in F}
 \bigl(W_2H^2(E_{c_0}^X)\bigr)^{\oplus2}.
\]
Here these spaces are polarizable because lowest-weight cohomology of a
smooth variety is an image of the cohomology of a smooth projective
compactification.

For all but countably many $b\in G$,
\begin{equation}\label{comp:fixedtypes}
 \Hom_{\mathrm{HS}}(U_b,H_F)=0.
\end{equation}
To check this precisely, exclude the countably many nonsimple fibers and
trivialize $\mathscr U$ on a countable cover by simply connected discs.
A nonzero map from a remaining $U_b$ is injective.  There are countably
many rational linear injections into the fixed rational vector space
$H_F$.  For each such injection, the condition of being a Hodge map
determines the entire Hodge filtration on $U_b$ uniquely as the inverse
image of the filtration on $H_F$.  The nonconstant period map meets this
one fixed filtration in a discrete set on a disc.  Taking the indicated
countable unions proves \eqref{comp:fixedtypes}.

Consider the error map
\[
 U_b(-1)\longrightarrow H^4(Y^X),
 \qquad x\longmapsto e\,a-d_X(\varphi_bx).
\]
Both terms are morphisms of Hodge structures: the cylinder and
$\widehat C$ act algebraically on the projective models, restriction is
a mixed-Hodge morphism, $e$ is algebraic on $Y^X$, and $\varphi_b$ is
Hodge by hypothesis.  Equation~\eqref{comp:smallopen} places the error
in the kernel of the last arrow of \eqref{comp:localization}.
Its source is pure of weight four.  Strictness for the weight filtration
\cite[Theorem~2.3.5\textup{(iii)}]{DeligneHodgeII} therefore puts its image in the
Gysin image of $W_4(H^2(D)(-1))\simeq H_F(-1)$.
This image is a quotient of a polarizable pure Hodge structure.
Semisimplicity splits that quotient, lifting the error map to
$U_b(-1)\to H_F(-1)$; by
\eqref{comp:fixedtypes} the lift is zero.  We have proved
\begin{equation}\label{comp:largeopen}
 e\,\widehat C_*d_S(x)=d_X(\varphi_bx)
 \quad\hbox{in }H^4(Y^X),\qquad x\in U_b.
\end{equation}

\paragraph{\emph{Returning to the surfaces.}}
The transpose cylinder can be evaluated on a class of $Y^X$ by
restriction to $X_b\times L_\pm$ and integration along $L_\pm$.
These are complete projective cylinders inside $Y^X$, so this operation
agrees with its action on classes restricted from $M^X$.
Moreover,
$e_b\,{}^td_X={}^td_X e$, since $j,\delta$ act only on the surface
coordinate.  Define the algebraic correspondence
\begin{equation}\label{comp:theta}
 \Theta_b=-\tfrac14\,e_b\circ{}^tD_X\circ\widehat C\circ D_S
 \in\CH^2(S_b\times X_b)_\Q.
\end{equation}
Applying the transpose to \eqref{comp:largeopen} and using
\eqref{comp:minusfour} gives, for every $x\in U_b$,
\[
 (\Theta_b)_*x
 =-\frac14\,{}^td_X e\,\widehat C_*d_S(x)
 =-\frac14\,{}^td_Xd_X(\varphi_bx)
 =\varphi_bx.
\]
In this calculation $e$ acts after restriction to $Y^X$.
Thus the action is exactly the specified map on the whole of $U_b$.
The codimension is $3+4+3-4-4=2$; composing with
$e_b\in\CH^2(X_b\times X_b)_\Q$ preserves it.
Every exceptional choice of $b$ used after fixing $G$ belongs to the countable set excluded
in \eqref{comp:fixedtypes}, together with the nonsimple fibers.  This
proves the theorem with the asserted order of choices.
\end{proof}

\subsection{Application to the abelian-cover pencil}

\begin{corollary}\label{comp:surface}
Fix abelian-cover data and a section pencil for which
\Cref{geom:structure,geom:period,pencil:realization} apply, with nonconstant $U$-period along
the pencil and simple $U_t$ outside a countable set.  On a nonempty open
of that pencil, for all but countably many $b$, there is an algebraic
correspondence
\[
 \Theta_b\in\CH^2(S_b^0\times X_b)_\Q
\]
whose action on $U_b$ is the given identification with
$H^2(X_b)^{j=+,\,\delta=-}$.
\end{corollary}

\begin{proof}
The variation, the projector, and the identification are supplied by
\Cref{geom:structure,pencil:realization}.
\Cref{pencil:realization} gives the uniruled threefold models on the K3 side.
For $X$, after $t=s+u^2$ the generic fiber over $B'$ is the conic
\[
 z^2-\pi^*d_1\,u^2=\pi^*(d_0+s d_1)
\]
over $\C(B')$.  A finite extension splits this conic, so a generically
finite surface cover of $B'$ yields a dominant rational map from a
product of a surface and $\PP^1$.  Every smooth projective model of the
threefold is therefore uniruled.  Theorem~\ref{comp:twosquares} applies.
It applies separately to every fixed pencil with these properties; its
countable set of excluded $b$ is chosen only after the pencil, the
degree-three correspondence, and the good $s$-open have been fixed.
\end{proof}
\section{Recovery of the prescribed Kuga--Satake tensor}
\label{final:section}

The comparison of the preceding section and the domination by curves
give an algebraic injection from the auxiliary K3 cohomology into a
tensor product of two curve cohomologies.  This does not yet realize
the standard tensor required by Theorem~\ref{ext:criterion}.  We bridge
the difference by proving that its compositions with rational Hodge maps
on the two curve factors, together with their interchanges, span every
Hodge map into the tensor square of the full even Clifford module.
This permits us to select the prescribed right factor and polarization.
The calculation applies at the construction periods selected in
Lemma~\ref{final:curves}, where the Hodge group is the full special
orthogonal group.
At the end of the section we transfer the curve source by a rational
Hodge isometry; Section~\ref{specialization:section} will use that transfer
to obtain seeds with $\NS=P$ and extend one fixed tensor to every ample
$P$-polarized target.

\subsection{An algebraic embedding into a tensor of curve cohomologies}

\begin{lemma}\label{final:curves}
Choose the abelian data and the pencil as in
Lemma~\ref{geom:genericity}, and then choose $b$ in the common good open
outside the countable exceptional sets of that lemma and
Corollary~\ref{comp:surface}.  Then there are smooth projective curves
$C_1,C_2$ and a rational algebraic correspondence
\[
 Z\in\CH^2(S_b^0\times C_1\times C_2)_{\Q}
\]
whose action restricts to an injection
\[
 f:U_b\longrightarrow H^1(C_1,\Q)\otimes H^1(C_2,\Q).
\]
Here $U_b=T_{S_b^0}$ and $\Hg(U_b)=\SO(U_b)$.
\end{lemma}

\begin{proof}
Let $\Theta_b\in\CH^2(S_b^0\times X_b)_\Q$ be the correspondence
of Corollary~\ref{comp:surface}; its restriction to $U_b$ is the
specified identification $\varphi_b:U_b\to U_{X_b}$.
At these periods $U_{X_b}$ has no nonzero rational $(1,1)$-class
by Proposition~\ref{geom:period}. Proposition~\ref{geom:domination}
therefore supplies curves $C_1,C_2$, a morphism
$g:C_1\times C_2\to X_b$, and the algebraic correspondence
\[
 Z_X=\bigl(\pi_1^{C_1}\boxtimes\pi_1^{C_2}\bigr)
       \circ{}^t\Gamma_g\circ e
 \in\CH^2(X_b\times C_1\times C_2)_\Q,
\]
where the curve projectors are those constructed there. Its action on
$U_{X_b}$ is the injection $g^*$, with image entirely in
$H^1(C_1,\Q)\otimes H^1(C_2,\Q)$. Set
\[
 Z=Z_X\circ\Theta_b.
\]
Its action on $U_b$ is $f=g^*\varphi_b$ and is injective. The
composition passes through the surface $X_b$, so its codimension is
$2+2-2=2$. The equality $U_b=T_{S_b^0}$ and the full Hodge group
follow from Propositions~\ref{geom:period} and~\ref{pencil:realization}
at these selected periods.
\end{proof}

The comparison just constructed fits into
\[
 (U_b=T_{S_b^0},q_Y)
 \xrightarrow{\ \varphi_b=(\Theta_b)_*\ }
 (U_{X_b},q_X)
 \xrightarrow{\ g^*\ }
 H^1(C_1,\Q)\otimes H^1(C_2,\Q).
\]
Here $q_X$ is the cup-product form on $X_b$, and $q_Y$ is the form
inherited from $Y$ in Proposition~\ref{geom:structure}; the latter is
exactly the K3 cup-product form by Proposition~\ref{pencil:realization}.
The resolved pullback satisfies $\varphi_b^*q_X=2q_Y$.  Thus the first
map is an algebraic Hodge identification with a factor-two change of
form, and the second is the algebraic injection supplied by the curves.
The rational isometry used below concerns $q_Y$, not the rescaled
form from $X_b$.

\subsection{The Clifford map and its tensor form}

We recall the tensor of Section~\ref{intro:section} for an
eight-dimensional quadratic space $(U,q)$.  With the Clifford convention
$v^2=q(v,v)$, the full even module is
\[
 W=C^+(U,q),\qquad \dim_\Q W=2^7=128.
\]
If $e_1,e_2$ form an oriented orthonormal basis of the positive period
plane, with $U^{2,0}=\C(e_1+i e_2)$, put $J=e_1e_2$.  Then $J^2=-1$,
and the Kuga--Satake Hodge circle acts on $W_\R$ by left
multiplication by $\exp(\theta J)$.  Conjugation by the same element
acts on $U^{2,0}$ by $\exp(2i\theta)$, as required by the weight-two
Hodge structure.

Fix the prescribed rational vector $w\in U$ with $q(w,w)\ne0$.  The
standard embedding on the even module is
\begin{equation}\label{final:specified-clifford}
 j_w:U(1)\longrightarrow\End(W),\qquad
 j_w(v)(a)=vaw\quad(a\in C^+(U,q)).
\end{equation}
The output is even, and $w^{-1}=w/q(w,w)$ shows, by evaluation at
$a=1$, that $j_w$ is injective.  For $g\in\Spin(U,q)$ and left
multiplication $L_g$ one has
\[
 L_gj_w(v)L_g^{-1}(a)=(gvg^{-1})aw=j_w(gvg^{-1})(a).
\]
Thus $j_w$ is a rational Hodge map for every period; the right factor
$w$ is fixed and need not have any particular Hodge type.

Let $E$ be the prescribed rational polarization of $W=H^1(A,\Q)$.
With $\End(W)=W\otimes W^\vee$, its identification
$I_E=(E^\flat(1))^{-1}:W^\vee\xrightarrow{\sim}W(1)$ converts $j_w$ into
\begin{equation}\label{final:specified-tensor}
 \iota_{w,E}:U\longrightarrow W\otimes W
                  \subset H^2(A\times A,\Q).
\end{equation}
Explicitly, $(\id\otimes I_E)j_w$ maps $U(1)$ into
$(W\otimes W)(1)$, and twisting by $\Q(-1)$ gives
Equation~\eqref{final:specified-tensor}.  This is the specified untwisted
tensor.  Its recovery will use the full rational space of Hodge maps
to $W\otimes W$.

\subsection{A rational spanning lemma for the two half-spin factors}

The Clifford tensor and the even-module half-spin decomposition are
classical; see \cite[Proposition~6.3 and Example~6.6]{VanGeemenKS}.
The spanning argument below also uses rational Hodge-Hom descent and
both tensor orientations.
The following statement is about effective weight-one Hodge
structures: their only types are $(1,0)$ and $(0,1)$.

\begin{lemma}\label{final:spin}
Let $(U,q)$ be an eight-dimensional polarized rational Hodge structure
of K3 type, of signature $(2,6)$, and suppose
$\Hg(U)=\SO(U,q)$.  Let $W_1,W_2$ be polarized effective weight-one
Hodge structures and let $W=C^+(U,q)$ have its Kuga--Satake Hodge
structure.  For every nonzero Hodge map
\[
 f:U\longrightarrow W_1\otimes W_2,
\]
put
\[
 \mathcal H=\Hom_{\mathrm{HS}}(W_1,W)\otimes_{\Q}
       \Hom_{\mathrm{HS}}(W_2,W).
\]
If $\tau:W\otimes W\to W\otimes W$ is the ordinary tensor flip, the
rational linear map
\[
 \Phi:\mathcal H\oplus\mathcal H\longrightarrow
       \Hom_{\mathrm{HS}}(U,W\otimes W),
\]
defined by
\[
 \begin{cases}
 \Phi(a\otimes b,0)=(a\otimes b)\circ f,\\
 \Phi(0,a\otimes b)=\tau\circ(a\otimes b)\circ f
 \end{cases}
\]
is surjective.

If $U\subset H^2(M,\Q)$ for a smooth projective surface $M$,
$W_i=H^1(C_i,\Q)$, $W=H^1(A,\Q)$, and $f$ is the restriction of a
codimension-two algebraic correspondence from $M$ to $C_1\times C_2$,
then every map in the target of $\Phi$ is induced by a cycle in
$\CH^2(M\times A\times A)_\Q$.
\end{lemma}

\begin{proof}
We prove surjectivity of the rational map $\Phi$ by spanning both
complex half-spin tensor orientations.  We first isolate
the rational orthogonal Lie algebra in the common Hodge group.
Weight-one effectivity will restrict the nontrivial constituents of
$W_1,W_2$ to half-spins and identify their rational intertwiners with
Hodge maps to $W$.  This allows the tensor-product
calculation to be made over $\C$ and its surjectivity to be descended
to $\Q$.

\emph{The rational Hodge group.}
Let
\[
 G=\Hg(W_1\oplus W_2\oplus W).
\]
The Hodge group is the connected rational algebraic group generated
by the Hodge circle.  Any rational invariant subspace of the
polarized pure weight-one structure $W_1\oplus W_2\oplus W$ has a
Hodge complement, obtained from its polarization.  The complement
is again $G$-invariant.  Its faithful representation is therefore
completely reducible, so $G$ is reductive.

The embedding in Equation~\eqref{final:specified-clifford} realizes
$U(1)$ in $\End(W)$; Tate twists do not change the Hodge-circle
action.  This gives a representation of $G$ on $U$ preserving $q$.
Since $G$ is connected, its image lies in $\SO(U,q)$.  It contains
the Hodge circle of $U$ and therefore $\Hg(U)=\SO(U,q)$, so the map
is surjective.  Write $\mathfrak g=\operatorname{Lie}(G)$.
The derived Lie algebra of $\mathfrak g$ is a direct sum of rational
simple ideals.  Since $\mathfrak{so}(U,q)$ is absolutely simple,
exactly one of those ideals maps nontrivially to it: two such images
would be commuting copies of the same nonabelian simple algebra.
The nonzero map is an isomorphism, and the remaining ideals and the
center form its kernel.  We obtain rational Lie ideals
\begin{equation}\label{final:lie-splitting}
 \mathfrak g=\mathfrak s\oplus\mathfrak k,
 \qquad
 \mathfrak s\xrightarrow{\sim}\mathfrak{so}(U,q),
 \qquad
 \mathfrak k=\ker\bigl(\mathfrak g\to\mathfrak{so}(U,q)\bigr).
\end{equation}
The ideal $\mathfrak s$ integrates to a connected rational normal
semisimple subgroup.  Finite central kernels are harmless here; we do
not identify a spin representation with a representation of
$\SO(U,q)$.

\emph{The allowed weight-one constituents.}
Normalize the complexified infinitesimal Hodge-circle generator $D$
so that its eigenvalues on weight one are $1,-1$.  Its image on $U$ has
eigenvalues $2,0,\ldots,0,-2$.  Write
$D=D_{\mathfrak s}+D_{\mathfrak k}$ according to
Equation~\eqref{final:lie-splitting}.  Each $\mathfrak s_\C$-isotypic
summand of $W_{i,\C}$ has the form $R_\lambda\otimes M$, with
$R_\lambda$ irreducible.  The commuting operator
$D_{\mathfrak k}$ acts only on $M$.  The two projections of the
semisimple element $D$ are semisimple, so their eigenvalues add on
simultaneous eigenspaces.

Identify the root system of $\mathfrak s_{\C}$ with $D_4$, with
roots $\pm e_i\pm e_j$, using a Cartan subalgebra in which
$D_{\mathfrak s}$ pairs with weights as $2e_1$.
A dominant highest weight has
\[
 \lambda_1\geq\lambda_2\geq\lambda_3\geq|\lambda_4|,
\]
with coordinates either all integral or all half-integral.  Its Weyl
orbit contains weights giving the eigenvalues
$2\lambda_1$ and $-2\lambda_1$ for $D_{\mathfrak s}$.  Translation by
any eigenvalue of the commuting factor cannot reduce this spread.
Effectivity of weight one therefore gives $4\lambda_1\leq2$.  The
only possibilities are the trivial representation and the two
half-spin representations $S_+,S_-$, with highest weights
$\tfrac12(e_1+e_2+e_3+e_4)$ and
$\tfrac12(e_1+e_2+e_3-e_4)$.  Their weights are all sign choices
$\tfrac12(\pm e_1\pm e_2\pm e_3\pm e_4)$ of the respective
parities.  Each half-spin therefore has both
$D_{\mathfrak s}$-eigenvalues $1,-1$.  If $\mu$ is an eigenvalue of
$D_{\mathfrak k}$ on its multiplicity space, both $\mu+1$ and
$\mu-1$ must belong to $\{1,-1\}$; hence $\mu=0$.
Semisimplicity now shows that $D_{\mathfrak k}$ is zero on that
multiplicity space.

The rational decompositions
\[
 W_i=W_i^{\mathfrak s}\oplus W_i^{\mathrm{nt}},
 \qquad W_i^{\mathrm{nt}}=\mathfrak s W_i
\]
are $G$-stable.  On $W_i^{\mathrm{nt}}$ the entire Hodge-circle
generator consequently acts through $\mathfrak s$.
The same assertion holds on $W$.  Its Hodge group lies in the
rational spin group acting by left multiplication on $C^+(U,q)$
and maps onto $\SO(U,q)$.  The spin covering has finite kernel,
so the Lie algebra of this image on $W$ maps isomorphically to
$\mathfrak{so}(U,q)$; in particular $\mathfrak k$ acts trivially
on $W$.

For completeness, write $U_\C=F\oplus F^\vee$ with $F$ maximal
isotropic of dimension four.  Exterior multiplication and contraction
give the Clifford action on $\bigwedge^\bullet F$.  These operators
generate its matrix units, and the even algebra preserves parity.
Thus, with $S_+=\bigwedge^{\mathrm{even}}F$ and
$S_-=\bigwedge^{\mathrm{odd}}F$, one has
\[
 C^+(U)_{\C}\simeq\End(S_+)\oplus\End(S_-).
\]
Since $\dim S_+=\dim S_-=8$, left multiplication gives eight
copies of each half-spin.
In particular $W$ has no $\mathfrak s$-invariant summand.

\emph{Rational Hodge maps and complex multiplicities.}
A rational $\mathfrak s$-intertwiner from $W_i^{\mathrm{nt}}$ to $W$
commutes with $D$ and hence with the connected Hodge circle.  Extended
by zero on $W_i^{\mathfrak s}$, it is therefore a Hodge map.
Conversely every rational Hodge map commutes with $G$ and hence with
$\mathfrak s$.  We have proved
\begin{equation}\label{final:hom-equality}
 \begin{aligned}
 \Hom_{\mathrm{HS}}(W_i,W)
   &=\Hom_{\mathfrak s}(W_i^{\mathrm{nt}},W),\\
 \Hom_{\mathrm{HS}}(U,W\otimes W)
   &=\Hom_{\mathfrak s}(U,W\otimes W).
 \end{aligned}
\end{equation}
In the second equality, both Hodge structures have weight two and
their Hodge-circle actions factor through $\mathfrak s$, so the
same commutation argument applies.  For any two rational
$\mathfrak s$-modules $V_1,V_2$ the intertwiner space is the kernel
of the rational linear map
\[
 \Hom_\Q(V_1,V_2)\longrightarrow
 \Hom_\Q(\mathfrak s\otimes V_1,V_2),\qquad
 a\longmapsto\bigl((x,v)\mapsto a(xv)-x\,a(v)\bigr).
\]
Extension of scalars is exact.  Hence
$\Hom_{\mathfrak s}(V_1,V_2)\otimes_\Q\C
=\Hom_{\mathfrak s_\C}(V_{1,\C},V_{2,\C})$.
In particular, after complexifying the rational Hodge Hom spaces in
Equation~\eqref{final:hom-equality}, all complex maps between the
half-spin multiplicity spaces are available.  This assertion does
not require either half-spin type to be defined separately over
$\Q$.

\emph{Spanning both tensor orientations.}
Write
\[
 \begin{aligned}
 W_{i,\C}^{\mathrm{nt}}
   &=(S_+\otimes M_{i,+})\oplus(S_-\otimes M_{i,-}),\\
 W_{\C}&=(S_+\otimes N_+)\oplus(S_-\otimes N_-),
       \qquad \dim N_+=\dim N_-=8.
 \end{aligned}
\]
Both half-spins are self-dual: the Weyl group contains simultaneous
negation of all four coordinates, preserving the two weight
parities.  Clifford antisymmetrization identifies
$C(U)_\C$ with $\bigwedge^\bullet U_\C$ as a module for conjugation.
The even part consists of the diagonal endomorphism blocks of
$S_+\oplus S_-$, and the odd part of the two off-diagonal blocks.
Choose a normalized volume element $\omega$ with $\omega^2=1$.
Left multiplication by $\omega$ exchanges exterior degrees $k$ and
$8-k$ and distinguishes the two blocks of each parity.  Each block
therefore contains one copy of each paired degree; degree four
splits into its two $35$-dimensional eigenspaces.  This gives
\begin{equation}\label{final:spin-products}
 \begin{aligned}
 S_+\otimes S_-&\simeq U_{\C}\oplus\bigwedge^3U_{\C},\\
 S_\pm\otimes S_\pm&\simeq
 \C\oplus\bigwedge^2U_{\C}\oplus\bigwedge^4_\pm U_{\C}.
 \end{aligned}
\end{equation}
The choice of the signs on $\bigwedge^4_\pm U_\C$ follows that of
$S_\pm$.  To check the needed multiplicity, $\bigwedge^3U_\C$ has
no standard summand.  Indeed, an equivariant map from $U_\C$ to
this exterior power would send a unit vector $u$ to an invariant
of its stabilizer $\SO(u^\perp)\simeq\SO_7$.  Neither
$\bigwedge^3u^\perp$ nor $u\wedge\bigwedge^2u^\perp$ has such an
invariant: determinant-one sign changes in an orthonormal basis
kill every exterior basis coefficient.  Thus the map vanishes on
the orbit of $u$, which spans $U_\C$.  The standard summand in each
opposite-sign product has multiplicity one.  There is none in a
same-sign product, since $-\id\in\SO(U_\C)$ acts trivially on all
even exterior powers and by $-1$ on $U_\C$.

Choosing a nonzero equivariant map from $U_\C$ into each opposite-sign
product identifies the complex target space with
\[
 \Hom_{\mathfrak s_\C}(U_\C,W_\C\otimes W_\C)
 \simeq (N_+\otimes N_-)\oplus(N_-\otimes N_+).
\]
The ordinary flip interchanges these two summands.  In particular,
Equation~\eqref{final:hom-equality} and base change give
\[
 \dim_\Q\Hom_{\mathrm{HS}}(U,W\otimes W)=8^2+8^2=128.
\]
It remains to generate both multiplicity spaces from the one map $f$.

A tensor involving a trivial $\mathfrak s$-constituent of either
$W_i$ is a sum of trivial or half-spin representations, so it also
contains no standard summand.  Therefore $f_{\C}$ is described by two
multiplicity tensors
\[
 t_{+-}\in M_{1,+}\otimes M_{2,-},
 \qquad
 t_{-+}\in M_{1,-}\otimes M_{2,+},
\]
at least one of which is nonzero.

Suppose, for example, that $t_{+-}\ne0$.  Choose linear functionals
$\ell_1,\ell_2$ with
$(\ell_1\otimes\ell_2)(t_{+-})\ne0$.  For arbitrary
$n_+\in N_+$ and $n_-\in N_-$, the rank-one maps
$m\mapsto\ell_1(m)n_+$ and $m\mapsto\ell_2(m)n_-$ send this component
to a nonzero scalar multiple of $n_+\otimes n_-$.  Taking the maps
zero on the other signs isolates the component.  Such decomposable
tensors span $N_+\otimes N_-$, and the independent second copy of $\mathcal H$
in $\Phi$ supplies the other orientation by $\tau$.  The case
$t_{-+}\ne0$ is identical.  Hence $\Phi_{\C}$ is surjective.  Its
rational cokernel has zero complexification and is therefore zero,
proving surjectivity of $\Phi$.  This descent accounts for arbitrary
rational forms and Schur multiplicities; neither half-spin type need
be defined separately over $\Q$.

\emph{Algebraic realization.}
Let $a:H^1(C_i,\Q)\to H^1(A,\Q)$ be a rational Hodge map.
Poincar\'e duality on $C_i$ identifies it with a unique rational
class $\alpha_a$ in the middle K\"unneth summand
$H^1(C_i,\Q)\otimes H^1(A,\Q)\subset H^2(C_i\times A,\Q)$,
normalized by
\[
 (p_A)_*(\alpha_a\smile p_{C_i}^*x)=a(x).
\]
The Hodge condition says precisely that $\alpha_a$ has type $(1,1)$.
After clearing denominators, the exponential sequence realizes it
as the first Chern class of a holomorphic line bundle: the
obstruction is its zero image in $H^2(\mathcal O_{C_i\times A})$.
Projective GAGA makes this line bundle algebraic.  Dividing by the
same denominator gives a divisor
$D_a\in\CH^1(C_i\times A)_\Q$ with $\cl(D_a)=\alpha_a$.
Thus its cohomology class has exactly the required middle component.

For maps $a:H^1(C_1)\to W$ and $b:H^1(C_2)\to W$, the exterior
product $D_a\boxtimes D_b$, with its factors ordered as
$(C_1\times C_2)\times(A\times A)$, has codimension two and acts
on $H^1(C_1)\otimes H^1(C_2)$ by $a\otimes b$.
If $Z$ is the given cycle from a surface, its composition with this
product has codimension $2+2-2=2$ and has image in $W\otimes W$.

Let $s:A\times A\to A\times A$ interchange the factors.  For
$u,v\in H^1(A,\Q)$,
\[
 s_*\bigl(p_1^*u\smile p_2^*v\bigr)
 =p_2^*u\smile p_1^*v
 =-p_1^*v\smile p_2^*u.
\]
Consequently $-\Gamma_s$ induces the ordinary tensor flip $\tau$
on this summand.  Given a target map $h$, choose a rational
preimage under $\Phi$ and write it as a finite sum
\[
 h=\sum_\nu c_\nu(a_\nu\otimes b_\nu)f
   +\sum_\mu d_\mu\tau(a'_\mu\otimes b'_\mu)f.
\]
It is realized exactly by the rational cycle
\[
 \sum_\nu c_\nu(D_{a_\nu}\boxtimes D_{b_\nu})\circ Z
 -\sum_\mu d_\mu\Gamma_s\circ
       (D_{a'_\mu}\boxtimes D_{b'_\mu})\circ Z.
\]
Composition with $\Gamma_s$, whose codimension is $\dim(A\times A)$,
also preserves codimension two.  On the specified source subspace
the resulting action equals $h$ in the full $H^2(A\times A,\Q)$.
\end{proof}

\subsection{Transfer by a rational Hodge isometry}

\begin{proposition}\label{final:generic}
Let $b$ be a construction point of Lemma~\ref{final:curves}, and let
$S$ be a projective K3 surface equipped with a rational Hodge isometry
$r_T:T_S\xrightarrow{\sim}U_b=T_{S_b^0}$ for the cup-product forms.
For every full even-Clifford model $A$ of $T_S$, every rational
nonisotropic right factor $w$, and every rational polarization $E$,
there is a cycle
\[
 \Gamma_S\in\CH^2(S\times A\times A)_\Q,
 \qquad (\Gamma_S)_*x=\iota_{w,E}(x)\quad(x\in T_S),
\]
where the equality holds in $H^2(A\times A,\Q)$.
\end{proposition}

\begin{proof}
The assumed $r_T$ uses the form on $U_b$ inherited from $Y$, which
the auxiliary K3 realizes isometrically.  The factor-two pullback to
$X_b$ enters only the curve correspondence of Lemma~\ref{final:curves}
and is not part of this isometry.

Choose markings of the full second cohomology spaces by
$\Lambda_{\mathrm{K3},\Q}$.  They turn $r_T$ into an isometry
between two nondegenerate rational subspaces of this one quadratic
space.  Witt's Extension Theorem
\cite[Corollary~7.2.8]{PetersSterk} extends it to an
isometry of $\Lambda_{\mathrm{K3},\Q}$.  Returning through the
markings gives
\[
 r:H^2(S,\Q)\xrightarrow{\sim}H^2(S_b^0,\Q).
\]
The isometry $r$ maps the orthogonal complement of $T_S$ to
the orthogonal complement of $T_{S_b^0}$.  These are exactly
$\NS(S)_\Q$ and $\NS(S_b^0)_\Q$, both of type $(1,1)$.
Since its restriction $r_T$ is Hodge, $r$ is therefore a rational
Hodge isometry on the full second cohomology.  Buskin's
Theorem~\cite[Theorem~1.1]{Buskin}, applied to these smooth projective
complex K3 surfaces, gives a cycle
\[
 R\in\CH^2(S\times S_b^0)_\Q,\qquad
 R_*|_{H^2(S,\Q)}=r.
\]
The hypotheses here are precisely full rational Hodge isometry and
projectivity; the integral N\'eron--Severi lattice of $S_b^0$ is not
required to equal $P$.

Compose $R$ with the cycle $Z$ of Lemma~\ref{final:curves}.
The resulting codimension-two correspondence induces the injection
\[
 f_S:T_S\xrightarrow{r_T}U_b
       \xrightarrow{f}H^1(C_1,\Q)\otimes H^1(C_2,\Q).
\]
At these periods $\Hg(T_S)=\SO(T_S)$.  Apply
Lemma~\ref{final:spin} to $f_S$ and to the given full
even-Clifford model $W=H^1(A,\Q)$ of dimension $128$.
The prescribed map $\iota_{w,E}$ of
Equation~\eqref{final:specified-tensor} belongs to the target of
$\Phi$.  A rational preimage under $\Phi$, realized by the divisor
products and signed interchange in that lemma, supplies $\Gamma_S$
with exactly the asserted action.  In particular the equality is
on all of $T_S$ and has no extra component in the other
K\"unneth summands of $H^2(A\times A,\Q)$.

The choices of abelian data, pencil, and final parameter retain the
successive order of Lemma~\ref{geom:genericity} and
Corollary~\ref{comp:surface}.  We have used the full Hodge group only
at these construction periods.  Proposition~\ref{specialization:ambient}
will realize one fixed standard tensor throughout a target period
ball, including Picard jumps.
\end{proof}

\begin{remark}
Buskin's theorem does not require integrality or preservation of an
ample chamber.  In its cyclic reduction, order one means an integral
isometry.  After a possible sign change and finitely many reflections
in algebraic $(-2)$-classes, it maps a K\"ahler chamber to a K\"ahler
chamber, so strong Torelli realizes it by an isomorphism; its graph
and the sign and reflection correspondences are algebraic.  In the
final rational-reflection chain, transport an initial integral ample
class through the partial isometries.  On each intermediate
period-realized K3 it is a rational $(1,1)$-class of positive square.
Clearing denominators gives an integral such class, so the K3
projectivity criterion makes every intermediate surface projective,
as required for the cyclic steps.
\end{remark}
\section{Algebraic cycles at every period}
\label{specialization:section}

Proposition~\ref{final:generic} constructs the prescribed correspondence
at the selected construction periods of Lemma~\ref{final:curves}.
To reach a Picard-jump target, we keep the perpendicular of $P$ fixed,
rather than varying the source with the transcendental space.

\begin{proposition}\label{specialization:ambient}
Let $S$ have an ample primitive $P$-polarization and put
$U_P=P_{\Q}^{\perp}\subset H^2(S,\Q)$.  There are a full even-Clifford
realization $W_U=C^+(U_P,q)=H^1(A_U,\Q)$, a rational polarization
$E_0$, a nonisotropic rational $w_0\in U_P$, and a cycle
\[
 \Gamma_U\in\CH^2(S\times A_U^2)_{\Q}
\]
whose full degree-two action on every $x\in U_P$ is exactly
$\iota_{w_0,E_0}(x)$.  This assertion concerns the whole fixed
rank-eight space, including its rational $(1,1)$ classes at a Picard
jump.
\end{proposition}

The parameter-space lemma below records equality of a cycle's action
on all of this fixed source.  Its proper images will supply actual
cycles at every target, without specializing the curves used to
construct a seed.

\subsection{Proper parameter spaces of successful cycles}

\begin{lemma}\label{final:hilbert}
Let $f:\mathcal Z\to B$ be a smooth proper holomorphic family of pure
relative dimension $n$ over a connected coordinate ball, with a closed
relative embedding $\mathcal Z\subset\PP^N\times B$.  Fix an integer
$0\leq r\leq n$, a rational local system $\mathcal V$ on $B$, a
morphism of local systems
\[
 L:R^{2r}f_*\Q\longrightarrow\mathcal V,
\]
and a flat section $\tau$ of $\mathcal V$.  The set of $b\in B$ admitting
a rational codimension-$r$ algebraic cycle $Z_b$ such that
\begin{equation}\label{final:cycle-condition}
 L_b\bigl(\cl(Z_b)\bigr)=\tau_b
\end{equation}
is a countable union of closed analytic subsets of $B$.  Each subset
in this union is the image of a proper analytic parameter space
carrying a fixed rational combination of universal flat subschemes;
every point of that space supplies a cycle satisfying
\eqref{final:cycle-condition}.

If the successful set contains a nonempty open subset of $B$, it equals
$B$.  The same conclusion holds if, for every countable union of proper
closed analytic subsets of $B$, there is a successful point outside
that union.
\end{lemma}

\begin{proof}
\emph{Embedded parameters over an analytic base.}
Let $R$ be the Hilbert polynomial of the embedded fibres of $f$; it is
constant because the family is flat.  For any fixed Hilbert polynomial
$Q$, write $K_Q=\operatorname{Hilb}^{Q}(\PP^N)$.  This is a projective
scheme with a flat universal subscheme, by the projective Hilbert-scheme
theorem \cite[Theorem~3.2]{GrothendieckHilbert}.  We use the incidence
space
\[
 H_Q=\{(b,z)\in B\times K_Q^{\mathrm{an}}:
                    \mathcal U_z\subset\mathcal Z_b\}.
\]
It is closed analytic, with its natural incidence structure.  To see
this directly, let $\mathcal J$ be the ideal sheaf of $\mathcal Z$ in
$\PP^N\times B$.  The bounded regularity for the fixed polynomials
$R$ and $Q$ supplies an integer $k$ such that $\mathcal J(k)$ is
relatively generated, its direct image is locally free with base
change, and the universal subscheme for $Q$ has the corresponding
vanishing and base-change properties.  On $B\times K_Q^{\mathrm{an}}$,
restriction defines a holomorphic vector-bundle map from sections of
$\mathcal J(k)$ to sections of $\mathcal O_{\mathcal U}(k)$.  Its zero
scheme is exactly $H_Q$: vanishing of these generating sections is
equivalent to containment of subschemes.  Since $K_Q^{\mathrm{an}}$ is
compact, $H_Q\to B$ is proper.  Its universal subscheme is flat over
$H_Q$, being the pullback of the universal Hilbert family.  No
algebraicity assumption on $B$ or on $f$ is required.

\emph{The ambient Betti class on singular parameters.}
For polynomials of degree $n-r$, associate to each universal fibre its
$(n-r)$-dimensional fundamental cycle, with the generic lengths of its
structure sheaf as multiplicities.  We claim that its class in the
ambient $H^{2r}(\mathcal Z_b,\Q)$ is locally constant on $H_Q$, including
at singular or nonreduced Hilbert parameters.

Here are details that justify this assertion in the analytic setting.
The same high-degree ideal construction gives a holomorphic
classifying map $B\to\operatorname{Hilb}^{R}(\PP^N)^{\mathrm{an}}$:
the bundle of degree-$k$ ideals defines the map to the Hilbert
scheme's Grassmannian realization.  Its image is in the open locus
$K_R^{\mathrm{sm}}$ of smooth, pure $n$-dimensional fibres.  Over this
algebraic locus let $\mathfrak Z\to K_R^{\mathrm{sm}}$ be the universal
smooth family, and let $\mathfrak H_Q$ be its relative Hilbert scheme.
Thus $H_Q$ and its universal pair are analytic pullbacks of this
algebraic universal pair.

On $\mathfrak Z\times_{K_R^{\mathrm{sm}}}\mathfrak H_Q$, the universal
structure sheaf $\mathcal F$ is flat over $\mathfrak H_Q$ and is a
perfect complex on the ambient space.  For this last assertion, take
successive locally free surjections onto $\mathcal F$ and their kernels
on the quasi-projective ambient space.  All the kernels remain flat
over $\mathfrak H_Q$.  The fibres of the ambient family are regular of
dimension $n$, so after at most $n$ syzygies the fibrewise kernels are
locally free.  The local flatness criterion then makes this kernel
locally free on the ambient space.  This gives a bounded vector-bundle
resolution even when $\mathfrak H_Q$ is singular.  Flatness over the
parameter space also shows that the resolution restricts to a
resolution of the ordinary structure sheaf on every ambient fibre.

The topological Chern character of this complex, pulled back to the
analytic universal pair, consequently satisfies
\begin{equation}\label{specialization:leading-character}
 \operatorname{ch}_r(\mathcal F)|_{\mathcal Z_b}
       =\cl\bigl([\mathcal U_z]_{n-r}\bigr)
       \quad\text{in }H^{2r}(\mathcal Z_b,\Q).
\end{equation}
This is the leading Chern-character identity for a coherent sheaf
supported in codimension at least $r$.  It follows from
Grothendieck--Riemann--Roch (GRR) and d\'evissage, followed by the
comparison of algebraic and topological Chern classes
\cite[Theorem~15.2, Example~15.2.16\textup{(a)}, and
Proposition~19.1.2]{Fulton}.
For the precise form needed here, remove a subset of codimension at least
$r+1$: d\'evissage then reduces to the generic lengths on smooth
codimension-$r$ supports, where the regular-embedding case of GRR gives
their fundamental classes as the leading terms.
The removed subset contributes nothing in degree $2r$, since its
cohomology with supports vanishes below degree $2r+2$
\cite[Section~19.1, equation~(3) and Lemma~19.1.1]{Fulton}.  This also
explains why lower-dimensional embedded components cause no problem.

A class on the total space of a pulled-back smooth proper family
restricts to a flat section of its fibre cohomology.  Indeed, locally
trivialize the original differentiable family over $B$ and pull that
trivialization back to the parameter space; inclusions of fibres
along a path are homotopic.  Complex analytic spaces are locally path
connected, and their nilpotents do not change their topology.
Equation~\eqref{specialization:leading-character} proves the claimed
constancy on all connected components, without assuming that the
universal subschemes themselves are smooth.

\emph{Successful components and their images.}
Take every finite list $Q_1,\ldots,Q_s$ of polynomials of degree $n-r$
and every list $c_1,\ldots,c_s\in\Q$, and form
\[
 H_{\boldsymbol Q}=H_{Q_1}\times_B\cdots\times_B H_{Q_s},
 \qquad \mathcal Z_{\boldsymbol c,z}
       =\sum_{i=1}^s c_i[\mathcal U_{i,z}]_{n-r}.
\]
The empty list represents the zero cycle.  These spaces are proper
over $B$, and the displayed universal combinations have constant
ambient Betti class on each connected component.  Because $B$ is
simply connected, both $L$ and $\tau$ are constant in local-system
trivializations.  Hence \eqref{final:cycle-condition} holds either
everywhere or nowhere on each component.

Retain exactly the successful components.  There are countably many
polynomial and coefficient lists and countably many components:
all the parameter spaces are second countable and locally connected,
so their connected components are open and closed.  Each retained
component is a closed analytic subspace, proper over $B$.  Its image
is closed analytic by Remmert's proper mapping theorem; for a primary
statement, see \cite[\S7.1, Satz~1]{GrauertProperImage}, applied to the
reduction of the component.  Reduction changes neither its points
nor its image.  For a textbook statement of the proper mapping theorem,
see \cite[Chapter~II, Section~8.2, Theorem~8.8]{Demailly2012}.

Every rational algebraic cycle on a fibre is a finite rational sum
of integral subvarieties and therefore occurs in this list.
Conversely, every point of any retained image has an actual universal
fibre combination above it.  Their supports in the projective fibre
are algebraic by Chow's theorem, so their fundamental cycles are
algebraic.  We have proved that the successful
set is exactly the stated countable union, with the asserted
cycle-theoretic meaning at every image point.

A proper closed analytic subset of a connected ball is nowhere dense.
If the union contains a nonempty open set, Baire's theorem forces
one image to have interior; the analytic identity principle makes
that image all of $B$.  If instead the union escapes every countable
union of proper closed analytic subsets, the same conclusion follows:
otherwise the successful images would themselves be such a union.
\end{proof}

\subsection{Realization in an ample period ball}

We now apply the proper-parameter lemma directly in a small period ball
through the desired target.  A rational change of period coordinates
places the generic geometric construction in that ball; it does not
change the algebraic pencil used in the construction.

\begin{lemma}\label{specialization:rational-density}
Let $(V,q)$ be a rational quadratic space of signature $(2,n)$.
The group $O(V)(\Q)$ is dense in $O(V_{\R})$, and the rational orbit
of every point of its type-IV period domain meets every nonempty open
subset of that domain.
\end{lemma}
\begin{proof}
Cartan--Dieudonn\'e expresses each real orthogonal transformation as
a finite product of reflections in nonisotropic real vectors.  The
reflection
\[
 r_v(x)=x-\frac{2q(x,v)}{q(v,v)}v
\]
depends continuously on $v$ on the nonisotropic open and is rational
for rational $v$.  Approximating the finitely many vectors proves
density.  The full real orthogonal group acts transitively on the
oriented positive two-planes, which are the period-domain points;
it also interchanges the two components.  Density proves the orbit
assertion.  We use the standard reflection theorem in
\cite[Theorem~7.2.4]{PetersSterk}.
\end{proof}

\begin{proof}[Proof of Proposition~\ref{specialization:ambient}]
\emph{The fixed local family and tensor.}
Choose a marking of the target and an integral ample class $h\in P$.
The local K3 period theorem gives a six-dimensional marked deformation
ball $B$ through $S$ on which $P$ remains algebraic, identified with
an open subset of the period domain of $U_P$.  The class $h$ lifts
locally to a line bundle; shrink $B$ so it is relatively ample and
a power embeds the family $\mathcal S\to B$ projectively.  Only $P$,
not the target's whole Picard group, is kept algebraic.  Thus $U_P$
is a fixed nondegenerate rational Hodge subvariation even at jumps.
These local assertions follow from local Torelli and the deformation
theory of polarized K3 surfaces \cite{Huybrechts}.

Fix any rational $w_0\in U_P$ with $q(w_0,w_0)\ne0$ and put
$W_U=C^+(U_P,q)$.  There is one rational alternating form polarizing
the left Clifford Hodge structures on the chosen period component.
Here is a construction that also fixes its characters.  Let $x^*$ be
reversal and $B_C(x,y)=\operatorname{scal}(x^*y)$.  Choose two
orthogonal positive rational vectors with oriented product $\beta$
in that component and set, up to an overall sign,
\[
 E_0(x,y)=B_C(x,y\beta).
\]
Reversal sends $\beta$ to $-\beta$, and cyclic invariance of scalar
part proves alternation.  Left spin multiplication preserves $E_0$
because $g^*g=1$.  At the chosen plane write $\beta=cJ$, $c>0$.
For an even orthogonal monomial containing $a$ of the two positive
generators, $B_C(x,x)$ has sign $(-1)^a$ and
$Jx\beta=(-1)^{a+1}cx$.  Thus $E_0(x,Jx)$ is definite after one
overall choice of sign.  Transitivity of the connected real spin
group on the component proves positivity throughout it.  This is the
standard Kuga--Satake polarization \cite{KS}.

The $(1,0)$ space of $W_U$ is $\ker L_\sigma$ for a period vector
$\sigma$, and varies holomorphically.  A rational symplectic basis for
$E_0$ and its self-dual integral lattice give, over a smaller $B$, a
polarized holomorphic abelian family $\mathcal A\to B$ with this fixed
rational $H^1$.  The usual torus construction over Siegel space and a
relative theta bundle give a relative projective embedding after
shrinking.  The tensor $\tau=\iota_{w_0,E_0}$ is a fixed rational
linear map $U_P\to W_U\otimes W_U$ and a Hodge map at every period,
by spin equivariance of $j_{w_0}$.

\emph{A geometric period chart in the target ball.}
Choose a geometric datum in the simultaneous regular-period and
domination open of Sections~\ref{geom:section} and
\ref{domination:section}.  On a local marking fix a rational isometry
$\theta:U\simeq U_P$ for the form inherited from the quotient surface
$Y$; this is possible by Proposition~\ref{geom:structure} and
Lemma~\ref{main:lattice}.  Lemma~\ref{specialization:rational-density}
gives a rational $g\in O(U_P)(\Q)$ moving its period into $B$.
Fix $g$ and $\theta$, and shrink the original normalized geometric
product chart so that its period map followed by $g\theta$ is a
submersion into $B$.
The change $g\theta$ acts only on rational period coordinates.
Pencils subsequently selected still have the form $d_0+td_1$ in the
normalized algebraic section chart, with directions in the original
good open at infinity.

\emph{Proper images for the full action.}
Apply Lemma~\ref{final:hilbert} to
$\mathcal Z=\mathcal S\times_B\mathcal A^2$, of relative dimension
$130$, with $r=2$ and
\[
 \mathcal V=\Hom_{\Q}\bigl(U_P,R^2(\mathcal A^2/B)_*\Q\bigr),
 \qquad
 L_b(\gamma)(x)=(p_{AA})_*(\gamma\smile p_S^*x).
\]
The success condition is equality to $\tau$ on all of $U_P$ in the
full target $H^2(A_b^2)$, not only after a projection or on the
holomorphic line.  It is a flat local-system condition.  The successful
set is therefore a countable union of closed analytic images $Y_i$
of proper cycle-parameter spaces, with actual successful cycles over
every point of each image.

\emph{A seed outside every proposed proper image.}
Suppose every $Y_i$ is proper.  In the regular construction period
chart, pull back the $Y_i$ by the fixed map $g\theta$ and also pull
back the extra-Picard hyperplanes in $B$.  These are countably many
proper closed analytic loci: the period submersion is open, and an
extra integral class has a nonzero $U_P$ projection because
$\Lambda_{\mathrm{K3}}\cap P_{\Q}=P$.  Its period equation is
therefore nontrivial.
Include also the proper marked Hodge-tensor conditions used in
Lemma~\ref{final:curves}.

Apply Lemma~\ref{geom:genericity} to this prescribed list, using the
admissible pencil open of the construction.  It selects abelian data,
then an actual admissible section pencil with good leading section.
Once that pencil, its first models, and its relative correspondence
are fixed, intersect its algebraic good open with that of
Corollary~\ref{comp:surface} and add the corollary's countable
exceptional parameters.  The intersection is still a nonempty
algebraic open of the pencil.  The last choice in the lemma supplies
a point in this open outside those exceptions and every prescribed
period locus.
In particular, the full special orthogonal Hodge group and the curve source
of Lemma~\ref{final:curves} hold at this construction point.

The target fiber at its moved period has $\NS=P$, so its full $U_P$
is transcendental.  The auxiliary K3 has transcendental space $U$.
The rational period isometry $g\theta$ identifies these
transcendental Hodge structures with their cup forms.
Proposition~\ref{final:generic} therefore supplies on the target
fiber a cycle realizing exactly the fixed tensor $\tau$.
That period lies in a successful image, contrary to its choice.
Thus some image is all of $B$, and its fiber over the original
target supplies $\Gamma_U$.

The rank-eight source and its full action have remained fixed
throughout the ball.  Neither a global moduli cycle nor specialization
of a smooth curve presentation is required.
\end{proof}

\section{Exact restriction and rational transfer}
\label{transfer:section}

We first convert the fixed-source result to the standard tensor of
the actual transcendental space at an ample $P$-polarized surface.
A separate rational period transfer will then reach every surface in
the quadratic locus of Theorem~\ref{main:theorem}.

\subsection{Restriction at Picard jumps}

At a Picard jump the inclusion $T_S\subset U_P$ induces an inclusion
of full even-Clifford modules.  The following identity compresses the
ambient tensor to this submodule with its prescribed polarization;
it does not require the ambient polarization to restrict to that choice.

\begin{lemma}\label{transfer:compression}
Let $i:W_T\hookrightarrow W_U$ be an inclusion of polarizable rational
weight-one Hodge structures, let $p:W_U\to W_T$ be a Hodge retraction,
and let $E_U,E_T$ be rational polarizations.  Let $H$ be a rational
Hodge structure of weight zero.  For a Hodge map
$j:H\to\End(W_U)$, write
$\iota_{j,E_U}:H(-1)\to W_U\otimes W_U$ for its tensor form.
Set
\[
 b=(E_T^\flat)^{-1}i^\vee(-1)E_U^\flat:W_U\longrightarrow W_T.
\]
Then
\begin{equation}\label{transfer:compression-identity}
 (p\otimes b)\iota_{j,E_U}=\iota_{pji,E_T}.
\end{equation}
If the weight-one structures are abelian $H^1$ spaces, both maps on
the tensor slots are induced by rational algebraic correspondences.
\end{lemma}
\begin{proof}
In the ordering $\End(W_U)=W_U\otimes W_U^\vee$, tensor conversion
is $\id\otimes(E_U^\flat)^{-1}$ after the common Tate twist.  Hence
the left side of \eqref{transfer:compression-identity} is
\[
 \bigl(p\otimes(E_T^\flat)^{-1}i^\vee(-1)\bigr)j,
\]
which is exactly the tensor conversion of $pji$.  All maps are rational
Hodge maps.  For completeness, if $f:H^1(B,\Q)\to H^1(A,\Q)$ is any
such map, its dual on homology commutes with complex structures.
For some positive integer $n$, $nf^\vee$ is integral and induces a
holomorphic homomorphism $\phi:A\to B$.  Projectivity and GAGA make
$\phi$ algebraic, and $f=\phi^*/n$.  Thus
$[\Gamma_\phi^t]/n\in\CH^{\dim B}(B\times A)_{\Q}$ induces $f$
on $H^1$.  No equality of dimensions or isogeny assumption is needed.
\end{proof}

\begin{proposition}\label{transfer:ample}
Let $S$ be an ample $P$-polarized K3 surface.  For every prescribed
nonisotropic $w\in T_S$, rational polarization $E$ and transported
even-Clifford realization of $C^+(T_S,q)$, its standard tensor
$\iota_{w,E}$ has a rational codimension-two algebraic realization.
\end{proposition}
\begin{proof}
Write $U=U_P$ and $M=U\cap\NS(S)_{\Q}$.  Since $P$ contains an ample
class, Hodge index makes $M$ negative definite and
\[
 U=T_S\perp M.
\]
The source projection to $T_S$ is algebraic: choose divisor bases
for $\NS(S)_{\Q}$ and use their inverse intersection matrix together
with the diagonal.  Proposition~\ref{specialization:ambient} gives
one exact tensor $\iota_{w_0,E_0}$ on all of $U$.

For the desired $w\in T_S$, put $r=w_0^{-1}w\in C^+(U,q)$.
It is invertible.  Right multiplication $R_r$ commutes with the left
Hodge complex structure, so is a rational Hodge automorphism, and
associativity gives
\[
 j_w=R_rj_{w_0},\qquad
 \iota_{w,E_0}=(R_r\otimes\id)\iota_{w_0,E_0}.
\]
This is a change of the first, output tensor slot.

The natural Clifford inclusion
$i:W_T=C^+(T_S,q)\hookrightarrow W_U=C^+(U,q)$ is injective by PBW
and is Hodge: the left period bivector lies in $C^+(T_S,q)$ on both
spaces.  Polarizability gives a rational Hodge retraction $p$.
Apply Lemma~\ref{transfer:compression} with the desired $E$ on $W_T$.
For $v,w\in T_S$, $j_w(v)$ preserves $W_T$ and restricts there to
the standard operator $a\mapsto vaw$.  Therefore
\[
 \bigl(pR_r\otimes(E^\flat)^{-1}i^\vee(-1)E_0^\flat\bigr)
       \iota_{w_0,E_0}
 =\iota^{T_S}_{w,E}.
\]
The order matters: $R_r$ itself need not preserve $W_T$ when
$w_0\notin T_S$, but the resulting operator $j_w$ does.

Both tensor-slot maps are algebraic on abelian $H^1$ by the lemma.
After clearing denominators, the single product homomorphism on the
two abelian factors and its normalized transpose graph realize their
tensor product, with no unrecorded permutation sign.  The normalized
transpose graph has codimension $2\dim A_U=\dim A_U^2$, so postcomposition
preserves codimension two.  Precompose with the algebraic source projector
and postcompose with algebraic degree-one projectors on the abelian factors.
The full degree-two action is then precisely the displayed tensor.
The same argument with the prescribed rational Hodge identification
transports it to any stipulated isogenous model.
\end{proof}

\subsection{Transfer to the rational quadratic locus}

\begin{lemma}\label{transfer:period}
Let $S$ be a projective K3 surface with a rational isometric embedding
$\iota:T_S\hookrightarrow V_P$.  There is an ample $P$-polarized K3
surface $R$ and a full rational Hodge isometry
\[
 \Phi:H^2(S,\Q)\xrightarrow{\sim}H^2(R,\Q),\qquad
 \Phi(T_S)=T_R.
\]
Both directions of this isometry are induced by algebraic correspondences.
\end{lemma}
\begin{proof}
Use the primitive embedding in Lemma~\ref{main:lattice} and an exact-$P$
surface supplied there.  Choose an integral ample $h\in P$ and a small
marked $P$-polarized period patch $\Omega$ about that surface, shrinking
so that $h$ stays ample.  This is an actual open image in the period
domain of $V_P$, and includes its Picard-jump fibers.

The line $z=\iota(H^{2,0}(S))$ is a point of this period domain.
Lemma~\ref{specialization:rational-density} gives $g\in O(V_P)(\Q)$
with $gz\in\Omega$.  Choose the fiber $R$ with period exactly $gz$.
Then $g\iota:T_S\to H^2(R,\Q)$ is an isometry taking the holomorphic
line to that of $R$.  The two full second-cohomology spaces have the
rational K3 form.  Witt's extension theorem extends this map to a
full rational isometry $\Phi$.  It preserves the holomorphic line,
its conjugate by rationality, and their orthogonal complement, which
is the $(1,1)$ space.  Hence $\Phi$ is Hodge.  Lefschetz $(1,1)$ gives
$\Phi(\NS(S)_{\Q})=\NS(R)_{\Q}$ and therefore $\Phi(T_S)=T_R$.
Buskin's theorem \cite[Theorem~1.1]{Buskin} algebraizes $\Phi$ and its
inverse.  It requires neither integrality nor preservation of an ample
class.  Density was used only to enter the patch; the comparison uses
the exact realized period.
\end{proof}

\begin{proposition}\label{transfer:tensor}
Under the isometric-embedding hypothesis of Theorem~\ref{main:theorem},
the exact tensor $\iota_{w,E}$ on $T_S$ is algebraic for every stipulated
$w$, $E$, and transported isogenous even-Clifford model.
\end{proposition}
\begin{proof}
Choose $R,\Phi$ by Lemma~\ref{transfer:period} and set
$\phi=\Phi|_{T_S}:T_S\simeq T_R$.  The isometry induces a unital
Clifford algebra isomorphism
\[
 c_\phi:C^+(T_S,q)\xrightarrow{\sim}C^+(T_R,q),
 \qquad c_\phi(vaw)=\phi(v)c_\phi(a)\phi(w).
\]
It is a weight-one Hodge isomorphism because it takes the oriented
period bivector defining left multiplication to the corresponding
bivector for $R$.  For the broader functoriality under Hodge similarities,
see \cite[Proposition~3.1 and Lemma~3.5]{VarescoSimilarities}; the
multiplier here is one, with the prescribed right factor transported
as in the displayed identity.
Proposition~\ref{transfer:ample} gives the standard
tensor on $T_R$ for $\phi(w)$ and the polarization transported from $E$.
Precompose its cycle with the algebraic correspondence for $\Phi$ and
transport both abelian slots by $c_\phi^{-1}$.  These rational Hodge
maps on $H^1$ are algebraic by Lemma~\ref{transfer:compression}.
The displayed Clifford identity proves equality with the specified
tensor, not merely with an abstract Hodge embedding.  Source and
degree-one target projectors isolate its full action.  All operations
are degree-preserving on their intermediate factors, so the resulting
cycle has codimension two.  The same algebraic transport treats the
stipulated isogenous model.
\end{proof}

\begin{proof}[Proof of Theorem~\ref{main:theorem}]
Proposition~\ref{transfer:tensor} proves the exact standard-tensor
assertion under the stated rational isometric-embedding hypothesis,
with every prescribed right vector, polarization, and transported model.
Choose any one such tensor and apply Theorem~\ref{ext:criterion}.
This gives rational HC and GHC in every degree of every self-power.
For an ample $P$-polarized surface, the actual transcendental space
is a nondegenerate subspace of its fixed perpendicular $U_P\simeq V_P$
by Lemma~\ref{main:lattice}, so this case is included as asserted.
\end{proof}

\subsection{Hilbert schemes and moduli spaces}

\begin{corollary}\label{transfer:moduli}
Let $S$ be a smooth projective complex K3 surface admitting a rational
isometric embedding $(T_S,q)\hookrightarrow V_P$.
For every $n,r\ge1$, rational HC and rational GHC hold in every degree
of $(S^{[n]})^r$, where $S^{[n]}$ is the Hilbert scheme of length-$n$
subschemes of $S$.

More generally, let $\alpha\in\operatorname{Br}(S)$ and let $M$ be either
a smooth projective moduli space of Gieseker-stable $\alpha$-twisted
sheaves on $S$, or a smooth projective moduli space of $\sigma$-stable
objects in $D^b(S,\alpha)$ for a generic Bridgeland stability condition
$\sigma$.
Then rational HC and rational GHC hold in every degree of $M^r$ for
every $r\ge1$.  If $M$ is hyperk\"ahler, the same conclusions hold for
every smooth projective hyperk\"ahler variety birational to $M$ and all
its self-powers.
\end{corollary}
\begin{proof}
The assertion for zero-dimensional $M$ is immediate.
Write $h(-)$ for the Chow motive with rational coefficients.
The decomposition of de Cataldo--Migliorini
\cite[Theorem~6.2.1]{deCataldoMigliorini} expresses $h(S^{[n]})$ as
a sum of Tate twists of summands of motives of powers of $S$; the
summands corresponding to symmetric products are cut out by rational
averaging projectors.  For either stated moduli space, B\"ulles
\cite[Theorem~0.1]{BullesModuli} likewise makes $h(M)$ a direct summand
of a finite sum of Tate twists of motives $h(S^k)$.
Tensoring these split correspondences gives analogous splittings for
every self-power, again through Tate twists of powers of $S$.
Thus these varieties are motivated by $S$ in the sense of Arapura.
Theorem~\ref{main:theorem} supplies HC and GHC for every power of $S$,
so Arapura's transfer lemma \cite[Lemma~4.2]{ArapuraMotivation},
which retains geometric coniveau with its Tate shifts, gives both
conjectures for every power of each indicated variety.
Finally, Rie\ss's specialization
correspondence for birational smooth projective hyperk\"ahler varieties
\cite[Theorem~3.2 and Section~3.3]{RiessBirational} gives an isomorphism
of Chow motives, as formulated in
\cite[Proposition~1.4]{BullesModuli}.  This proves the last assertion.
\end{proof}


\begin{thebibliography}{99}
\raggedright

\bibitem{Andre}
Andr{\'e}, Yves, \emph{Pour une th{\'e}orie inconditionnelle des motifs}, Publications Math{\'e}matiques de l'IH{\'E}S \textbf{83} (1996), 5--49. \href{https://doi.org/10.1007/BF02698643}{doi:10.1007/BF02698643}.

\bibitem{ArapuraMotivation}
Arapura, Donu, \emph{Motivation for {Hodge} cycles}, Advances in Mathematics \textbf{207} (2006), no.~2, 762--781. Cited locators follow arXiv:math/0501348v3 (16 December 2005). \href{https://doi.org/10.1016/j.aim.2006.01.005}{doi:10.1016/j.aim.2006.01.005}; \url{https://arxiv.org/abs/math/0501348v3}.

\bibitem{ArtinResolution}
Artin, Michael, \emph{Algebraic construction of {Brieskorn}'s resolutions}, Journal of Algebra \textbf{29} (1974), no.~2, 330--348. \href{https://doi.org/10.1016/0021-8693(74)90102-1}{doi:10.1016/0021-8693(74)90102-1}.

\bibitem{BirkenhakeLange}
Birkenhake, Christina; Lange, Herbert, \emph{Complex Abelian Varieties}, second edition, Grundlehren der mathematischen Wissenschaften \textbf{302}, Springer-Verlag, 2004. \href{https://doi.org/10.1007/978-3-662-06307-1}{doi:10.1007/978-3-662-06307-1}.

\bibitem{BolognesiLaterveer}
Bolognesi, Michele; Laterveer, Robert, \emph{A 9-dimensional family of {K3} surfaces with finite-dimensional motive}, Rendiconti del Circolo Matematico di Palermo, Series 2 \textbf{73} (2024), 2313--2331. Cited locators follow arXiv:2310.11981v2 (16 April 2024). \href{https://doi.org/10.1007/s12215-024-01036-0}{doi:10.1007/s12215-024-01036-0}; \url{https://arxiv.org/abs/2310.11981v2}.

\bibitem{BDPP}
Boucksom, S{\'e}bastien; Demailly, Jean-Pierre; P{\u a}un, Mihai; Peternell, Thomas, \emph{The pseudo-effective cone of a compact {K{\"a}hler} manifold and varieties of negative {Kodaira} dimension}, Journal of Algebraic Geometry \textbf{22} (2013), no.~2, 201--248. Cited locators follow arXiv:math/0405285v1. \href{https://doi.org/10.1090/S1056-3911-2012-00574-8}{doi:10.1090/S1056-3911-2012-00574-8}; \url{https://arxiv.org/abs/math/0405285v1}.

\bibitem{Brieskorn}
Brieskorn, Egbert, \emph{Die Aufl{\"o}sung der rationalen Singularit{\"a}ten holomorpher Abbildungen}, Mathematische Annalen \textbf{178} (1968), 255--270. \href{https://doi.org/10.1007/BF01352140}{doi:10.1007/BF01352140}.

\bibitem{BullesModuli}
B{\"u}lles, Tim-Henrik, \emph{Motives of moduli spaces on {K3} surfaces and of special cubic fourfolds}, Manuscripta Mathematica \textbf{161} (2020), no.~1--2, 109--124. Cited locators follow arXiv:1806.08284v1 (21 June 2018). \href{https://doi.org/10.1007/s00229-018-1086-0}{doi:10.1007/s00229-018-1086-0}; \url{https://arxiv.org/abs/1806.08284v1}.

\bibitem{Buskin}
Buskin, Nikolay, \emph{Every rational {Hodge} isometry between two {K3} surfaces is algebraic}, Journal f{\"u}r die reine und angewandte Mathematik \textbf{755} (2019), 127--150. Cited locators follow arXiv:1510.02852v3. \href{https://doi.org/10.1515/crelle-2017-0027}{doi:10.1515/crelle-2017-0027}; \url{https://arxiv.org/abs/1510.02852v3}.

\bibitem{deCataldoMigliorini}
de Cataldo, Mark Andrea A.; Migliorini, Luca, \emph{The {Chow} groups and the motive of the {Hilbert} scheme of points on a surface}, Journal of Algebra \textbf{251} (2002), no.~2, 824--848. \href{https://doi.org/10.1006/jabr.2001.9105}{doi:10.1006/jabr.2001.9105}.

\bibitem{DeligneHodgeII}
Deligne, Pierre, \emph{Th{\'e}orie de {Hodge}, {II}}, Publications Math{\'e}matiques de l'IH{\'E}S \textbf{40} (1971), 5--57. \href{https://doi.org/10.1007/BF02684692}{doi:10.1007/BF02684692}.

\bibitem{Deligne}
Deligne, Pierre, \emph{La conjecture de {Weil} pour les surfaces {K3}}, Inventiones Mathematicae \textbf{15} (1972), 206--226. \href{https://doi.org/10.1007/BF01404126}{doi:10.1007/BF01404126}.

\bibitem{DeligneHodgeIII}
Deligne, Pierre, \emph{Th{\'e}orie de {Hodge}, {III}}, Publications Math{\'e}matiques de l'IH{\'E}S \textbf{44} (1974), 5--77. \href{https://doi.org/10.1007/BF02685881}{doi:10.1007/BF02685881}.

\bibitem{DeligneMilne}
Deligne, Pierre; Milne, James S., \emph{Tannakian Categories}, in \emph{Hodge Cycles, Motives, and Shimura Varieties}, Lecture Notes in Mathematics \textbf{900}, Springer, 1982, 101--228. Cited locators follow the corrected author version of 15 August 2012, which retains the original numbering. \href{https://doi.org/10.1007/978-3-540-38955-2_4}{doi:10.1007/978-3-540-38955-2\_4}; \url{https://www.jmilne.org/math/xnotes/tc.pdf}.

\bibitem{Demailly2012}
Demailly, Jean-Pierre, \emph{Complex Analytic and Differential Geometry}, 2012 version. Chapter~II, Section~8.2, Theorem~8.8. \url{https://people.math.harvard.edu/~demarco/Math274/Demailly_ComplexAnalyticDiffGeom.pdf}.

\bibitem{DiLorenzoFringuelliVistoli}
Di Lorenzo, Andrea; Fringuelli, Roberto; Vistoli, Angelo, \emph{Integral {Picard} group of moduli of polarized {K3} surfaces}, Bulletin of the London Mathematical Society \textbf{56} (2024), no.~2, 833--846. Cited locators follow arXiv:2305.07574v2 (3 November 2023). \href{https://doi.org/10.1112/blms.12968}{doi:10.1112/blms.12968}; \url{https://arxiv.org/abs/2305.07574v2}.

\bibitem{DonagiPrym}
Donagi, Ron, \emph{The fibers of the Prym map}, arXiv:alg-geom/9206008v2, 23 June 1992. \url{https://arxiv.org/abs/alg-geom/9206008v2}.

\bibitem{FloccariFuOG6}
Floccari, Salvatore; Fu, Lie, \emph{The Hodge conjecture for Weil fourfolds with discriminant 1 via singular OG6-varieties}, Journal de Math\'ematiques Pures et Appliqu\'ees \textbf{210} (2026), Article~103876. Cited locators follow the author-hosted version. \href{https://doi.org/10.1016/j.matpur.2026.103876}{doi:10.1016/j.matpur.2026.103876}; \url{https://irma.math.unistra.fr/~lfu/articles/HCWeilViaOG6.pdf}.

\bibitem{Fulton}
Fulton, William, \emph{Intersection Theory}, second edition, Ergebnisse der Mathematik und ihrer Grenzgebiete, 3.~Folge \textbf{2}, Springer-Verlag, 1998. \href{https://doi.org/10.1007/978-1-4612-1700-8}{doi:10.1007/978-1-4612-1700-8}.

\bibitem{GrauertProperImage}
Grauert, Hans, \emph{Ein {Theorem} der analytischen {Garbentheorie} und die {Modulr{\"a}ume} komplexer {Strukturen}}, Publications Math{\'e}matiques de l'IH{\'E}S \textbf{5} (1960), 5--64. \href{https://doi.org/10.1007/BF02684746}{doi:10.1007/BF02684746}.

\bibitem{GrothendieckHilbert}
Grothendieck, Alexander, \emph{Techniques de construction et th{\'e}or{\`e}mes d'existence en g{\'e}om{\'e}trie alg{\'e}brique {IV}: les sch{\'e}mas de {Hilbert}}, S{\'e}minaire Bourbaki, 1960/61, Expos{\'e} 221, 249--276. \url{https://numdam.org/item/SB_1960-1961__6__249_0/}. See also the 1962 erratum, which does not change the statement of Theorem~3.2: \url{https://www.numdam.org/item/SB_1961-1962__7__302_1.pdf}.

\bibitem{Huybrechts}
Huybrechts, Daniel, \emph{Lectures on {K3} Surfaces}, Cambridge Studies in Advanced Mathematics \textbf{158}, Cambridge University Press, 2016. Cited locators follow the 2015 prepublication copy; the published book incorporates further corrections. \href{https://doi.org/10.1017/CBO9781316594193}{doi:10.1017/CBO9781316594193}; \url{https://ncatlab.org/nlab/files/HuybrechtsLecturesOnK3.pdf}.

\bibitem{KraftProcesi}
Kraft, Hanspeter; Procesi, Claudio, \emph{Classical Invariant Theory: A Primer}, preliminary version, July 1996. \url{https://dmi.unibas.ch/fileadmin/user_upload/dmi/Personen/Kraft_Hanspeter/Classical_Invariant_Theory.pdf}.

\bibitem{KS}
Kuga, Michio; Satake, Ichiro, \emph{Abelian varieties attached to polarized {K3}-surfaces}, Mathematische Annalen \textbf{169} (1967), 239--242. \href{https://doi.org/10.1007/BF01399540}{doi:10.1007/BF01399540}.

\bibitem{LehrerZhangSO}
Lehrer, Gustav I.; Zhang, Ruibin, \emph{Invariants of the special orthogonal group and an enhanced {Brauer} category}, L'Enseignement Math{\'e}matique \textbf{63} (2017), no.~1--2, 181--200. Cited locators follow arXiv:1612.03998v1 (13 December 2016). \href{https://doi.org/10.4171/LEM/63-1/2-6}{doi:10.4171/LEM/63-1/2-6}; \url{https://arxiv.org/abs/1612.03998v1}.

\bibitem{LiBlowup}
Li, Li, \emph{Chow motive of Fulton--MacPherson configuration spaces and wonderful compactifications}, Michigan Mathematical Journal \textbf{58} (2009), no.~2, 565--598. Cited locators follow arXiv:math/0611459v2 (1 January 2009). \url{https://arxiv.org/abs/math/0611459v2}.

\bibitem{MilnorStasheff}
Milnor, John W.; Stasheff, James D., \emph{Characteristic Classes}, Annals of Mathematics Studies \textbf{76}, Princeton University Press, 1974.

\bibitem{MumfordTheta}
Mumford, David, \emph{Tata Lectures on Theta {I}}, Progress in Mathematics \textbf{28}, Birkh{\"a}user Boston, 1983. \href{https://doi.org/10.1007/978-1-4899-2843-6}{doi:10.1007/978-1-4899-2843-6}.

\bibitem{Paranjape}
Paranjape, Kapil, \emph{Abelian varieties associated to certain {K3} surfaces}, Compositio Mathematica \textbf{68} (1988), no.~1, 11--22. \url{https://numdam.org/item/CM_1988__68_1_11_0/}.

\bibitem{PetersSterk}
Peters, Chris; Sterk, Hans, \emph{Symmetric and Quadratic Forms, with Applications to Coding Theory, Algebraic Geometry and Topology}, author-hosted draft, version label ``JUNE 2024,'' dated 16 September 2024. \url{https://www-fourier.univ-grenoble-alpes.fr/~peters/Books/QuadraticForms/QuadForms.pdf}.

\bibitem{Poon}
Poon, Flora, \emph{{Kuga--Satake} construction on families of {K3} surfaces of {Picard} rank 14}, Mathematische Nachrichten \textbf{299} (2026), no.~8, 1894--1916. \href{https://doi.org/10.1002/mana.70173}{doi:10.1002/mana.70173}.

\bibitem{RiessBirational}
Rie\ss, Ulrike, \emph{On the Chow ring of birational irreducible symplectic varieties}, Manuscripta Mathematica \textbf{145} (2014), no.~3--4, 473--501. Cited locators follow arXiv:1304.4404v2 (11 September 2014). \href{https://doi.org/10.1007/s00229-014-0698-2}{doi:10.1007/s00229-014-0698-2}; \url{https://arxiv.org/abs/1304.4404v2}.

\bibitem{SaitoMHMIntro}
Saito, Morihiko, \emph{Introduction to mixed Hodge modules}, Ast{\'e}risque \textbf{179--180} (1989), 145--162. \url{https://www.numdam.org/article/AST_1989__179-180__145_0.pdf}.

\bibitem{SaitoMHC}
Saito, Morihiko, \emph{Mixed Hodge complexes on algebraic varieties}, Mathematische Annalen \textbf{316} (2000), no.~2, 283--331. Cited locators follow arXiv:math/9906088v1 (13 June 1999). \href{https://doi.org/10.1007/s002080050014}{doi:10.1007/s002080050014}; \url{https://arxiv.org/abs/math/9906088v1}.

\bibitem{SchnuererSoergel}
Schn{\"u}rer, Olaf M.; Soergel, Wolfgang, \emph{Proper base change for separated locally proper maps}, Rendiconti del Seminario Matematico della Universit{\`a} di Padova \textbf{135} (2016), 223--250. Cited locators follow arXiv:1404.7630v2 (5 November 2014). \href{https://doi.org/10.4171/RSMUP/135-13}{doi:10.4171/RSMUP/135-13}; \url{https://arxiv.org/abs/1404.7630v2}.

\bibitem{SerreGAGA}
Serre, Jean-Pierre, \emph{G{\'e}om{\'e}trie alg{\'e}brique et g{\'e}om{\'e}trie analytique}, Annales de l'Institut Fourier \textbf{6} (1956), 1--42. \href{https://doi.org/10.5802/aif.59}{doi:10.5802/aif.59}.

\bibitem{VanGeemenKS}
van Geemen, Bert, \emph{Kuga--Satake varieties and the Hodge conjecture}, in \emph{The Arithmetic and Geometry of Algebraic Cycles}, B. B. Gordon et al. (eds.), Kluwer Academic Publishers, 2000, 51--82. Cited locators follow arXiv:math/9903146v1 (24 March 1999). \url{https://arxiv.org/abs/math/9903146v1}.

\bibitem{VanGeemenReal}
van Geemen, Bert, \emph{Real multiplication on K3 surfaces and Kuga--Satake varieties}, Michigan Mathematical Journal \textbf{56} (2008), no.~2, 375--399. Cited locators follow arXiv:math/0609839v1. \href{https://doi.org/10.1307/mmj/1224783519}{doi:10.1307/mmj/1224783519}; \url{https://arxiv.org/abs/math/0609839v1}.

\bibitem{VanGeemenVerra}
van Geemen, Bert; Verra, Alessandro, \emph{Quaternionic {Pryms} and {Hodge} classes}, Topology \textbf{42} (2003), no.~1, 35--53. Cited locators follow the author-hosted arXiv:math/0103111v1 (17 March 2001). \href{https://doi.org/10.1016/S0040-9383(02)00004-6}{doi:10.1016/S0040-9383(02)00004-6}; \url{https://sites.unimi.it/vangeemen/0103111.pdf}.

\bibitem{VarescoSimilarities}
Varesco, Mauro, \emph{Hodge similarities, algebraic classes, and Kuga--Satake varieties}, Mathematische Zeitschrift \textbf{305} (2023), Article~69. Cited locators follow arXiv:2304.02519v3 (2 November 2023). \href{https://doi.org/10.1007/s00209-023-03390-8}{doi:10.1007/s00209-023-03390-8}; \url{https://arxiv.org/abs/2304.02519v3}.

\bibitem{VarescoPowers}
Varesco, Mauro, \emph{The {Hodge} conjecture for powers of {K3} surfaces of {Picard} number 16}, Michigan Mathematical Journal \textbf{75} (2025), no.~5, 1085--1116. Cited locators follow arXiv:2203.09778v3 (13 May 2024). \href{https://doi.org/10.1307/mmj/20236349}{doi:10.1307/mmj/20236349}; \url{https://arxiv.org/abs/2203.09778v3}.

\bibitem{Wlodarczyk}
W{\l}odarczyk, Jaros{\l}aw, \emph{Simple {Hironaka} resolution in characteristic zero}, Journal of the American Mathematical Society \textbf{18} (2005), no.~4, 779--822. \href{https://doi.org/10.1090/S0894-0347-05-00493-5}{doi:10.1090/S0894-0347-05-00493-5}.

\bibitem{Zarhin1983}
Zarhin, Yuri G., \emph{{Hodge} groups of {K3} surfaces}, Journal f{\"u}r die reine und angewandte Mathematik \textbf{341} (1983), 193--220. \href{https://doi.org/10.1515/crll.1983.341.193}{doi:10.1515/crll.1983.341.193}.

\end{thebibliography}
\end{document}
